% ============================================================================
% Fichier     : FAV_main.tex
% Titre       : Forensics et analyse de vulnérabilités
% Auteur      : MINKA MI NGUIDJOÏ Thierry Emmanuel
% Institution : Chaire UNESCO d'Afrique Centrale Accès TIC, ENSP / LIMSI
% Version     : 1.0 -- novembre 2024 (traité complet : P00 à C16, annexes A à F)
% ============================================================================

\documentclass[12pt, twoside, openright, a4paper]{book}

\usepackage{Style_FAV}
\raggedbottom

\addbibresource{bibliographie/FAV_references.bib}

\hypersetup{
    pdftitle={Forensics et analyse de vulnérabilités},
    pdfauthor={MINKA MI NGUIDJOÏ Thierry Emmanuel},
    pdfsubject={MSN-2S3-SEC07 · Master Professionnel en Sécurité des Systèmes Numériques · Chaire UNESCO d'Afrique Centrale Accès TIC},
    pdfkeywords={investigation numérique, forensique, ISO/IEC 27037, analyse mémoire, analyse disque, analyse réseau, CVSS, EPSS, KEV, gestion des vulnérabilités, NEK}
}

\begin{document}

% --------------------------------------------------------------------------
% ÉLÉMENTS PRÉLIMINAIRES (pages romaines)
% --------------------------------------------------------------------------
\frontmatter

    % ==============================================================================
% Fichier : front/00_page_titre.tex
% Titre   : Page de titre - Forensics et Analyse de Vulnérabilités
% Auteur  : MINKA MI NGUIDJOÏ Thierry Emmanuel
% Version : 1.0 -- novembre 2024
% ==============================================================================

\begin{titlepage}
\thispagestyle{empty}

% --- Fond décoratif : bandeaux bleu marine et filets or ---
\begin{tikzpicture}[remember picture, overlay]
    \fill[CouleurPrimaire]
        (current page.north west) rectangle
        ([yshift=-3.6cm] current page.north east);
    \fill[CouleurAccent]
        ([yshift=-3.6cm] current page.north west) rectangle
        ([yshift=-4.0cm] current page.north east);
    \fill[CouleurPrimaire]
        (current page.south west) rectangle
        ([yshift=2.5cm] current page.south east);
    \fill[CouleurAccent]
        ([yshift=2.5cm] current page.south west) rectangle
        ([yshift=2.9cm] current page.south east);
    \node[align=center, text=white] at ([yshift=-1.8cm] current page.north) {%
        {\fontsize{11}{13}\selectfont
        \textsc{République du Cameroun - Paix, Travail, Patrie}}\\[0.3ex]
        {\color{CouleurAccent!60!white}\rule{0.4\paperwidth}{0.4pt}}\\[0.6ex]
        {\fontsize{10}{12}\selectfont\textbf{\institutionprimaire}}\\[0.2ex]
        {\fontsize{8}{10}\selectfont \tutelle}\\[0.2ex]
        {\fontsize{9}{11}\selectfont \unitepedago}};
    \node[fill=white, rounded corners=4pt, inner sep=2pt]
        at ([xshift=1.75cm, yshift=-1.8cm] current page.north west)
        {\includegraphics[height=2.6cm]{figures/logos/enspy_logo.jpg}};
    \node[fill=white, rounded corners=4pt, inner sep=2pt]
        at ([xshift=-1.75cm, yshift=-1.8cm] current page.north east)
        {\includegraphics[height=2.6cm]{figures/logos/unesco_logo.jpeg}};
    \node[align=center] at ([yshift=1.25cm] current page.south) {%
        {\color{white}\bfseries\large Année académique \anneescolaire}\\[0.3ex]
        {\color{CouleurAccent!70!white}\small
        Version 1.0 | Document pédagogique de l'UE \codeUE}};
\end{tikzpicture}

\vspace*{0.6cm}

\begin{center}
    {\color{CouleurPrimaire}\rule{0.85\textwidth}{1.2pt}}\\[1.2ex]
    {\fontsize{26}{32}\selectfont\bfseries\color{CouleurPrimaire}
        Forensics et Analyse\\[0.4ex]
        de Vulnérabilités}\\[0.8ex]
    {\color{CouleurAccent}\rule{0.85\textwidth}{2pt}}\\[0.5ex]
    {\color{CouleurPrimaire}\rule{0.85\textwidth}{0.5pt}}
\end{center}

\vspace{0.8cm}

\begin{center}
    {\fontsize{13}{16}\selectfont\itshape\color{CouleurSecondaire}
        De la trace à la faille~:\\
        enquêter, prouver, corriger}
\end{center}

\vspace{0.5cm}

\begin{center}
    \begin{tikzpicture}
        \node[draw=CouleurPrimaire, fill=CouleurPrimaire, rounded corners=6pt,
              inner xsep=14pt, inner ysep=6pt]{%
            \color{white}\bfseries\large \sigleCours{} - UE~\codeUE
        };
    \end{tikzpicture}
\end{center}

\vspace{1.0cm}

% --- Schéma de la discipline : de la préservation au rapport, en passant par
%     la vulnérabilité et sa correction ---
\begin{center}
\begin{tikzpicture}[
    etape/.style = {draw=CouleurBordDef, fill=CouleurDef, rounded corners=3pt,
        minimum width=2.5cm, minimum height=0.8cm, align=center,
        font=\small\bfseries\color{CouleurPrimaire}, inner sep=2pt},
    arrow/.style = {-{Stealth[length=4pt]}, thick, color=CouleurSecondaire}
]
\node[etape] (n1) at (0, 0)      {Préserver};
\node[etape] (n2) at (3.2, 0)    {Acquérir};
\node[etape] (n3) at (6.4, 0)    {Analyser};
\node[etape] (n4) at (6.4, -1.4) {Vulnérabilité};
\node[etape] (n5) at (3.2, -1.4) {Correction};
\node[etape] (n6) at (0, -1.4)   {Rapport};
\draw[arrow] (n1) -- (n2); \draw[arrow] (n2) -- (n3); \draw[arrow] (n3) -- (n4);
\draw[arrow] (n4) -- (n5); \draw[arrow] (n5) -- (n6); \draw[arrow] (n6) -- (n1);
\node[draw=CouleurPrimaire, fill=CouleurPrimaire, rounded corners=4pt,
      align=center, font=\footnotesize\bfseries\color{white}, inner xsep=8pt]
      at (3.2, -2.8)
      {Chaque conclusion~: \OUI{} / \NON{} / \NEK{}, la constatation séparée de l'analyse};
\end{tikzpicture}
\end{center}

\vspace{0.6cm}

\begin{center}
    {\large\bfseries\color{CouleurPrimaire}
        \auteurcours}\\[0.5ex]
    {\normalsize\color{CouleurSecondaire}
       Enseignant Chercheur | \laboratoire{} | ENSPY}\\[0.2ex]
    {\small\color{CouleurBordInfo}
        \institutionprimaire}\\[0.2ex]
    {\small\color{CouleurBordInfo}
        \href{mailto:maletyon@proton.me}{maletyon@proton.me}}
\end{center}

\end{titlepage}

    \cleardoublepage

    % ============================================================================
% Fichier  : front/01_avant_propos.tex (Avant-propos)
% Version  : 1.0 -- novembre 2024 (proposition à relire par l'auteur)
% ============================================================================

\chapter*{Avant-propos}
\addcontentsline{toc}{chapter}{Avant-propos}
\markboth{Avant-propos}{Avant-propos}

Une microfinance de Bafoussam perd quarante-sept millions de francs CFA en
une nuit. Trois mois plus tôt, l'éditeur de sa passerelle d'accès à distance
avait publié un correctif pour la faille que les attaquants ont empruntée.
L'enquête qui suit doit dire ce qui s'est passé, qui l'a fait si on peut le
savoir, et par où ils sont entrés. La dernière question renvoie à une autre
discipline, l'analyse de vulnérabilités, qui aurait pu éviter la première.

Ce cours réunit ces deux disciplines parce que, sur le terrain, elles se
répondent. L'investigation numérique part d'une trace et remonte vers une
cause~; la gestion des vulnérabilités part d'une faiblesse connue et cherche
à la fermer avant qu'elle ne laisse une trace. L'une sans l'autre laisse un
travail à moitié fait~: une enquête qui ne débouche sur aucune correction
prépare l'incident suivant, et un programme de correctifs qui ignore ce que
les enquêtes révèlent corrige au hasard.

Il s'adresse aux étudiants du Master Professionnel en Sécurité des Systèmes
Numériques de la Chaire UNESCO d'Afrique centrale, et aux praticiens qui
exercent ou exerceront comme analystes forensiques, analystes en réponse à
incident, analystes en gestion des vulnérabilités ou RSSI. Il suppose les
bases de la sécurité des systèmes. Il laisse au cours d'Ethical Hacking et
Pentesting (UE MSN-2S3-SEC03) l'exploitation des failles, et au cours de
droit du numérique (UE MSN-2S3-JUR01) la protection des données
personnelles.

Trois partis pris le traversent. Le premier est de tenir la preuve pour
fragile~: une copie mal faite, un horodatage mal lu ou un scellé mal tenu
suffisent à faire écarter une constatation exacte. Le deuxième est de
séparer, à chaque page, ce que l'on a constaté de ce que l'on en déduit, et
de dire avec quelle confiance on le déduit. Le troisième est de coter ce que
l'on ne sait pas~: l'instrument \OUI{}~/ \NON{}~/ \NEK{} du programme de
recherche de l'auditologie sert ici à dire, pour chaque conclusion
d'enquête, si elle est établie, réfutée, ou hors de portée des preuves
disponibles.

Les deux organisations qui accompagnent le lecteur, NYANGA CRÉDIT et
SOLITRANS LOGISTIQUE, sont fictives, comme leurs dirigeants, leurs
prestataires, les dates, les adresses et les montants. La faille de la
passerelle porte un identifiant fictif. Les mécanismes, les outils, les
normes et les textes cités sont exacts. Les articles de loi sont cités dans
leur texte officiel~; les normes sous licence sont paraphrasées.

Le livre est accompagné d'un livret des corrigés, qui explique et annote
chaque exercice, et de NDAP SUKULU, où les fiches de questions et les dix
laboratoires du cours permettent de pratiquer chaque chapitre.

\bigskip
\hfill\begin{minipage}{0.5\textwidth}
\raggedleft
\auteurcours\\
\laboratoire, Yaoundé, novembre 2024
\end{minipage}

    \cleardoublepage

    % ============================================================================
% Fichier  : front/02_mode_emploi.tex (Mode d'emploi)
% Version  : 1.0 -- novembre 2024
% ============================================================================
% 30 h de cours magistral = 15 séances de 2 h ; TD (fiches, exercices) et TP
% (dix labs) faits par l'étudiant chez lui, hors des 30 h.
% ============================================================================

\chapter*{Mode d'emploi du cours}
\addcontentsline{toc}{chapter}{Mode d'emploi du cours}
\markboth{Mode d'emploi}{Mode d'emploi}

\section*{Structure générale}

Le cours compte un prologue et six parties, soit seize chapitres, pour
30~heures de cours magistral réparties en quinze séances de deux heures.
Les travaux dirigés et les travaux pratiques se font à la maison, hors de
ces 30~heures~: pour les travaux dirigés, la fiche de questions de chaque
chapitre et ses exercices de fin de chapitre (environ dix heures au
total)~; pour les travaux pratiques, les dix laboratoires de NDAP SUKULU,
d'une heure chacun. Les corrigés des exercices sont réunis dans un livret
séparé.

\begin{tableaucours}{cL{6.4cm}cL{3.2cm}}{%
    \tableauHeader{Partie} &
    \tableauHeader{Thème} &
    \tableauHeader{Chap.} &
    \tableauHeader{Niveau}
}
0 & Prologue~: la nuit du 14 mars à Bafoussam & P & introductif \\
I & Fondements~: enquêter sur des traces numériques & 1--2 & fondamental \\
II & Méthode~: acquérir et analyser & 3--7 & fondamental à
intermédiaire \\
III & Cadre normatif et juridique & 8 & fondamental \\
IV & Vulnérabilités~: identifier, coter, corriger & 9--12 & intermédiaire \\
V & Défis avancés & 13--15 & intermédiaire à expert \\
VI & Recherche~: mesurer l'incertitude d'une enquête & 16 & recherche \\
\end{tableaucours}

\section*{Les quinze séances}

\begin{tableaucours}{cL{6.8cm}L{4.2cm}}{%
    \tableauHeader{Séance} &
    \tableauHeader{Contenu} &
    \tableauHeader{Travail à la maison}
}
1 & Prologue et Chap.~1~: l'incident, l'investigation numérique & Fiches P et 1 \\
2 & Chap.~2~: la preuve et la chaîne de possession & Fiche 2, lab 01 \\
3 & Chap.~3~: acquérir et préserver & Fiche 3, lab 02 \\
4 & Chap.~4~: l'analyse du disque & Fiche 4, lab 03 \\
5 & Chap.~5~: l'analyse de la mémoire & Fiche 5, lab 04 \\
6 & Chap.~6~: l'analyse du réseau & Fiche 6, lab 05 \\
7 & Chap.~7~: chronologie et corrélation & Fiche 7, lab 06 \\
8 & Chap.~8~: normes et droit de l'investigation & Fiche 8 \\
9 & Chap.~9~: qu'est-ce qu'une vulnérabilité & Fiche 9, lab 08 \\
10 & Chap.~10~: identifier sans exploiter & Fiche 10, lab 07 \\
11 & Chap.~11~: gérer les vulnérabilités & Fiche 11, lab 09 \\
12 & Chap.~12~: de l'artefact à la vulnérabilité & Fiche 12 \\
13 & Chap.~13~: anti-forensique, nuage et conteneurs & Fiche 13 \\
14 & Chap.~14~: l'enquête intégratrice & Fiche 14, lab 10 \\
15 & Chap.~15 et~16~: le rapport, l'incertitude d'une enquête & Fiches 15
et 16 \\
\end{tableaucours}

Cette répartition est indicative~; l'enseignant l'adapte au calendrier et au
rythme du groupe.

\section*{Le gabarit des chapitres}

Chaque chapitre suit le même plan~: une épigraphe, le niveau et les
prérequis, les objectifs, une situation d'ouverture tirée de l'un des deux
cas, des sections numérotées, ce qu'il faut retenir, des points de
contrôle, trois à quatre exercices gradués, des lectures, et une
transition.

\section*{Les encadrés}

\begin{tableaucours}{L{3.6cm}L{8.6cm}}{%
    \tableauHeader{Encadré} &
    \tableauHeader{Rôle}
}
Objectifs & Ce que l'on doit savoir faire à la fin du chapitre \\
Définition & Les notions, en gras dans le texte et indexées \\
Violation typique & Une erreur fréquente, et la règle qui l'évite \\
Information & Un complément, un renvoi, une nuance \\
Opération (NYANGA, SOLITRANS) & Le cas~: les constatations, puis
l'analyse, avec un niveau de confiance \\
Artefact & Une trace décrite~: source, emplacement, empreinte, ce qu'elle
prouve et ce qu'elle ne prouve pas \\
Chaîne de possession & Ce que le dossier doit contenir pour qu'un autre
puisse refaire l'examen \\
Lecture & Un décompte de cotes \OUI{}~/ \NON{}~/ \NEK{} et sa lecture \\
Métier & Ce que la notion change pour chacun des métiers du cours \\
Texte juridique & Un article cité dans son texte officiel, et sa portée \\
Ce qu'il faut retenir & Les faits et les règles du chapitre \\
Exercice & Trois niveaux de difficulté~: \niveaubase, \niveaumoyen,
\niveauexpert \\
\end{tableaucours}

\section*{Les conventions}

\begin{itemize}
    \item Les cotes~: \OUI{} (établi par une preuve), \NON{} (absent ou
          contredit par une preuve), \NEK{} (connaissance insuffisante pour
          trancher). Une conclusion d'enquête reçoit l'une de ces trois
          cotes.
    \item Constatation et analyse sont toujours séparées~; l'analyse porte
          un niveau de confiance (faible, moyenne, élevée).
    \item Les heures sont données en heure locale du Cameroun (UTC+1), sauf
          mention UTC explicite.
    \item Les textes~: seuls les articles lus dans leur version officielle
          sont cités entre guillemets~; les normes sous licence sont
          paraphrasées~; le degré de vérification figure à
          l'annexe~\ref{ann:textes}.
    \item Les adresses IP des cas appartiennent aux plages réservées à la
          documentation (RFC~5737), et l'identifiant de la faille de la
          passerelle, VULN-2024-4471, est fictif.
    \item Les cas sont fictifs~: noms en majuscules, montants en francs
          CFA, villes réelles.
\end{itemize}

    \cleardoublepage

    % ==============================================================================
% Fichier : front/03_ndap_sukulu.tex
% Titre   : NDAP SUKULU, l'environnement compagnon du cours
% Auteur  : MINKA MI NGUIDJOÏ Thierry Emmanuel
% Version : 1.0 -- novembre 2024
% ==============================================================================
% Texte commun aux cours de l'auteur, mis à jour en novembre 2024 (version en
% ligne, banques de 100 questions, compte de démonstration) ; le tableau des
% labs est propre au cours (front/03b_labs.tex).
% ==============================================================================

\chapter*{NDAP SUKULU, la salle de classe compagnon}
\addcontentsline{toc}{chapter}{NDAP SUKULU, la salle de classe compagnon}
\markboth{NDAP SUKULU}{NDAP SUKULU}

\begin{center}
    \includegraphics[width=0.62\textwidth]{figures/logos/logo_ndap_sukulu.png}
\end{center}

\section*{Un nom venu du pays bassa}

\begin{quote}\itshape
Le soir tombe sur Éséka. Un enfant revient des champs avec son grand-père et
lui demande pourquoi il doit marcher si loin, chaque matin, pour aller à
l'école. Le vieil homme s'arrête au bord du chemin et répond, en bassa~:
« \textup{Ndap sukulu}, la salle de classe, ce ne sont pas les murs. Ce sont
les murs qui viennent après. La salle de classe, c'est l'endroit où quelqu'un
t'attend pour t'apprendre ce qu'il sait. » L'enfant insiste~: « Et qui
m'attend~? » Le grand-père sourit~: « \textup{Malet yon}. Ton enseignant. »
\end{quote}

En bassa, la langue paternelle de l'auteur, \textit{ndap sukulu} désigne la
salle de classe et \textit{malet yon} signifie « ton enseignant ».
L'environnement compagnon de ce cours porte ce nom parce qu'il prolonge la
salle de classe au-delà des heures de cours~: il attend l'étudiant le soir,
le week-end, à distance, avec des questions et des pièces à examiner. Il
est édité par MaletYon et mis à la disposition des étudiants et des
praticiens qui souhaitent se perfectionner
(\href{mailto:maletyon@proton.me}{maletyon@proton.me}).

\section*{Ce que l'on y trouve}

\begin{itemize}
    \item Une fiche de questions par chapitre. Chaque tentative tire vingt
          questions d'une banque de cent (QCM, vrai ou faux, textes à
          trous) et dure trente minutes.
    \item Dix laboratoires. Chaque lab place l'étudiant devant les pièces
          d'une affaire du cours (images, captures, journaux, tableurs de
          balayage) et se termine par des points de contrôle.
    \item Un espace enseignant, qui corrige les remises, établit la
          synthèse de la classe et publie les résultats.
\end{itemize}

Les fiches et les labs sont des travaux à faire hors des heures de cours.
Ils ne remplacent ni le livre ni les exercices de fin de chapitre~: ils en
vérifient la maîtrise.

\section*{Comment y accéder}

\begin{enumerate}
    \item En ligne, à l'adresse
          \url{https://maletyon.github.io/cours-fav/ndap-sukulu/}, dans un
          navigateur récent~; sur téléphone, ajouter \texttt{telephone/} à
          la fin de l'adresse.
    \item Hors ligne, en téléchargeant l'archive de l'édition ordinateur
          depuis le dépôt \texttt{MaletYon/cours-fav}, dossier
          \texttt{ndap-sukulu}.
    \item Se connecter avec le matricule et le mot de passe remis par
          l'enseignant. Pour essayer avant d'être inscrit, un compte de
          démonstration est ouvert à tous~: matricule \texttt{ETU001}, mot
          de passe \texttt{eTU001}~; ses essais ne sont jamais transmis à
          l'enseignant.
    \item Pour un lab, vérifier d'abord l'empreinte SHA-256 des pièces
          (\texttt{sha256sum -c MANIFESTE\_SHA256.txt}), puis travailler sur
          une copie~: on n'examine jamais l'original.
\end{enumerate}

\begin{boxinfo}
Les fiches de questions sont aussi fournies au format texte dans le dossier
\texttt{fiches} du dépôt du cours, prêtes à être distribuées. Leurs
corrections figurent dans le livret de correction, à la suite des corrigés
des exercices de chaque chapitre.
\end{boxinfo}

\section*{Les laboratoires du cours}

% Tableau des labs du cours FAV dans NDAP SUKULU (propre au cours).
Le cours compte dix laboratoires, disponibles sous le code \codeUE. Les six
premiers reprennent les pièces de l'enquête NYANGA, les trois suivants le
programme de SOLITRANS, et le dernier conduit une enquête complète sur la
compromission du portail de SOLITRANS en février 2026, de la première trace
à la faille.

\begin{tableaucours}{lL{6.2cm}L{1.6cm}L{2.2cm}}{%
    \tableauHeader{Lab} &
    \tableauHeader{Sujet} &
    \tableauHeader{Chapitre} &
    \tableauHeader{Opération}
}
FAV-LAB-01 & Chaîne de possession et scellés & 2 & NYANGA \\
FAV-LAB-02 & Duplication et hachage & 3 & NYANGA \\
FAV-LAB-03 & Disque~: système de fichiers et artefacts & 4 & NYANGA \\
FAV-LAB-04 & Mémoire~: processus et connexions & 5 & NYANGA \\
FAV-LAB-05 & Réseau~: l'exfiltration & 6 & NYANGA \\
FAV-LAB-06 & Chronologie unifiée & 7 & NYANGA \\
FAV-LAB-07 & Lire un résultat de balayage & 10 & SOLITRANS \\
FAV-LAB-08 & Coter~: CVSS v4.0, EPSS, KEV & 9 & SOLITRANS \\
FAV-LAB-09 & Prioriser et suivre la remédiation & 11 & SOLITRANS \\
FAV-LAB-10 & Enquête intégratrice & 12, 14 & SOLITRANS \\
\end{tableaucours}
\vspace{-0.8em}
\begin{center}{\small\itshape Les laboratoires du cours dans NDAP SUKULU}\end{center}


    \cleardoublepage

    \tableofcontents
    \cleardoublepage

    \listoffigures
    \listoftables
    \cleardoublepage

% CORPS (pages arabes). Numérotation continue : le prologue est le chapitre 0,
% puis C01 à C16, un numéro par fichier.
% --------------------------------------------------------------------------
\mainmatter

    % PARTIE 0 · PROLOGUE OPÉRATIONNEL (niveau introductif)
    \part*{Prologue Opérationnel}
    \addcontentsline{toc}{part}{Prologue Opérationnel}

        \setcounter{chapter}{-1}
        % ============================================================================
% Fichier  : partie0/P00_nuit_bafoussam.tex
% Titre    : La Nuit du 14 Mars à Bafoussam
% Auteur   : MINKA MI NGUIDJOÏ Thierry Emmanuel
% Version  : 1.0 -- novembre 2024
% Corrigés : corriges/Corr_P00.tex
% ============================================================================
% Faits posés (bible des cas, Claude outputs/01_fiche_plan_bible.md) :
% NYANGA CRÉDIT, Bafoussam, 22 agences, 210 000 clients ; 23 virements
% frauduleux de 02 h 05 à 03 h 31 le 14/03/2025, 47 300 000 FCFA, 11
% portefeuilles ; exfiltration de 3,2 Go de SRV-FIC01 vers 203.0.113.47 de
% 01 h 12 à 03 h 47 ; faille VULN-2024-4471 (fictive) de la passerelle VPN,
% correctif du 15/12/2024, KEV le 19/12/2024, premier accès le 02/03/2025.
% ============================================================================

\chapter{La Nuit du 14 Mars à Bafoussam}
\label{chap:prologue}
\chaptermark{La nuit du 14 mars}

\epigraph{%
    « Ce que l'on fait dans la première heure\\
    décide de ce que l'on pourra prouver dans un an. »%
}{Principe du cours}

\niveauChapitre{Introductif}{%
    Aucun prérequis~: ce prologue ouvre le cours et en pose les questions.
}

\begin{boxobjectifs}
\begin{itemize}
    \item Reconnaître, dans un incident réel, les questions qui relèvent de
          l'investigation numérique et celles qui relèvent de l'analyse de
          vulnérabilités.
    \item Distinguer les gestes qui préservent la preuve de ceux qui la
          détruisent pendant la première heure.
    \item Séparer, dans un récit d'incident, la constatation de l'analyse.
    \item Coter une conclusion d'enquête \OUI{}, \NON{} ou \NEK{}.
    \item Situer les métiers qui interviennent, de l'alerte à la correction.
\end{itemize}
\end{boxobjectifs}

\bigskip

% ============================================================================
\section*{Avant de commencer : un vendredi matin}
% ============================================================================

Bafoussam, vendredi 14 mars 2025, 7~h~30. Au service des paiements de NYANGA
CRÉDIT\footnote{NYANGA CRÉDIT, SOLITRANS LOGISTIQUE, leurs dirigeants, leurs
prestataires, les adresses, les dates et les montants sont fictifs~;
l'identifiant VULN-2024-4471 est fictif. Les mécanismes techniques, les
outils, les normes et les textes cités sont exacts.}, une agente de
permanence ouvre le rapprochement de la nuit. Elle compte vingt-trois
virements sortants passés entre 2~h~05 et 3~h~31, tous par l'interface qui
relie l'établissement aux opérateurs de paiement mobile, tous vers onze
portefeuilles que personne ne connaît. Le total atteint 47\,300\,000 francs
CFA. Aucun agent n'était en poste à ces heures-là.

NYANGA CRÉDIT est un établissement de microfinance de deuxième catégorie,
installé à Bafoussam, présent dans vingt-deux agences de l'Ouest et du
Littoral, au service d'environ 210\,000 clients. Son informatique tient dans
une salle du siège~: un contrôleur de domaine, un serveur comptable, un
serveur de fichiers, l'interface de paiement mobile, et une passerelle
d'accès à distance qui permet aux agences et aux techniciens de se connecter
au siège.

À 8~h~15, le directeur des systèmes d'information, M.~KAMGA, arrive. Le
serveur de fichiers ralentit depuis la veille. Son premier réflexe est de le
redémarrer, puis de restaurer la sauvegarde de la veille. La responsable de la
sécurité, Mme~NDOUMBE, l'arrête~: elle fait débrancher le câble réseau du
serveur, sans l'éteindre, et demande que personne ne touche plus aux
machines du siège avant son retour. Elle revient vingt minutes plus tard avec
une clé USB qui contient un outil de capture de la mémoire vive.

\medskip

Ce prologue raconte les premières heures de l'incident et les questions
qu'elles posent. Les seize chapitres qui suivent y répondent une à une.

\begin{boxalert}
\textbf{Violation typique.} Redémarrer, nettoyer ou restaurer une machine
compromise avant d'avoir préservé ce qu'elle contient.

\smallskip
\textbf{Règle~:} on isole avant d'agir, on préserve avant d'analyser, et on
consigne chaque geste avec son heure et son auteur. Une machine redémarrée a
perdu sa mémoire vive, ses connexions ouvertes et une partie de ses
journaux~; une machine restaurée a perdu tout le reste.
\end{boxalert}

\separateur

% ============================================================================
\section{Ce que l'on sait à 9 heures}
\label{sec:p00-9h}
% ============================================================================

À 9~h~10, le directeur général, M.~TCHOUA, réunit une cellule de crise~:
M.~KAMGA, Mme~NDOUMBE et la juriste de l'établissement, Me~ETOUNDI. Mme~NDOUMBE
pose sur la table ce qui est établi, et seulement cela.

\begin{tableaucours}{L{5.2cm}L{4.4cm}c}{%
    \tableauHeader{Constatation} &
    \tableauHeader{Source} &
    \tableauHeader{Cote}
}
23 virements entre 2~h~05 et 3~h~31, 47,3~M~FCFA, 11 portefeuilles &
Journal de l'interface de paiement & \OUI{} \\
Les virements ont utilisé la clé d'accès de l'établissement à l'interface &
Journal de l'interface de paiement & \OUI{} \\
Aucun agent n'était connecté au siège entre 1~h et 5~h & Registre des
présences, journal de la passerelle & \OUI{} \\
Le serveur de fichiers a émis un volume sortant inhabituel dans la nuit &
Supervision du lien Internet & \OUI{} \\
Ce volume correspond à des fichiers de clients & Aucune pièce à ce stade &
\NEK{} \\
L'attaquant est entré par la passerelle d'accès à distance & Hypothèse du
DSI & \NEK{} \\
L'attaque vient d'un employé & Rumeur dans les couloirs & \NEK{} \\
\end{tableaucours}
\vspace{-0.8em}
\begin{center}{\small\itshape Tableau~0.1. Le point de 9~h~10}\end{center}

Trois lignes sont cotées \NEK{}~: aucune preuve ne permet encore de les
établir ou de les écarter. Une hypothèse, même plausible, ne devient pas une
constatation parce qu'un responsable la formule avec assurance. Mme~NDOUMBE
le dit à la cellule~: « Nous savons ce qui est sorti du compte. Nous ne
savons pas encore par où ils sont entrés. »

\begin{boxdef}
\textbf{\termecle{Constatation}}~: fait vérifiable, rattaché à une source
identifiée (un journal, une image, une capture), que toute personne
compétente retrouverait en refaisant l'examen.

\smallskip
\textbf{\termecle{Analyse}}~: interprétation d'une ou de plusieurs
constatations~; elle porte un niveau de confiance (faible, moyenne, élevée)
et peut être révisée par une constatation nouvelle.
\end{boxdef}

\begin{boxdef}
\textbf{\termecle{Cote OUI, NON, NEK}}~: réponse à une question d'enquête
fermée. \OUI{}~: la proposition est établie par une preuve. \NON{}~: elle est
contredite par une preuve. \NEK{} (\textit{not enough knowledge})~: les
preuves disponibles ne permettent pas de trancher. L'instrument vient du
programme de recherche de l'auditologie~\cite{minka2026nek}.
\end{boxdef}

% ============================================================================
\section{Les gestes de la première heure}
\label{sec:p00-gestes}
% ============================================================================

Ce que Mme~NDOUMBE fait entre 8~h~15 et 10~h~30 n'a rien d'improvisé~: elle
applique un ordre que l'on retrouve dans la RFC~3227~\cite{rfc3227} et dans
la norme ISO/IEC~27037~\cite{iso27037_2012}, celui de la volatilité. On
recueille d'abord ce qui disparaît le plus vite.

\begin{tableaucours}{cL{3.6cm}L{4.0cm}L{3.4cm}}{%
    \tableauHeader{Heure} &
    \tableauHeader{Geste} &
    \tableauHeader{Ce qu'il préserve} &
    \tableauHeader{Ce qu'il évite}
}
8~h~20 & Débrancher le câble réseau du serveur de fichiers, sans l'éteindre &
La mémoire vive, les processus, les connexions récentes & La poursuite de
l'exfiltration, l'effacement à distance \\
8~h~25 & Interdire toute connexion aux serveurs du siège & L'état des
disques et des journaux & Les modifications par l'équipe elle-même \\
8~h~40 & Désactiver la passerelle d'accès à distance & Les journaux de la
passerelle, figés & Un nouvel accès \\
10~h~02 & Capturer la mémoire du contrôleur de domaine (16~Go), calculer
son empreinte & Les processus et les identifiants en mémoire & La perte au
prochain redémarrage \\
10~h~30 & Ouvrir un registre des actions~: heure, auteur, geste, résultat &
La traçabilité de chaque intervention & La contestation ultérieure \\
\end{tableaucours}
\vspace{-0.8em}
\begin{center}{\small\itshape Tableau~0.2. Les gestes de la première heure}\end{center}

Chaque geste a un coût. Débrancher le serveur arrête aussi le partage de
fichiers des agences~; désactiver la passerelle prive les techniciens
d'accès à distance. Le DSI le fait remarquer. La réponse de Mme~NDOUMBE tient
en une phrase~: une heure d'interruption se rattrape, une preuve détruite
ne se reconstitue pas.

\begin{boxobjet}{Capture mémoire de SRV-AD01, 14 mars 2025, 10~h~02}
\textbf{Source~:} contrôleur de domaine SRV-AD01, salle serveurs du siège,
Bafoussam.

\textbf{Pièce~:} fichier \texttt{SRV-AD01\_20250314\_1002.raw}, 16~Go,
copié sur un disque externe neuf, étiqueté et placé sous enveloppe.

\textbf{Empreinte~:} SHA-256 calculée à 10~h~41 et inscrite au registre,
contresignée par M.~KAMGA.

\textbf{Ce qu'elle peut prouver~:} les processus actifs à 10~h~02, les
connexions réseau ouvertes, les traces d'outils exécutés en mémoire.

\textbf{Ce qu'elle ne prouve pas~:} ce qui s'est passé avant le dernier
redémarrage de la machine, ni l'identité d'une personne.
\end{boxobjet}

% ============================================================================
\section{L'enquête : ce qu'elle établira}
\label{sec:p00-enquete}
% ============================================================================

L'après-midi même, NYANGA CRÉDIT mandate le cabinet SANKOFA EXPERTISE, à
Yaoundé, pour conduire la réponse à l'incident~; Mme~MBARGA dirige la
mission. Le lendemain, l'établissement dépose plainte au parquet de
Bafoussam. Le 17~mars, le procureur de la République réquisitionne un expert
indépendant, M.~ABOUBAKAR, expert près la cour d'appel de l'Ouest. Deux
enquêtes avancent désormais en parallèle~: celle du cabinet, au service de
l'établissement, et celle de l'expert, au service de la justice. Elles
examinent les mêmes pièces et n'ont pas les mêmes obligations.

Trois semaines plus tard, les deux enquêtes établissent la même séquence.

\begin{boxscenario}{Opération NYANGA. Ce que l'enquête établit}
\textbf{Constatations.}
\begin{itemize}
    \item Le 2~mars à 23~h~48, le journal de la passerelle enregistre une
          première session du compte de service \texttt{svc\_backup}, depuis
          l'adresse 203.0.113.18, sans échec d'authentification préalable.
    \item Le 9~mars à 1~h~37, un fichier \texttt{lsass.dmp}, copie de la
          mémoire du processus qui gère les authentifications Windows, est
          créé dans un dossier temporaire de SRV-AD01~; la capture mémoire
          du 14~mars montre encore un outil d'accès à distance actif et une
          connexion vers 203.0.113.47.
    \item Le 12~mars, le journal de SRV-CPT01 montre la lecture du fichier
          de configuration qui contient la clé d'accès à l'interface de
          paiement.
    \item Le 14~mars, de 1~h~12 à 3~h~47, la capture réseau montre 3,2~Go
          envoyés de SRV-FIC01 vers 203.0.113.47.
    \item La passerelle fonctionnait avec une version de micrologiciel
          antérieure au correctif publié le 15~décembre 2024 pour la faille
          VULN-2024-4471, qui permet de contourner l'authentification.
\end{itemize}

\textbf{Analyse} (confiance élevée). L'attaquant est entré par la
passerelle, en exploitant VULN-2024-4471 pour ouvrir une session sous un
compte de service~; il a récupéré des identifiants d'administration en
mémoire, lu la clé de l'interface de paiement, copié des fichiers, puis
déclenché les virements.

\textbf{Ce qui reste \NEK{}.} L'identité des auteurs. Les adresses
203.0.113.18 et 203.0.113.47 désignent des machines, pas des personnes.
\end{boxscenario}

La dernière constatation change la nature de l'affaire. Le correctif était
disponible depuis 89~jours le jour de l'attaque, et depuis 77~jours lors du
premier accès~; le 19~décembre, la faille avait été inscrite au catalogue
des vulnérabilités exploitées tenu par l'agence américaine CISA. L'enquête,
partie d'un virement frauduleux, aboutit à une faiblesse connue, publiée et
corrigeable. Pourquoi n'a-t-elle pas été corrigée~? Parce que la passerelle
ne figurait pas dans l'inventaire de l'outil de balayage de vulnérabilités~:
personne ne l'a jamais analysée.

% ============================================================================
\section{Deux disciplines, une même affaire}
\label{sec:p00-deux}
% ============================================================================

L'affaire NYANGA fait se rejoindre deux disciplines que l'on enseigne
souvent séparément.

\begin{tableaucours}{L{3.2cm}L{4.6cm}L{4.6cm}}{%
    \tableauHeader{} &
    \tableauHeader{Investigation numérique} &
    \tableauHeader{Analyse de vulnérabilités}
}
Point de départ & Une trace, un incident & Une faiblesse connue, un actif \\
Question & Que s'est-il passé, comment, depuis quand~? & Où sommes-nous
exposés, et que corriger d'abord~? \\
Matière & Images de disque, mémoire, captures, journaux & Inventaires,
résultats de balayage, bulletins, scores \\
Produit & Un rapport qui peut servir de preuve & Un plan de correction
suivi \\
Temps & Après l'incident & Avant l'incident, en continu \\
\end{tableaucours}
\vspace{-0.8em}
\begin{center}{\small\itshape Tableau~0.3. Deux disciplines}\end{center}

Elles se nourrissent l'une l'autre. L'enquête révèle les faiblesses que
l'attaquant a réellement empruntées~; elle dit au programme de gestion des
vulnérabilités ce qu'il a manqué. Le programme, en retour, réduit le nombre
d'enquêtes à conduire et fournit aux enquêteurs l'inventaire et l'historique
des versions qui leur permettent de dater une exposition. À Douala, un autre
acteur du cours, SOLITRANS LOGISTIQUE, montrera dans la partie~IV comment on
construit un tel programme sur 1\,380 actifs.

\begin{boxinfo}
\textbf{Ce que le cours ne traite pas.} L'exploitation des failles relève
de l'UE MSN-2S3-SEC03 (Ethical Hacking et Pentesting)~: ce cours apprend à
identifier, coter et faire corriger une vulnérabilité, pas à s'en servir.
La protection des données personnelles exposées par l'exfiltration relève de
l'UE MSN-2S3-JUR01. L'audit et la gestion du risque relèvent de l'UE
MSN-2S3-SEC08.
\end{boxinfo}

% ============================================================================
\section{Les questions du cours}
\label{sec:p00-questions}
% ============================================================================

Chaque chapitre répond à l'une des questions que la cellule de crise s'est
posées, ou aurait dû se poser.

\begin{tableaucours}{L{7.4cm}L{4.8cm}}{%
    \tableauHeader{Question de la cellule} &
    \tableauHeader{Chapitre}
}
Qu'est-ce qu'enquêter sur des traces numériques~? &
Chapitre~\ref{chap:investigation} \\
Comment faire pour que ce que l'on trouve serve de preuve~? &
Chapitre~\ref{chap:preuve} \\
Comment copier un disque ou une mémoire sans les altérer~? &
Chapitre~\ref{chap:acquisition} \\
Que contient le disque du serveur comptable~? & Chapitre~\ref{chap:disque} \\
Que révèle la mémoire du contrôleur de domaine~? &
Chapitre~\ref{chap:memoire} \\
Qu'est-il sorti, et vers où~? & Chapitre~\ref{chap:reseau} \\
Dans quel ordre les choses se sont-elles passées~? &
Chapitre~\ref{chap:chronologie} \\
Qui a le droit de faire quoi, et selon quelles règles~? &
Chapitre~\ref{chap:normes-droit} \\
Qu'est-ce qu'une vulnérabilité, et comment la coter~? &
Chapitre~\ref{chap:vulnerabilite} \\
Comment les trouver sans attaquer~? & Chapitre~\ref{chap:identifier} \\
Comment décider de ce qu'on corrige d'abord~? & Chapitre~\ref{chap:gerer} \\
Comment remonter de la trace à la faille~? &
Chapitre~\ref{chap:artefact-vuln} \\
Et si l'attaquant a effacé ses traces, ou travaille dans le nuage~? &
Chapitre~\ref{chap:antiforensique} \\
Comment conduire l'enquête de bout en bout~? & Chapitre~\ref{chap:enquete} \\
Comment l'écrire, et la défendre devant un juge~? &
Chapitre~\ref{chap:rapport} \\
Comment mesurer ce que l'enquête ne sait pas~? &
Chapitre~\ref{chap:auditologie} \\
\end{tableaucours}
\vspace{-0.8em}
\begin{center}{\small\itshape Tableau~0.4. Les questions et les chapitres}\end{center}

% ============================================================================
\section{Les métiers de l'affaire}
\label{sec:p00-metiers}
% ============================================================================

Quatre métiers ont travaillé sur l'affaire NYANGA. Le cours les suit
chapitre après chapitre.

\begin{boxmetier}{RSSI}
Mme~NDOUMBE a préservé avant d'analyser, ouvert le registre des actions et
refusé de coter une hypothèse comme un fait. Son rôle, dans les jours qui
suivent, est d'arbitrer entre la reprise de l'activité et la conservation
des preuves, et de faire corriger la faille.
\end{boxmetier}

\begin{boxmetier}{Analyste en réponse à incident}
Mme~MBARGA et son équipe contiennent l'incident, recueillent les pièces,
reconstituent la séquence et remettent à l'établissement un rapport qui
sert à décider. Ils travaillent vite, mais chaque pièce qu'ils touchent
peut finir devant un tribunal.
\end{boxmetier}

\begin{boxmetier}{Application OPJ/Expert Judiciaire}
L'adjudant-chef NANA, officier de police judiciaire, conduit l'enquête
pénale~; M.~ABOUBAKAR, réquisitionné par le procureur, examine les pièces
pour la justice. L'expert ne répond qu'aux questions de sa mission, et
chacune de ses conclusions doit pouvoir être refaite par un autre expert.
\end{boxmetier}

\begin{boxmetier}{Analyste en gestion des vulnérabilités}
Personne n'occupait ce rôle à NYANGA CRÉDIT. C'est la leçon de l'affaire~:
un actif hors de l'inventaire est un actif que personne ne corrige.
\end{boxmetier}

% ============================================================================
\section{Cinq réflexes}
\label{sec:p00-reflexes}
% ============================================================================

\begin{enumerate}
    \item Isoler sans éteindre, puis préserver dans l'ordre de la
          volatilité.
    \item Consigner chaque geste~: heure, auteur, action, résultat,
          empreinte.
    \item Séparer la constatation de l'analyse, et coter \NEK{} ce que les
          pièces ne permettent pas de trancher.
    \item Remonter jusqu'à la faiblesse exploitée~: une enquête qui s'arrête
          aux auteurs laisse la porte ouverte.
    \item Vérifier que la correction est faite, sur la machine elle-même, et
          pas sur la foi d'un courriel.
\end{enumerate}

\separateur

\begin{boxresume}
\textbf{Ce qu'il faut retenir de ce chapitre~:}
\begin{itemize}
    \item Les gestes de la première heure décident des preuves
          disponibles~: isoler, préserver dans l'ordre de la volatilité,
          consigner.
    \item Une constatation se rattache à une source~; une analyse porte un
          niveau de confiance. Une hypothèse reste \NEK{} tant qu'aucune
          pièce ne la tranche.
    \item À NYANGA CRÉDIT, l'enquête remonte du virement frauduleux à une
          faille corrigeable depuis 89~jours, absente de l'inventaire.
    \item L'investigation numérique et l'analyse de vulnérabilités se
          répondent~: l'une explique l'incident, l'autre empêche le
          suivant.
\end{itemize}
\end{boxresume}

\begin{boxinfo}
\textbf{Points de contrôle.} Vous avez compris ce prologue si vous savez~:
\begin{itemize}
    \item citer trois gestes qui préservent la preuve et deux qui la
          détruisent~;
    \item réécrire une phrase de rapport en séparant la constatation et
          l'analyse~;
    \item coter une conclusion \OUI{}, \NON{} ou \NEK{} et justifier la
          cote par une source.
\end{itemize}
\end{boxinfo}

% ============================================================================
\section*{Exercices}
% ============================================================================

\begin{boxexo}[title={\ding{45}\quad Exercice~0.1. Préserver ou détruire \niveaubase}]
Pour chacun des gestes suivants, dites s'il préserve ou s'il détruit de la
preuve, et laquelle.
\begin{enumerate}[(a)]
    \item Lancer un antivirus complet sur le serveur de fichiers.
    \item Photographier l'écran du serveur avant toute action.
    \item Redémarrer le contrôleur de domaine pour « couper l'attaquant ».
    \item Exporter les journaux de la passerelle sur un support neuf et en
          calculer l'empreinte.
    \item Changer immédiatement le mot de passe de \texttt{svc\_backup}
          depuis le serveur compromis.
\end{enumerate}
\end{boxexo}

\begin{boxexo}[title={\ding{45}\quad Exercice~0.2. Constatation ou analyse \niveaumoyen}]
Le premier compte rendu du DSI contient les phrases suivantes. Pour
chacune, dites si c'est une constatation, une analyse ou une hypothèse, et
proposez une cote \OUI{}, \NON{} ou \NEK{} avec sa source.
\begin{enumerate}[(a)]
    \item « Les virements ont été faits depuis l'interface de paiement avec
          notre clé. »
    \item « Un employé a forcément donné son mot de passe. »
    \item « Le serveur de fichiers a envoyé 3,2~Go vers l'extérieur. »
    \item « Les fichiers volés sont ceux des clients. »
    \item « L'attaque vient d'un groupe étranger. »
\end{enumerate}
\end{boxexo}

\begin{boxexo}[title={\ding{45}\quad Exercice~0.3. Le compte des jours \niveaumoyen}]
\begin{enumerate}[(a)]
    \item Vérifiez les deux durées données dans le chapitre~: 77~jours entre
          le correctif et le premier accès, 89~jours entre le correctif et
          l'exfiltration.
    \item La faille a été inscrite au catalogue des vulnérabilités
          exploitées le 19~décembre 2024. Que change cette inscription pour
          un établissement qui suit ce catalogue~?
    \item Donnez deux raisons pour lesquelles un actif peut échapper à
          l'inventaire d'un outil de balayage.
\end{enumerate}
\end{boxexo}

\begin{boxexo}[title={\ding{45}\quad Exercice~0.4. Deux enquêtes, deux obligations \niveauexpert}]
Le cabinet SANKOFA EXPERTISE et l'expert judiciaire M.~ABOUBAKAR examinent
les mêmes pièces.
\begin{enumerate}[(a)]
    \item Pour chacun, dites à qui il rend compte et quelle question il doit
          trancher.
    \item Le directeur général demande au cabinet de ne pas mentionner le
          retard de correction dans son rapport. Que doit faire Mme~MBARGA~?
          L'expert judiciaire serait-il tenu par la même demande~?
    \item Rédigez, pour la cellule de crise, une note de cinq lignes qui
          distingue ce qui est établi, ce qui est réfuté et ce qui reste
          \NEK{} au soir du 14~mars.
\end{enumerate}
\end{boxexo}

\bigskip

\begin{boxinfo}
\textbf{Pour aller plus loin.}
\begin{itemize}
    \item La RFC~3227~\cite{rfc3227}, pour l'ordre de volatilité.
    \item Le NIST SP~800-86~\cite{nist80086}, pour l'intégration de la
          forensique à la réponse aux incidents.
    \item Le catalogue KEV de la CISA~\cite{cisakev}.
\end{itemize}
\end{boxinfo}

\begin{boxinfo}
\textbf{Transition.} Avant de toucher aux pièces de l'affaire, il faut
savoir ce qu'est une enquête numérique~: son objet, ses principes, ses
étapes. C'est l'objet du Chapitre~\ref{chap:investigation}.
\end{boxinfo}


    % PARTIE 1 · FONDEMENTS (ISO/IEC 27037, 27043 ; RFC 3227)
    \part{Fondements : Enquêter sur des Traces Numériques}
    \label{part:fondements}

        % ============================================================================
% Fichier  : partie1/P01_C01_investigation_numerique.tex
% Titre    : L'Investigation Numérique
% Auteur   : MINKA MI NGUIDJOÏ Thierry Emmanuel
% Version  : 1.0 -- novembre 2026
% Corrigés : corriges/Corr_C01.tex
% ============================================================================
% Sources : ISO/IEC 27037:2012 (principes, paraphrasés), ISO/IEC 27043:2015
% (classes de processus, paraphrasées), NIST SP 800-86 (quatre phases),
% NIST SP 800-61 r2, Locard (1920), Casey (2004), Heuer (1999, analyse des
% hypothèses concurrentes).
% Faits nouveaux (bible) : hypothèses H1 à H4 de la cellule du 15/03/2025.
% ============================================================================

\chapter{L'Investigation Numérique}
\label{chap:investigation}
\chaptermark{L'investigation numérique}

\epigraph{%
    « Une enquête ne cherche pas à confirmer une idée.\\
    Elle cherche ce qui pourrait la réfuter. »%
}{Principe du cours}

\niveauChapitre{Fondamental}{%
    Lecture du Chapitre~\ref{chap:prologue}
}

\begin{boxobjectifs}
\begin{itemize}
    \item Définir l'investigation numérique, la trace et la preuve
          numériques.
    \item Énoncer les principes qui rendent un examen acceptable~:
          traçabilité, répétabilité, reproductibilité, justification.
    \item Décrire les phases d'une enquête selon le NIST SP~800-86 et les
          classes de processus de l'ISO/IEC~27043.
    \item Distinguer enquête judiciaire, enquête interne et réponse à
          incident, et leurs exigences de preuve.
    \item Conduire un raisonnement par hypothèses concurrentes et coter
          chaque hypothèse.
\end{itemize}
\end{boxobjectifs}

\bigskip

% ============================================================================
\section*{Avant de commencer : quatre hypothèses sur la table}
% ============================================================================

Bafoussam, samedi 15 mars 2025, 8~heures. Mme~MBARGA, de SANKOFA EXPERTISE,
ouvre la première réunion de la mission en demandant à chacun d'écrire, sur
une feuille, comment il pense que l'attaquant est entré. Elle recueille
quatre réponses.

\begin{itemize}
    \item H1~: par la passerelle d'accès à distance, avec un compte volé ou
          une faille.
    \item H2~: par un poste d'agence infecté par un courriel piégé, puis par
          le réseau interne.
    \item H3~: par un employé du siège qui aurait agi de l'intérieur,
          depuis un poste du siège.
    \item H4~: par le prestataire de l'interface de paiement, dont la clé
          aurait fuité de son côté.
\end{itemize}

Le DSI tient pour H3, la rumeur du vendredi. Le directeur général tient pour
H4, qui met l'établissement hors de cause. Mme~MBARGA ne tient pour aucune~:
elle écrit les quatre au tableau et annonce que l'enquête consistera à
chercher, pour chacune, la pièce qui la réfuterait.

\begin{boxalert}
\textbf{Violation typique.} Choisir l'hypothèse la plus probable dès le
premier jour, puis ne chercher que les traces qui la confirment.

\smallskip
\textbf{Règle~:} on garde ouvertes toutes les hypothèses compatibles avec les
constatations, et l'on cherche en priorité la pièce capable d'en éliminer
une. Une hypothèse survit parce qu'on n'a pas réussi à la réfuter, pas parce
qu'on a trouvé des indices en sa faveur.
\end{boxalert}

\separateur

% ============================================================================
\section{Trace, Preuve et Investigation}
\label{sec:c01-definitions}
% ============================================================================

\begin{boxdef}
\textbf{\termecle{Trace numérique}}~: toute donnée produite par un système
d'information, conservée sur un support ou en transit, qui garde la marque
d'une action ou d'un état passé~: une ligne de journal, une date de
modification, un fichier supprimé mais récupérable, une connexion en mémoire.

\smallskip
\textbf{\termecle{Preuve numérique}}~: trace recueillie, préservée et
examinée de façon à pouvoir être présentée à un décideur (un juge, une
direction, un régulateur) et soutenir ou réfuter une proposition de fait.

\smallskip
\textbf{\termecle{Investigation numérique}}~: démarche méthodique qui
identifie, recueille, préserve, examine et analyse des traces numériques
pour répondre à des questions de fait, et en rend compte de façon
vérifiable. On parle aussi de \textit{digital forensics}.
\end{boxdef}

Toute trace n'est pas une preuve. Une trace devient preuve quand trois
conditions sont réunies~: on sait d'où elle vient, on peut montrer qu'elle
n'a pas changé depuis son recueil, et on peut expliquer ce qu'elle signifie.
Le Chapitre~\ref{chap:preuve} traite des deux premières~; la troisième
occupe toute la partie~II.

\subsection{Le principe d'échange}

Edmond Locard, fondateur du laboratoire de police technique de Lyon, a
formulé au début du \textsc{xx}\textsuperscript{e}~siècle l'idée que tout
contact laisse une trace~\cite{locard1920}. Le numérique la vérifie d'une
manière particulière. Un attaquant qui ouvre une session laisse une entrée
dans le journal de la passerelle~; qui exécute un programme laisse une trace
dans la mémoire, dans les journaux du système et souvent dans des fichiers
de préchargement~; qui copie des données laisse un volume sortant dans la
capture du lien Internet. Mais l'échange joue aussi dans l'autre sens~:
l'enquêteur qui examine une machine allumée y laisse ses propres traces. Le
principe d'échange oblige donc l'enquêteur autant que l'attaquant.

\subsection{Ce qui distingue la trace numérique}

\begin{tableaucours}{L{3.6cm}L{8.6cm}}{%
    \tableauHeader{Propriété} &
    \tableauHeader{Conséquence pour l'enquête}
}
Copiable à l'identique & On travaille sur une copie, et l'empreinte
cryptographique prouve que la copie égale l'original \\
Modifiable sans marque visible & Une date ou un journal se changent sans
rature~; seule la préservation prouve l'intégrité \\
Volatile & La mémoire vive et les connexions disparaissent à l'extinction~;
l'ordre de recueil compte \\
Abondante & Un serveur contient des millions de traces~; il faut des
questions précises pour savoir quoi chercher \\
Dépendante de son contexte & Une heure n'a de sens qu'avec son fuseau et la
dérive de l'horloge qui l'a produite \\
\end{tableaucours}
\vspace{-0.8em}
\begin{center}{\small\itshape Tableau~1.1. Propriétés de la trace numérique}\end{center}

% ============================================================================
\section{Les Principes d'un Examen Acceptable}
\label{sec:c01-principes}
% ============================================================================

La norme ISO/IEC~27037~\cite{iso27037_2012} fixe des lignes directrices pour
identifier, recueillir, acquérir et préserver la preuve numérique. Elle
retient quatre principes, que l'on paraphrase ici.

\begin{tableaucours}{L{3.2cm}L{5.4cm}L{3.6cm}}{%
    \tableauHeader{Principe} &
    \tableauHeader{Sens} &
    \tableauHeader{Comment le prouver}
}
Traçabilité (\textit{auditability}) & Un tiers peut reconstituer chaque geste
de l'enquêteur & Registre des actions, notes horodatées \\
Répétabilité & La même personne, avec les mêmes outils sur la même pièce,
obtient le même résultat & Empreintes identiques, versions d'outils notées \\
Reproductibilité & Une autre personne, avec d'autres outils, obtient un
résultat compatible & Méthode décrite, pièces conservées \\
Justification & Chaque choix de méthode ou d'outil peut être expliqué &
Raisons écrites dans le dossier \\
\end{tableaucours}
\vspace{-0.8em}
\begin{center}{\small\itshape Tableau~1.2. Quatre principes (d'après l'ISO/IEC~27037)}\end{center}

À ces quatre principes s'ajoutent deux règles de conduite que l'on retrouve
dans les guides de bonnes pratiques policiers, comme celui de l'ACPO
britannique~\cite{acpo2012}. La première est l'intervention minimale~: on ne
modifie pas une pièce, et si l'on y est contraint, par exemple pour capturer
la mémoire d'une machine allumée, on le fait en personne compétente et l'on
explique pourquoi. La seconde est la proportionnalité~: on recueille ce qui
répond aux questions de l'enquête, pas tout ce qui est techniquement
accessible.

\begin{boxinfo}
\textbf{Répétable, reproductible.} La distinction n'est pas un détail. Un
résultat répétable mais non reproductible dépend peut-être d'une
particularité de l'outil. Quand deux outils indépendants donnent la même
liste de fichiers supprimés, la conclusion gagne en confiance~; quand ils
divergent, la divergence elle-même devient un objet d'examen.
\end{boxinfo}

% ============================================================================
\section{Les Phases d'une Enquête}
\label{sec:c01-phases}
% ============================================================================

\subsection{Quatre phases selon le NIST}

Le guide NIST SP~800-86~\cite{nist80086} décrit l'intégration des techniques
forensiques à la réponse aux incidents en quatre phases.

\begin{center}
\begin{tikzpicture}[
    etape/.style = {draw=CouleurBordDef, fill=CouleurDef, rounded corners=3pt,
        minimum width=2.6cm, minimum height=1.0cm, align=center,
        font=\small\bfseries\color{CouleurPrimaire}},
    arrow/.style = {-{Stealth[length=4pt]}, thick, color=CouleurSecondaire}
]
\node[etape] (a) at (0,0)    {Collecte};
\node[etape] (b) at (3.3,0)  {Examen};
\node[etape] (c) at (6.6,0)  {Analyse};
\node[etape] (d) at (9.9,0)  {Rapport};
\draw[arrow] (a) -- (b); \draw[arrow] (b) -- (c); \draw[arrow] (c) -- (d);
\node[font=\footnotesize, align=center, text width=2.6cm] at (0,-1.1) {Support, données, intégrité};
\node[font=\footnotesize, align=center, text width=2.6cm] at (3.3,-1.1) {Extraire ce qui est pertinent};
\node[font=\footnotesize, align=center, text width=2.6cm] at (6.6,-1.1) {Relier, dater, conclure};
\node[font=\footnotesize, align=center, text width=2.6cm] at (9.9,-1.1) {Rendre compte, recommander};
\end{tikzpicture}
\end{center}
\begin{center}{\small\itshape Figure~1.1. Les quatre phases du NIST SP~800-86}\end{center}

La collecte identifie et recueille les supports et les données en préservant
leur intégrité. L'examen extrait, de la masse recueillie, les éléments
pertinents~: on décode un système de fichiers, on filtre des journaux. L'analyse
relie ces éléments entre eux pour répondre aux questions posées. Le rapport
en rend compte et formule des recommandations.

\subsection{Cinq classes de processus selon l'ISO/IEC 27043}

La norme ISO/IEC~27043~\cite{iso27043_2015} propose un modèle plus large,
qui commence avant l'incident. On en retient cinq classes de processus,
paraphrasées.

\begin{tableaucours}{L{3.4cm}L{8.8cm}}{%
    \tableauHeader{Classe} &
    \tableauHeader{Ce qu'elle couvre}
}
Préparation (\textit{readiness}) & Ce que l'organisation fait avant tout
incident pour qu'une enquête soit possible~: journaux conservés, horloges
synchronisées, procédure de première intervention \\
Initialisation & Le déclenchement~: détection, première évaluation,
décision d'enquêter, mandat \\
Acquisition & L'identification, le recueil, l'acquisition, le transport et
le stockage des pièces \\
Investigation & L'examen, l'analyse, l'interprétation, le rapport, la
présentation et la clôture \\
Processus concurrents & Ce qui accompagne toute l'enquête~: documentation,
chaîne de possession, préservation, autorisations \\
\end{tableaucours}
\vspace{-0.8em}
\begin{center}{\small\itshape Tableau~1.3. Classes de processus (d'après l'ISO/IEC~27043)}\end{center}

\begin{boxscenario}{Opération NYANGA. La préparation qui manquait}
\textbf{Constatations.} La sonde réseau installée en 2024 conserve 48~heures
de capture~; les journaux de la passerelle sont gardés 60~jours~; les
horloges des serveurs sont synchronisées sur un serveur de temps interne~;
l'établissement n'a aucune procédure écrite de première intervention, et
aucun journal des postes d'agence n'est centralisé.

\textbf{Analyse} (confiance élevée). La préparation existait pour le siège,
pas pour les agences. Si l'attaquant était passé par un poste d'agence (H2),
les traces auraient été dispersées sur vingt-deux sites et peut-être déjà
écrasées. La réaction de Mme~NDOUMBE a compensé l'absence de procédure~; un
autre responsable, ce vendredi-là, aurait pu laisser le DSI redémarrer.
\end{boxscenario}

% ============================================================================
\section{Trois Cadres d'Enquête}
\label{sec:c01-cadres}
% ============================================================================

Les mêmes pièces peuvent servir trois enquêtes différentes, qui n'ont ni les
mêmes questions, ni les mêmes exigences, ni les mêmes destinataires.

\begin{tableaucours}{L{2.6cm}L{3.2cm}L{3.2cm}L{3.2cm}}{%
    \tableauHeader{} &
    \tableauHeader{Enquête judiciaire} &
    \tableauHeader{Enquête interne} &
    \tableauHeader{Réponse à incident}
}
Qui la mandate & Le procureur, le juge, l'OPJ & La direction de
l'organisation & L'organisation victime \\
Question & Une infraction a-t-elle été commise, par qui~? & Un salarié
a-t-il manqué à ses obligations~? & Que s'est-il passé, comment arrêter et
reprendre~? \\
Exigence de preuve & La plus haute~: débat contradictoire & Celle du
contentieux du travail & Suffisante pour décider vite \\
Priorité & La vérité judiciaire & La décision disciplinaire & La reprise
de l'activité \\
Risque principal & Pièce écartée pour vice de forme & Atteinte à la vie
privée du salarié & Preuve détruite par la remise en route \\
\end{tableaucours}
\vspace{-0.8em}
\begin{center}{\small\itshape Tableau~1.4. Trois cadres d'enquête}\end{center}

La réponse à incident, décrite par le guide NIST SP~800-61~\cite{nist80061r2},
a pour objet premier de contenir l'attaque et de rétablir le service. Elle
entre en tension avec la préservation~: remettre en route un serveur,
c'est souvent effacer ce qu'il contenait. La règle pratique est de toujours
travailler comme si l'enquête pouvait finir devant un juge. Un incident qui
paraît interne au matin du vendredi peut devenir une affaire pénale la semaine suivante,
comme à NYANGA CRÉDIT~; les pièces mal recueillies le vendredi ne seront pas
recevables la semaine suivante.

\begin{boxinfo}
\textbf{À NYANGA, les trois cadres coexistent.} SANKOFA EXPERTISE conduit la
réponse à incident pour l'établissement~; l'OPJ et l'expert réquisitionné
conduisent l'enquête judiciaire~; si H3 avait survécu, la direction aurait dû
ouvrir une enquête interne, avec les garanties dues au salarié mis en cause.
Le Chapitre~\ref{chap:normes-droit} précise ce que la loi permet à chacun.
\end{boxinfo}

% ============================================================================
\section{Raisonner par Hypothèses Concurrentes}
\label{sec:c01-hypotheses}
% ============================================================================

Richards Heuer, analyste du renseignement américain, a décrit une méthode
simple pour résister à la tentation de confirmer sa première idée~: l'analyse
des hypothèses concurrentes~\cite{heuer1999}. On liste toutes les hypothèses
plausibles, on liste les pièces, puis on remplit une matrice en notant, pour
chaque couple, si la pièce est compatible ou incompatible avec l'hypothèse.
On retient l'hypothèse qui a le moins d'incompatibilités, pas celle qui a le
plus de compatibilités. Une pièce compatible avec toutes les hypothèses n'aide
pas à choisir, quelle que soit son importance apparente.

\begin{boxscenario}{Opération NYANGA. La matrice du 29 mars}
\begin{tableaucours}{L{5.2cm}cccc}{%
    \tableauHeader{Pièce} &
    \tableauHeader{H1} & \tableauHeader{H2} & \tableauHeader{H3} &
    \tableauHeader{H4}
}
Virements avec la clé de l'établissement & C & C & C & C \\
Session VPN \texttt{svc\_backup} depuis 203.0.113.18, le 2~mars & C & I & I &
I \\
Aucun échec d'authentification avant cette session & C & N & N & N \\
Aucun courriel piégé dans les boîtes des agences sur la période & C & I & C &
C \\
Lecture du fichier de clé sur SRV-CPT01, le 12~mars, par \texttt{adm.kamga}
& C & C & C & I \\
\texttt{adm.kamga} utilisé depuis le contrôleur de domaine, jamais depuis
un poste du siège & C & I & I & N \\
\end{tableaucours}

\smallskip
{\footnotesize C~: compatible~; I~: incompatible~; N~: neutre (ne distingue
pas).}

\medskip
\textbf{Constatations.} H1 n'a aucune incompatibilité~; H2 en a trois, H3
deux, H4 deux. La première ligne, la plus spectaculaire, est compatible avec
tout et ne sert à rien pour choisir.

\textbf{Analyse} (confiance élevée). H1 est retenue~; H2, H3 et H4 sont
incompatibles avec au moins deux pièces. Mais H3 n'est réfutée que telle
qu'elle était formulée~: un employé du siège n'a pas agi depuis un poste du
siège. Cela ne dit pas qu'aucun employé n'a participé~: un employé aurait pu
utiliser la passerelle depuis l'extérieur, comme l'attaquant. Cette question
est distincte de celle du chemin d'entrée, et aucune pièce ne la tranche.
\end{boxscenario}

\begin{boxlecture}{Les quatre hypothèses au 29 mars}
\begin{tableaucours}{L{6.0cm}cL{4.4cm}}{%
    \tableauHeader{Hypothèse} &
    \tableauHeader{Cote} &
    \tableauHeader{Pièce décisive}
}
H1, entrée par la passerelle & \OUI{} & Session du 2~mars, puis usage de
\texttt{adm.kamga} depuis le contrôleur de domaine \\
H2, poste d'agence infecté & \NON{} & Aucune session interne, aucun courriel
piégé \\
H3, employé agissant depuis le siège & \NON{} & Aucune session de
\texttt{adm.kamga} depuis un poste du siège \\
H4, fuite chez le prestataire & \NON{} & Clé lue sur SRV-CPT01 avant les
virements \\
Participation d'un employé, quel que soit le chemin & \NEK{} & Aucune pièce
ne l'établit ni ne l'exclut \\
\end{tableaucours}

\medskip
\textbf{Lecture.} Une proposition établie, trois réfutées, une non
tranchée. Le rapport intermédiaire ne doit pas dire « l'attaque vient de
l'extérieur » sans réserve~: il doit dire que l'entrée par la passerelle est
établie, et que la participation d'un employé n'est ni établie ni exclue.
\end{boxlecture}

% ============================================================================
\section{Les Biais de l'Enquêteur}
\label{sec:c01-biais}
% ============================================================================

L'enquêteur est un instrument de mesure, et cet instrument a des défauts
connus.

\begin{tableaucours}{L{3.4cm}L{4.6cm}L{4.2cm}}{%
    \tableauHeader{Biais} &
    \tableauHeader{Manifestation} &
    \tableauHeader{Parade}
}
Confirmation & On ne cherche que ce qui confirme l'hypothèse préférée &
Matrice des hypothèses concurrentes \\
Ancrage & La première information reçue pèse trop & Noter l'hypothèse
initiale et la réexaminer à chaque pièce \\
Autorité & L'avis du supérieur ou du client devient un fait & Exiger une
source pour chaque constatation \\
Contexte & Savoir qui est suspect oriente la lecture des traces & Examiner
les pièces avant de lire le dossier, quand c'est possible \\
Outil & Ce que l'outil n'affiche pas est tenu pour inexistant & Contrôler
avec un second outil, noter les limites \\
\end{tableaucours}
\vspace{-0.8em}
\begin{center}{\small\itshape Tableau~1.5. Biais et parades}\end{center}

Le biais d'outil mérite une attention particulière. Un logiciel d'analyse
affiche ce qu'il sait décoder. Un format de journal qu'il ne connaît pas, un
système de fichiers qu'il lit mal, une partition qu'il n'a pas montée
n'apparaissent nulle part, et l'absence d'affichage passe pour une absence de
trace. L'enquêteur qui écrit « aucune trace de connexion » doit pouvoir dire
« dans quelles sources, lues avec quels outils ». Sinon, la bonne cote est
\NEK{}, pas \NON{}.

% ============================================================================
\section{Ce que l'Investigation demande à chaque Métier}
\label{sec:c01-metiers}
% ============================================================================

\begin{boxmetier}{Analyste en réponse à incident}
\textbf{Ouvrir une enquête}
\begin{enumerate}
    \item Écrire les questions de l'enquête avant de toucher aux pièces.
    \item Lister toutes les hypothèses compatibles avec les premières
          constatations.
    \item Identifier, pour chacune, la pièce qui la réfuterait.
    \item Préserver ces pièces en priorité.
    \item Tenir le registre des actions dès la première minute.
\end{enumerate}
\end{boxmetier}

\begin{boxmetier}{Application OPJ/Expert Judiciaire}
L'expert ne choisit pas ses questions~: elles figurent dans la réquisition
ou l'ordonnance qui le désigne. S'il découvre un fait hors de sa mission, il
le signale à l'autorité qui l'a désigné, sans l'instruire de lui-même.
\end{boxmetier}

\begin{boxmetier}{RSSI}
La préparation est son affaire~: durée de conservation des journaux,
synchronisation des horloges, procédure de première intervention connue de
l'équipe de permanence. À NYANGA, il manquait la dernière.
\end{boxmetier}

\separateur

\begin{boxresume}
\textbf{Ce qu'il faut retenir de ce chapitre~:}
\begin{itemize}
    \item Une trace devient preuve quand on sait d'où elle vient, qu'on peut
          montrer qu'elle n'a pas changé, et qu'on peut expliquer ce qu'elle
          signifie.
    \item Quatre principes rendent un examen acceptable~: traçabilité,
          répétabilité, reproductibilité, justification.
    \item Le NIST décrit quatre phases (collecte, examen, analyse, rapport)~;
          l'ISO/IEC~27043 ajoute la préparation, avant l'incident.
    \item Enquête judiciaire, enquête interne et réponse à incident n'ont ni
          les mêmes questions ni les mêmes exigences~; on travaille toujours
          comme si l'affaire pouvait aller devant un juge.
    \item À NYANGA, la matrice du 29~mars établit H1, réfute H2, H3 et H4,
          et laisse \NEK{} la participation éventuelle d'un employé.
\end{itemize}
\end{boxresume}

\begin{boxinfo}
\textbf{Points de contrôle.} Vous avez compris ce chapitre si vous savez~:
\begin{itemize}
    \item expliquer la différence entre répétabilité et reproductibilité~;
    \item placer une action d'enquête dans l'une des phases du NIST~;
    \item remplir une matrice d'hypothèses concurrentes et en tirer une
          cote pour chaque hypothèse.
\end{itemize}
\end{boxinfo}

% ============================================================================
\section*{Exercices}
% ============================================================================

\begin{boxexo}[title={\ding{45}\quad Exercice~1.1. Trace ou preuve \niveaubase}]
Dites, pour chaque élément, s'il s'agit d'une simple trace ou d'une preuve
au sens du chapitre, et ce qui lui manque le cas échéant.
\begin{enumerate}[(a)]
    \item Une capture d'écran du journal de la passerelle, envoyée par
          WhatsApp par un technicien.
    \item L'export du même journal, sur support neuf, avec son empreinte
          SHA-256 inscrite au registre.
    \item Un fichier \texttt{lsass.dmp} trouvé sur SRV-AD01, dont personne
          n'a noté ni l'heure ni l'emplacement exacts.
    \item La capture mémoire de SRV-AD01 décrite au Chapitre~\ref{chap:prologue}.
\end{enumerate}
\end{boxexo}

\begin{boxexo}[title={\ding{45}\quad Exercice~1.2. Les phases et les classes \niveaubase}]
Rattachez chaque action à une phase du NIST SP~800-86 et à une classe de
processus de l'ISO/IEC~27043~: (a) synchroniser les horloges des serveurs
en 2024~; (b) débrancher le serveur de fichiers le 14~mars~; (c) extraire
du journal de la passerelle les sessions de \texttt{svc\_backup}~; (d)
relier ces sessions aux connexions de \texttt{adm.kamga}~; (e) remettre le
rapport à la direction.
\end{boxexo}

\begin{boxexo}[title={\ding{45}\quad Exercice~1.3. Une pièce qui change tout \niveaumoyen}]
On découvre le 2~avril que \texttt{adm.kamga} a ouvert, le 11~mars, une
session depuis un poste du siège, à 14~h~20, pendant les heures de bureau.
\begin{enumerate}[(a)]
    \item Ajoutez cette pièce à la matrice du 29~mars.
    \item Les cotes des hypothèses changent-elles~? Justifiez.
    \item Quelle vérification feriez-vous avant de modifier une cote~?
\end{enumerate}
\end{boxexo}

\begin{boxexo}[title={\ding{45}\quad Exercice~1.4. Trois cadres pour une affaire \niveauexpert}]
Six mois après l'incident, NYANGA CRÉDIT découvre qu'un agent comptable a
consulté, sans raison de service, des dossiers de clients figurant parmi les
fichiers exfiltrés.
\begin{enumerate}[(a)]
    \item Quels cadres d'enquête peuvent s'ouvrir~? Qui mandate chacun~?
    \item Pour chacun, formulez la question de fait et l'exigence de preuve.
    \item Quelles précautions prendre pour que les pièces recueillies dans
          l'enquête interne restent utilisables dans une enquête judiciaire~?
\end{enumerate}
\end{boxexo}

\bigskip

\begin{boxinfo}
\textbf{Pour aller plus loin.}
\begin{itemize}
    \item E.~Casey, \textit{Digital Evidence and Computer
          Crime}~\cite{casey2004}, pour les fondements.
    \item Le NIST SP~800-86~\cite{nist80086} et l'ISO/IEC~27043~\cite{iso27043_2015}.
    \item R.~Heuer, \textit{Psychology of Intelligence
          Analysis}~\cite{heuer1999}, chapitre sur les hypothèses
          concurrentes.
\end{itemize}
\end{boxinfo}

\begin{boxinfo}
\textbf{Transition.} Une hypothèse se tranche avec des pièces, et une pièce
ne vaut que si l'on peut prouver qu'elle n'a pas changé. C'est l'objet du
Chapitre~\ref{chap:preuve}~: la preuve numérique et la chaîne de possession.
\end{boxinfo}

        % ============================================================================
% Fichier  : partie1/P01_C02_preuve_chaine_possession.tex
% Titre    : La Preuve Numérique et la Chaîne de Possession
% Auteur   : MINKA MI NGUIDJOÏ Thierry Emmanuel
% Version  : 1.0 -- novembre 2024
% Corrigés : corriges/Corr_C02.tex
% ============================================================================
% Sources : ISO/IEC 27037:2012 (paraphrasée) ; loi n° 2010/012, art. 53 et 54
% (texte officiel, cité) ; FIPS 180-4 et état de l'art sur les collisions
% MD5/SHA-1 (présenté sans citation d'article).
% Faits nouveaux (bible) : chaîne de possession de la capture mémoire
% (scellé NC-01) ; trou de 17 h 40 dans la nuit du 14 au 15 mars ; scellés
% NC-01 à NC-06.
% ============================================================================

\chapter{La Preuve Numérique et la Chaîne de Possession}
\label{chap:preuve}
\chaptermark{Preuve et chaîne de possession}

\epigraph{%
    « Une pièce dont on ne peut pas dire où elle était\\
    ne vaut pas mieux qu'une pièce perdue. »%
}{Principe du cours}

\niveauChapitre{Fondamental}{%
    Maîtrise du Chapitre~\ref{chap:investigation}
}

\begin{boxobjectifs}
\begin{itemize}
    \item Distinguer la recevabilité d'une preuve et sa force probante.
    \item Expliquer ce qu'une empreinte cryptographique prouve, et ce
          qu'elle ne prouve pas.
    \item Choisir une fonction de hachage adaptée et justifier ce choix.
    \item Tenir une chaîne de possession complète, du recueil à la
          restitution.
    \item Identifier une rupture de la chaîne et en mesurer les
          conséquences.
    \item Appliquer les articles~53 et~54 de la loi n°~2010/012 à une
          perquisition portant sur des données.
\end{itemize}
\end{boxobjectifs}

\bigskip

% ============================================================================
\section*{Avant de commencer : un disque sur un bureau}
% ============================================================================

Bafoussam, samedi 15 mars 2025, 9~h~30. Mme~MBARGA demande où se trouve le
disque externe qui contient la capture mémoire du contrôleur de domaine. Il
est sur le bureau de M.~KAMGA, dans une enveloppe kraft fermée par un ruban
adhésif. Personne ne l'a mis sous clé pendant la nuit. La porte du bureau
du DSI ne ferme pas, et l'équipe de ménage passe le vendredi soir.

L'empreinte de la capture avait été calculée la veille à 10~h~41 et inscrite
au registre. Mme~MBARGA la recalcule devant M.~KAMGA et Mme~NDOUMBE~: elle
est identique. Le contenu n'a donc pas changé. Mais l'avocat de la défense,
dans dix-huit mois, posera une autre question~: pendant ces dix-sept heures,
qui pouvait toucher ce disque~?

\begin{boxalert}
\textbf{Violation typique.} Croire qu'une empreinte identique suffit à
défendre une pièce, et négliger la trace écrite de ceux qui l'ont détenue.

\smallskip
\textbf{Règle~:} l'empreinte prouve que le contenu n'a pas changé entre deux
calculs~; la chaîne de possession prouve qui détenait la pièce, où, et dans
quelles conditions. Les deux se complètent et aucune ne remplace l'autre.
\end{boxalert}

\separateur

% ============================================================================
\section{Recevable et Probante}
\label{sec:c02-recevable}
% ============================================================================

Deux questions se posent devant toute pièce, et elles ne se confondent pas.

\begin{boxdef}
\textbf{\termecle{Recevabilité}}~: aptitude d'une pièce à être versée au
débat. Une pièce obtenue en violation de la loi, par exemple par un accès
non autorisé à un système, peut être écartée.

\smallskip
\textbf{\termecle{Force probante}}~: poids que le décideur accorde à une
pièce recevable. Elle dépend de la fiabilité de la source, de l'intégrité de
la pièce, de la qualité de la méthode et de la clarté de son interprétation.
\end{boxdef}

En matière pénale, le droit camerounais, comme la plupart des droits issus
de la tradition civiliste, admet en principe la preuve par tout moyen et
laisse le juge apprécier librement la valeur des pièces [disposition du Code
de procédure pénale à vérifier par le lecteur sur le texte officiel]. Ce
principe ouvre la porte à la preuve numérique, mais il ne la rend pas
automatiquement convaincante. Une pièce numérique contestée sur son
intégrité ou sur sa provenance sera recevable et peu probante~: le juge la
lira, puis lui accordera peu de poids. L'enquêteur travaille donc moins pour
la recevabilité, que la loi accorde largement, que pour la force probante,
qu'il doit construire.

\begin{tableaucours}{L{4.0cm}L{8.2cm}}{%
    \tableauHeader{Ce qui fonde la force probante} &
    \tableauHeader{Ce que l'enquêteur doit pouvoir montrer}
}
Provenance & De quel système, de quel emplacement, à quelle heure la pièce
a été tirée \\
Intégrité & Que le contenu examiné est celui qui a été recueilli \\
Continuité & Qui a détenu la pièce, à chaque instant, et dans quelles
conditions \\
Méthode & Comment la pièce a été acquise et examinée, avec quels outils et
quelles versions \\
Interprétation & Ce que la pièce signifie, et les autres explications
possibles \\
\end{tableaucours}
\vspace{-0.8em}
\begin{center}{\small\itshape Tableau~2.1. Les fondements de la force probante}\end{center}

% ============================================================================
\section{L'Empreinte Cryptographique}
\label{sec:c02-empreinte}
% ============================================================================

\begin{boxdef}
\textbf{\termecle{Fonction de hachage cryptographique}}~: fonction qui
associe à une donnée de taille quelconque une valeur de taille fixe,
l'\textbf{empreinte} (\textit{hash}), de sorte qu'il soit pratiquement
impossible de trouver deux données différentes de même empreinte
(résistance aux collisions), ou de construire une donnée qui ait une
empreinte donnée (résistance aux préimages).
\end{boxdef}

Si l'on modifie un seul bit d'une image disque de 500~Go, son empreinte
SHA-256 change entièrement. Recalculer l'empreinte à chaque étape et
constater qu'elle n'a pas varié démontre donc que le contenu est resté le
même, sans avoir à comparer les deux copies octet par octet.

\subsection{Quelle fonction choisir}

\begin{tableaucours}{L{2.0cm}cL{4.6cm}L{3.8cm}}{%
    \tableauHeader{Fonction} &
    \tableauHeader{Taille} &
    \tableauHeader{État} &
    \tableauHeader{Usage en enquête}
}
MD5 & 128 bits & Collisions calculables depuis 2004~; on fabrique deux
fichiers distincts de même empreinte & Complément pour comparer à
d'anciennes bases, jamais seule \\
SHA-1 & 160 bits & Collision pratique publiée en 2017 & Complément
seulement \\
SHA-256 & 256 bits & Aucune attaque pratique connue & Empreinte de
référence \\
SHA-512, SHA-3 & 512 bits et plus & Aucune attaque pratique connue &
Possibles, moins répandues dans les outils \\
\end{tableaucours}
\vspace{-0.8em}
\begin{center}{\small\itshape Tableau~2.2. Fonctions de hachage et usage en enquête}\end{center}

La faiblesse de MD5 porte sur les collisions, pas directement sur les
préimages~: un attaquant peut fabriquer deux fichiers de même empreinte, il
ne peut pas, à ce jour, fabriquer un fichier qui ait l'empreinte d'un
fichier donné. Une image acquise avec sa seule empreinte MD5 reste donc
défendable en pratique, mais elle offre à la partie adverse un argument
qu'une empreinte SHA-256 lui refuse. On calcule aujourd'hui SHA-256, et l'on
ajoute MD5 quand un outil ou une base de référence l'exige.

\subsection{Ce que l'empreinte ne prouve pas}

\begin{itemize}
    \item Elle ne prouve pas que la copie est complète~: une acquisition qui
          a sauté des secteurs illisibles a une empreinte stable et un
          contenu lacunaire. Le journal de l'outil d'acquisition le dit.
    \item Elle ne prouve pas que la pièce vient du système annoncé~: c'est
          le rôle de la provenance, consignée au recueil.
    \item Elle ne prouve rien sur la période qui précède son premier
          calcul~: une capture modifiée avant d'être hachée garde, ensuite,
          une empreinte parfaitement stable.
\end{itemize}

La troisième limite explique pourquoi on calcule l'empreinte le plus tôt
possible, idéalement pendant l'acquisition, et pourquoi on la fait
constater par une seconde personne.

% ============================================================================
\section{La Chaîne de Possession}
\label{sec:c02-chaine}
% ============================================================================

\begin{boxdef}
\textbf{\termecle{Chaîne de possession}} (\textit{chain of custody})~:
documentation continue qui retrace, pour chaque pièce, chaque personne qui
l'a détenue, chaque transfert, chaque lieu de conservation et chaque
opération réalisée, depuis le recueil jusqu'à la restitution ou la
destruction.
\end{boxdef}

L'ISO/IEC~27037~\cite{iso27037_2012} demande que cette documentation permette
de répondre, pour chaque pièce, à des questions simples~: qui l'a recueillie,
quand, où, comment et pourquoi~; qui l'a détenue ensuite et à quel moment
elle a changé de mains~; quelles opérations ont été faites sur elle. On
paraphrase ici ces exigences sous la forme d'un formulaire.

\begin{boxprocedure}
\textbf{Le formulaire de chaîne de possession d'une pièce}
\begin{enumerate}
    \item Identifiant unique de la pièce et numéro de scellé.
    \item Description~: nature, marque, modèle, numéro de série, capacité.
    \item Provenance~: système, emplacement physique, propriétaire.
    \item Recueil~: date, heure, lieu, personne, méthode, outil et version.
    \item Empreintes~: algorithme, valeur, heure du calcul, témoin.
    \item Conditionnement~: enveloppe, sac antistatique, scellé, photo.
    \item Pour chaque transfert~: remettant, destinataire, date, heure,
          motif, état du scellé, signatures.
    \item Pour chaque examen~: personne, date, opération, copie utilisée.
    \item Conservation~: lieu, conditions (local fermé, armoire forte).
    \item Sort final~: restitution, destruction, archivage, avec
          autorisation.
\end{enumerate}
\end{boxprocedure}

\subsection{Le scellé}

Un scellé est un conditionnement qui ne s'ouvre pas sans laisser de marque~:
sac numéroté à fermeture inviolable, enveloppe signée sur les rabats et
recouverte d'un adhésif transparent. Le numéro du scellé figure dans le
formulaire. À chaque ouverture, on note qui l'a ouvert, pourquoi, et le
numéro du nouveau scellé. Dans la procédure pénale, la mise sous scellé
relève de l'officier de police judiciaire, comme le rappelle l'article~53 de
la loi n°~2010/012.

\begin{boxjuris}
\textbf{\loicam{n°~2010/012}, article~53}~\cite{loi2010012}~:
\textit{« (1) Les perquisitions en matière de cybercriminalité sont
susceptibles de porter sur des données qui peuvent être des supports
physiques ou des copies réalisées en présence des personnes qui assistent à
la perquisition. (2) Lorsqu'une copie des données saisies a été faite,
celle-ci peut être détruite sur instruction du Procureur de la République
pour des raisons de sécurité. (3) Sur accord du Procureur de la République,
seuls seront gardés sous scellé par l'Officier de Police Judiciaire, les
objets, documents et données utilisées à la manifestation de la vérité.
(4) Les personnes présentes lors de la perquisition peuvent être
réquisitionnées de fournir les renseignements sur les objets, documents et
données saisis. »}

\medskip
\textbf{\loicam{n°~2010/012}, article~54}~\cite{loi2010012}~:
\textit{« Les perquisitions et les saisies sont effectuées conformément aux
dispositions du Code de Procédure Pénale en tenant compte du dépérissement
des preuves. »}

\medskip
\textbf{Portée} (analyse). L'article~53 admet que la saisie porte sur une
copie des données, réalisée devant les personnes présentes~: c'est la base
légale de l'acquisition sur place, sans emporter le serveur. L'article~54
fait du dépérissement des preuves un critère de la conduite des saisies, ce
qui justifie de capturer d'abord ce qui disparaît le plus vite.
\end{boxjuris}

% ============================================================================
\section{La Chaîne de NYANGA}
\label{sec:c02-nyanga}
% ============================================================================

\begin{boxscenario}{Opération NYANGA. Le parcours du scellé NC-01}
\begin{tableaucours}{L{2.6cm}L{2.6cm}L{3.0cm}L{3.8cm}}{%
    \tableauHeader{Date et heure} &
    \tableauHeader{Détenteur} &
    \tableauHeader{Lieu} &
    \tableauHeader{Opération}
}
14/03, 10~h~02 & Mme NDOUMBE & Salle serveurs & Capture mémoire de
SRV-AD01 \\
14/03, 10~h~41 & Mme NDOUMBE & Salle serveurs & Empreinte SHA-256,
contresignée par M.~KAMGA \\
14/03, 15~h~50 & M. KAMGA & Bureau du DSI & Enveloppe kraft, ruban
adhésif, sans numéro \\
15/03, 9~h~30 & Mme MBARGA & Bureau du DSI & Empreinte recalculée,
identique~; mise sous scellé NC-01 \\
15/03, 11~h~00 & Mme MBARGA & Salle de crise & Copie de travail, empreinte
vérifiée \\
16/03, 10~h~15 & Adjudant-chef NANA & Brigade de recherches & Remise sur
réquisition, scellé intact, procès-verbal \\
17/03, 14~h~00 & M. ABOUBAKAR & Laboratoire de l'expert & Ouverture du
scellé, copie, nouveau scellé NC-01-E1 \\
\end{tableaucours}

\medskip
\textbf{Constatations.} La pièce a été détenue par cinq personnes
successives. Entre le 14~mars à 15~h~50 et le 15~mars à 9~h~30, soit
17~h~40, elle est restée dans un bureau ouvert, sans scellé numéroté et sans
consignation. Les empreintes calculées le 14 à 10~h~41 et le 15 à 9~h~30
sont identiques.

\textbf{Analyse} (confiance élevée sur l'intégrité, moyenne sur la
continuité). Le contenu n'a pas changé~: l'égalité des empreintes l'établit.
La continuité, elle, présente un trou. Il n'affaiblit pas l'intégrité, mais
il affaiblit la réponse à une autre objection~: la substitution d'un disque
par un autre de même contenu est sans effet, mais la manipulation du disque
pendant la nuit, par exemple pour en copier le contenu, ne peut être exclue.
Le rapport doit mentionner ce trou, et non le taire.
\end{boxscenario}

\begin{boxlecture}{Les six scellés de l'affaire au 17 mars}
\begin{tableaucours}{L{1.4cm}L{4.6cm}ccc}{%
    \tableauHeader{Scellé} &
    \tableauHeader{Pièce} &
    \tableauHeader{Intégrité} &
    \tableauHeader{Continuité} &
    \tableauHeader{Provenance}
}
NC-01 & Capture mémoire de SRV-AD01 & \OUI{} & \NEK{} & \OUI{} \\
NC-02 & Journaux de la passerelle (export) & \OUI{} & \OUI{} & \OUI{} \\
NC-03 & Image disque de SRV-CPT01 & \OUI{} & \OUI{} & \OUI{} \\
NC-04 & Capture réseau de la sonde & \OUI{} & \OUI{} & \OUI{} \\
NC-05 & Journaux de l'interface de paiement & \OUI{} & \OUI{} & \NEK{} \\
NC-06 & Capture d'écran WhatsApp d'un technicien & \NEK{} & \NON{} & \NEK{} \\
\end{tableaucours}

\medskip
\textbf{Lecture.} Sur dix-huit cotes, treize \OUI{}, une \NON{} et quatre
\NEK{}. NC-05 a été envoyé par le prestataire de l'interface sous forme de
fichier joint à un courriel~: son intégrité depuis l'envoi est établie, mais
rien ne dit comment le prestataire l'a extrait. NC-06 ne sert qu'à orienter
l'enquête~; on ne fonde aucune conclusion sur elle. NC-01 reste utilisable,
avec la réserve sur la nuit du 14 au 15.
\end{boxlecture}

% ============================================================================
\section{Les Ruptures et leurs Remèdes}
\label{sec:c02-ruptures}
% ============================================================================

\begin{tableaucours}{L{3.6cm}L{4.2cm}L{4.4cm}}{%
    \tableauHeader{Rupture} &
    \tableauHeader{Conséquence} &
    \tableauHeader{Remède}
}
Empreinte calculée tard & Rien ne prouve l'état d'avant & L'écrire, et
dater l'acquisition par d'autres sources (journal de l'outil) \\
Trou dans la détention & Une manipulation ne peut être exclue & Le dire
dans le rapport, recueillir les témoignages, vérifier les journaux d'accès
au local \\
Original examiné & L'original a pu changer & Recalculer l'empreinte,
documenter l'opération, travailler ensuite sur copie \\
Pièce transmise par courriel & Provenance invérifiable avant l'envoi &
Demander l'extraction en présence d'un témoin ou par procès-verbal \\
Scellé ouvert sans trace & La continuité est rompue & Nouveau scellé,
procès-verbal d'ouverture, signatures \\
\end{tableaucours}
\vspace{-0.8em}
\begin{center}{\small\itshape Tableau~2.3. Ruptures de chaîne et remèdes}\end{center}

Une rupture ne rend pas la pièce inutilisable. Elle en réduit la force
probante, d'autant plus que la pièce est seule à soutenir une conclusion.
Quand plusieurs pièces indépendantes convergent, la faiblesse de l'une est
compensée par les autres~: à NYANGA, l'outil d'accès à distance vu dans la
capture mémoire (NC-01) est aussi visible dans la capture réseau (NC-04),
par ses connexions vers 203.0.113.47. La conclusion ne repose pas sur NC-01
seule.

% ============================================================================
\section{Ce que la Preuve demande à chaque Métier}
\label{sec:c02-metiers}
% ============================================================================

\begin{boxmetier}{Application OPJ/Expert Judiciaire}
\textbf{Recevoir un scellé}
\begin{enumerate}
    \item Vérifier le numéro et l'état du scellé avant toute ouverture.
    \item Ouvrir en présence d'un témoin, photographier, consigner.
    \item Recalculer l'empreinte et la comparer à celle du formulaire.
    \item Travailler sur une copie, et resceller l'original.
    \item Signaler dans le rapport toute rupture constatée dans la chaîne
          reçue.
\end{enumerate}
\end{boxmetier}

\begin{boxmetier}{Analyste en réponse à incident}
Le formulaire de chaîne de possession s'ouvre dès le recueil, même quand
personne ne parle encore de plainte. Le coût est de quelques minutes par
pièce~; une chaîne reconstituée après coup ne convainc personne.
\end{boxmetier}

\begin{boxmetier}{RSSI}
Prévoir, avant l'incident, le matériel de conditionnement (sacs numérotés,
étiquettes, disques neufs) et un lieu de conservation fermé. À NYANGA, il
n'y avait ni l'un ni l'autre, d'où l'enveloppe kraft sur le bureau du DSI.
\end{boxmetier}

\separateur

\begin{boxresume}
\textbf{Ce qu'il faut retenir de ce chapitre~:}
\begin{itemize}
    \item La recevabilité dit si une pièce entre au débat~; la force
          probante dit ce qu'elle pèse. L'enquêteur construit la seconde.
    \item L'empreinte SHA-256 prouve que le contenu n'a pas changé entre
          deux calculs~; elle ne prouve ni la complétude, ni la provenance,
          ni l'état d'avant le premier calcul.
    \item La chaîne de possession retrace chaque détenteur, chaque
          transfert, chaque opération~; un trou se dit, il ne se cache pas.
    \item Les articles~53 et~54 de la loi n°~2010/012 admettent la saisie
          par copie et imposent de tenir compte du dépérissement des
          preuves.
    \item À NYANGA, NC-01 a une intégrité établie et une continuité \NEK{}
          sur 17~h~40~; la conclusion qu'elle soutient est confirmée par
          NC-04.
\end{itemize}
\end{boxresume}

\begin{boxinfo}
\textbf{Points de contrôle.} Vous avez compris ce chapitre si vous savez~:
\begin{itemize}
    \item expliquer pourquoi MD5 seul ne suffit plus, et ce qui reste vrai
          à son sujet~;
    \item remplir un formulaire de chaîne de possession~;
    \item coter l'intégrité, la continuité et la provenance d'une pièce.
\end{itemize}
\end{boxinfo}

% ============================================================================
\section*{Exercices}
% ============================================================================

\begin{boxexo}[title={\ding{45}\quad Exercice~2.1. Ce que dit l'empreinte \niveaubase}]
Vrai ou faux~? Justifiez.
\begin{enumerate}[(a)]
    \item Deux empreintes SHA-256 identiques prouvent que deux fichiers sont
          identiques.
    \item Une empreinte identique prouve que l'image disque est complète.
    \item Une image hachée en MD5 seulement est irrecevable.
    \item Il suffit de calculer l'empreinte au moment du rapport.
\end{enumerate}
\end{boxexo}

\begin{boxexo}[title={\ding{45}\quad Exercice~2.2. Remplir le formulaire \niveaumoyen}]
Le 15~mars à 16~h~00, SANKOFA EXPERTISE acquiert l'image du disque de
SRV-CPT01 (scellé NC-03). Rédigez le formulaire de chaîne de possession de
cette pièce jusqu'à sa remise à l'OPJ le 16~mars, en inventant de façon
cohérente les éléments manquants (outil, version, empreintes, témoins).
\end{boxexo}

\begin{boxexo}[title={\ding{45}\quad Exercice~2.3. La nuit du 14 au 15 \niveaumoyen}]
\begin{enumerate}[(a)]
    \item Rédigez le paragraphe du rapport qui mentionne le trou de la
          chaîne de NC-01, en séparant constatation et analyse.
    \item Quelles vérifications complémentaires proposez-vous pour réduire
          l'incertitude sur cette nuit~?
    \item Le trou change-t-il une conclusion de l'enquête~? Justifiez à
          l'aide de NC-04.
\end{enumerate}
\end{boxexo}

\begin{boxexo}[title={\ding{45}\quad Exercice~2.4. NC-05 et le prestataire \niveauexpert}]
Le prestataire de l'interface de paiement a envoyé ses journaux en pièce
jointe d'un courriel. L'avocat de la défense soutiendra qu'ils ont pu être
modifiés chez le prestataire.
\begin{enumerate}[(a)]
    \item Quelles cotes donnez-vous à cette pièce, et pourquoi~?
    \item Comment l'OPJ aurait-il dû obtenir ces journaux~? Appuyez-vous
          sur les articles~53 et~54.
    \item Quelles pièces indépendantes permettent de corroborer les heures
          des virements~?
\end{enumerate}
\end{boxexo}

\bigskip

\begin{boxinfo}
\textbf{Pour aller plus loin.}
\begin{itemize}
    \item L'ISO/IEC~27037~\cite{iso27037_2012}, sur l'identification, le
          recueil et la préservation.
    \item Le guide de l'ACPO~\cite{acpo2012}, pour ses quatre principes de
          conduite.
\end{itemize}
\end{boxinfo}

\begin{boxinfo}
\textbf{Transition.} La chaîne de possession garde la pièce~; encore
faut-il l'avoir acquise sans l'altérer. Copier un disque, capturer une
mémoire, figer un journal~: c'est l'objet du Chapitre~\ref{chap:acquisition}.
\end{boxinfo}


    % PARTIE 2 · MÉTHODE (NIST SP 800-86 ; ISO/IEC 27042)
    \part{Méthode : Acquérir et Analyser}
    \label{part:methode}

        % ============================================================================
% Fichier  : partie2/P02_C03_acquerir_preserver.tex
% Titre    : Acquérir et Préserver
% Auteur   : MINKA MI NGUIDJOÏ Thierry Emmanuel
% Version  : 1.0 -- novembre 2026
% Corrigés : corriges/Corr_C03.tex
% ============================================================================
% Sources : RFC 3227 (ordre de volatilité) ; ISO/IEC 27037 (identification,
% recueil, acquisition, préservation, paraphrasée) ; NIST SP 800-86.
% Outils cités (existants) : bloqueurs d'écriture, formats brut, E01 (EWF),
% AFF4 ; outils de capture mémoire WinPmem, DumpIt, LiME, AVML.
% Faits nouveaux (bible) : acquisition de SRV-CPT01 le 15/03, 16 h 00 à
% 19 h 25, E01, NC-03 ; SRV-FIC01 chiffré BitLocker, acquisition logique à
% chaud le 18/03 (NC-07) ; journaux de la passerelle : rotation à 60 jours.
% ============================================================================

\chapter{Acquérir et Préserver}
\label{chap:acquisition}
\chaptermark{Acquérir et préserver}

\epigraph{%
    « On n'examine jamais l'original.\\
    On examine une copie dont on peut prouver qu'elle lui est égale. »%
}{Principe du cours}

\niveauChapitre{Fondamental}{%
    Maîtrise du Chapitre~\ref{chap:preuve}
}

\begin{boxobjectifs}
\begin{itemize}
    \item Appliquer l'ordre de volatilité à une situation concrète.
    \item Choisir entre acquisition à froid et acquisition à chaud, et
          justifier ce choix.
    \item Acquérir un disque avec un bloqueur d'écriture, dans un format
          adapté, et vérifier l'acquisition.
    \item Capturer la mémoire vive en connaissant l'empreinte que laisse
          l'outil.
    \item Préserver des journaux avant leur rotation et des données
          détenues par des tiers.
    \item Documenter une acquisition de façon qu'un autre puisse la
          contrôler.
\end{itemize}
\end{boxobjectifs}

\bigskip

% ============================================================================
\section*{Avant de commencer : un serveur chiffré}
% ============================================================================

Bafoussam, mardi 18 mars 2025. Mme~MBARGA veut l'image du serveur de
fichiers SRV-FIC01, celui qui a envoyé 3,2~Go vers l'extérieur. Le serveur
est toujours allumé, débranché du réseau depuis le vendredi. M.~KAMGA
rappelle un détail~: son disque de données est chiffré par BitLocker, et la
clé de récupération est conservée chez le prestataire qui a installé la
machine, à Douala. Si l'on éteint le serveur pour en extraire le disque, on
obtiendra une copie parfaite d'un volume chiffré, illisible tant que la clé
n'arrive pas. Si on le laisse allumé, le volume est déverrouillé et lisible,
mais toute copie faite à chaud modifie un peu la machine.

\begin{boxalert}
\textbf{Violation typique.} Appliquer mécaniquement la règle « on éteint
puis on copie le disque » à une machine dont le chiffrement rendra le disque
illisible une fois éteinte.

\smallskip
\textbf{Règle~:} on choisit la méthode d'acquisition d'après ce que l'on
perdrait avec chacune, et l'on écrit la raison du choix. La meilleure
méthode est celle qui conserve le plus de preuve utile au prix de la plus
petite altération documentée.
\end{boxalert}

\separateur

% ============================================================================
\section{L'Ordre de Volatilité}
\label{sec:c03-volatilite}
% ============================================================================

La RFC~3227~\cite{rfc3227} recommande de recueillir les traces de la plus
volatile à la moins volatile. On en tire l'ordre suivant pour un serveur.

\begin{tableaucours}{cL{4.6cm}L{2.6cm}L{4.2cm}}{%
    \tableauHeader{Rang} &
    \tableauHeader{Trace} &
    \tableauHeader{Durée de vie} &
    \tableauHeader{Moyen d'acquisition}
}
1 & Registres, cache du processeur & Nanosecondes & Hors de portée en
pratique \\
2 & Mémoire vive~: processus, connexions, tables de routage et ARP, clés
en mémoire & Jusqu'à l'extinction & Capture de la mémoire \\
3 & Fichiers temporaires, fichier d'échange & Heures à jours & Image du
disque \\
4 & Disques & Mois à années & Image du disque \\
5 & Journaux distants, données de supervision & Selon la rotation &
Export, réquisition \\
6 & Configuration physique, topologie & Stable & Photographies, schéma \\
7 & Supports d'archive & Années & Copie \\
\end{tableaucours}
\vspace{-0.8em}
\begin{center}{\small\itshape Tableau~3.1. L'ordre de volatilité (d'après la RFC~3227)}\end{center}

Le rang~5 surprend souvent~: un journal distant paraît durable, alors que
sa rotation peut l'effacer en quelques jours. Les journaux de la passerelle
de NYANGA CRÉDIT ne gardent que 60~jours~: le premier accès du 2~mars
disparaîtrait le 1\textsuperscript{er}~mai. Leur export, le 14~mars, était
une urgence au même titre que la capture mémoire.

% ============================================================================
\section{À Froid ou à Chaud}
\label{sec:c03-froid-chaud}
% ============================================================================

\begin{boxdef}
\textbf{\termecle{Acquisition à froid}} (\textit{dead acquisition})~: copie
d'un support après extinction de la machine, à l'aide d'un bloqueur
d'écriture. Elle ne modifie pas le support et perd tout ce qui était en
mémoire.

\smallskip
\textbf{\termecle{Acquisition à chaud}} (\textit{live acquisition})~:
recueil sur une machine allumée (mémoire, volume déverrouillé, état du
système). Elle conserve le volatil et laisse une empreinte de l'outil sur la
machine, qu'il faut documenter.
\end{boxdef}

\begin{tableaucours}{L{4.4cm}L{3.8cm}L{4.0cm}}{%
    \tableauHeader{Situation} &
    \tableauHeader{Méthode} &
    \tableauHeader{Raison}
}
Serveur allumé, attaque possiblement en cours & À chaud d'abord (mémoire),
puis à froid & Processus et connexions perdus à l'extinction \\
Volume chiffré, clé indisponible & À chaud (copie logique du volume
déverrouillé) & Volume illisible une fois éteint \\
Poste éteint à l'arrivée & À froid & Rien de volatil à préserver \\
Machine virtuelle & Instantané de l'hyperviseur, puis copie des fichiers
de disque & Mémoire et disque figés ensemble, sans outil dans la machine \\
Service en nuage & Export par l'interface du fournisseur, réquisition &
Aucun accès au matériel \\
Serveur critique qui ne peut s'arrêter & À chaud, copie logique ciblée &
Arrêt impossible~; périmètre écrit et justifié \\
\end{tableaucours}
\vspace{-0.8em}
\begin{center}{\small\itshape Tableau~3.2. Choisir la méthode d'acquisition}\end{center}

% ============================================================================
\section{Acquérir un Disque}
\label{sec:c03-disque}
% ============================================================================

\subsection{Le bloqueur d'écriture}

Un système d'exploitation écrit sur tout support qu'on lui présente~: il met
à jour des dates d'accès, monte des volumes, écrit des journaux. Le
\termecle{bloqueur d'écriture} s'interpose entre le support et la station
d'acquisition et laisse passer les lectures en refusant les écritures. Le
bloqueur matériel est un boîtier placé sur le câble~; le bloqueur logiciel
est un réglage du système de la station. Le premier est préféré en enquête
judiciaire, parce que son comportement se teste et se documente
indépendamment du système. On vérifie un bloqueur avant usage, sur un
support de test, et l'on note la vérification.

\subsection{Les formats d'image}

\begin{tableaucours}{L{2.2cm}L{4.8cm}L{5.2cm}}{%
    \tableauHeader{Format} &
    \tableauHeader{Caractéristiques} &
    \tableauHeader{Usage}
}
Brut (\texttt{dd}, \texttt{.raw}) & Copie secteur par secteur, sans
métadonnées ni compression & Universel, lisible par tout outil~; empreinte
et notes dans des fichiers séparés \\
E01 (EWF) & Métadonnées intégrées (enquêteur, numéro de pièce, notes),
compression, contrôle par blocs, empreinte stockée & Format courant des
outils d'enquête~; bon compromis \\
AFF4 & Format ouvert, métadonnées riches, adapté aux très gros volumes et à
la mémoire & Moins répandu, utile pour les acquisitions volumineuses \\
\end{tableaucours}
\vspace{-0.8em}
\begin{center}{\small\itshape Tableau~3.3. Formats d'image}\end{center}

\subsection{Les pièges du stockage moderne}

Les disques à mémoire flash (SSD) gèrent eux-mêmes l'effacement des blocs
libérés~: la commande TRIM permet au contrôleur d'effacer réellement des
données supprimées, parfois quelques minutes après la suppression. Deux
acquisitions successives d'un SSD allumé peuvent donc différer, sans que
personne ait écrit dessus, et la récupération de fichiers supprimés y est
moins fructueuse que sur un disque magnétique. Les secteurs défectueux d'un
disque magnétique usé, eux, produisent des lectures incomplètes~: l'outil
d'acquisition les journalise et l'enquêteur reporte leur nombre dans son
rapport.

\begin{boxprocedure}
\textbf{Acquérir un disque à froid}
\begin{enumerate}
    \item Photographier la machine, ses câbles et ses étiquettes~; noter le
          numéro de série du disque.
    \item Vérifier le bloqueur d'écriture sur un support de test, et le
          consigner.
    \item Relier le disque au bloqueur, puis à la station d'acquisition.
    \item Choisir le format (E01 par défaut), saisir les métadonnées~:
          numéro de pièce, enquêteur, description.
    \item Acquérir, en calculant SHA-256 (et MD5 si une base l'exige).
    \item Lire le journal de l'outil~: secteurs illisibles, erreurs,
          durée.
    \item Vérifier l'image~: recalculer l'empreinte de l'image et la
          comparer à celle du support.
    \item Faire une copie de travail, vérifier son empreinte, et sceller
          l'image de référence.
\end{enumerate}
\end{boxprocedure}

% ============================================================================
\section{Capturer la Mémoire}
\label{sec:c03-memoire}
% ============================================================================

La mémoire vive contient ce qu'aucun disque ne garde~: la liste des
processus actifs, les connexions réseau ouvertes, les commandes saisies, les
clés de chiffrement des volumes ouverts, et parfois des programmes qui
n'existent qu'en mémoire. On la capture avec un outil exécuté sur la machine
elle-même~: WinPmem ou DumpIt sous Windows, LiME ou AVML sous Linux. Ces
outils sont libres ou gratuits et largement documentés.

\begin{boxinfo}
\textbf{Une acquisition qui ne se répète pas.} La mémoire change pendant la
capture~: les processus continuent de tourner, et l'outil lui-même occupe de
la mémoire. On ne pourra jamais recalculer l'empreinte de la source pour la
comparer à celle de la capture. L'empreinte de la mémoire capturée prouve
seulement que la capture n'a pas changé depuis. C'est pourquoi l'on
documente avec soin l'heure, l'outil, sa version, le support utilisé et ce
qui tournait sur la machine.
\end{boxinfo}

\begin{tableaucours}{L{4.0cm}L{8.2cm}}{%
    \tableauHeader{Précaution} &
    \tableauHeader{Raison}
}
Lancer l'outil depuis un support externe & Ne pas écrire l'outil sur le
disque à examiner \\
Écrire la capture sur un support externe & Ne pas écraser l'espace libre
du disque, où dorment des fichiers supprimés \\
Ne rien installer & Une installation modifie le registre et le disque \\
Noter l'heure système et l'heure réelle & Mesurer la dérive de l'horloge \\
Consigner chaque commande & L'empreinte de l'enquêteur doit pouvoir être
distinguée de celle de l'attaquant \\
\end{tableaucours}
\vspace{-0.8em}
\begin{center}{\small\itshape Tableau~3.4. Précautions de la capture mémoire}\end{center}

% ============================================================================
\section{Les Journaux et les Données des Tiers}
\label{sec:c03-journaux}
% ============================================================================

Les journaux sont souvent la seule trace d'un événement ancien. On les
préserve par un export complet, avec son empreinte, plutôt que par une
recherche filtrée~: un filtre écrit le 14~mars sur les mots « svc\_backup »
aurait laissé de côté les autres comptes, alors qu'on ne savait pas encore
lesquels l'attaquant avait utilisés.

Les données détenues par des tiers (opérateur de paiement mobile, fournisseur
d'accès, hébergeur, éditeur d'un service en nuage) ne se copient pas~: on les
demande, et en matière pénale on les requiert. Le Chapitre~\ref{chap:normes-droit}
précise les fondements de ces demandes. Deux réflexes valent dans tous les
cas~: demander très tôt la conservation des données, avant même leur
communication, et demander comment elles ont été extraites.

% ============================================================================
\section{Les Acquisitions de NYANGA}
\label{sec:c03-nyanga}
% ============================================================================

\begin{boxscenario}{Opération NYANGA. Trois serveurs, trois méthodes}
\textbf{Constatations.}
\begin{itemize}
    \item SRV-AD01~: capture mémoire à chaud le 14~mars (NC-01)~; le serveur,
          indispensable à toutes les connexions, est ensuite remplacé par un
          contrôleur de secours, puis acquis à froid le 16~mars.
    \item SRV-CPT01~: non chiffré~; arrêt contrôlé le 15~mars à 15~h~55,
          acquisition à froid par bloqueur matériel de 16~h~00 à 19~h~25, au
          format E01 (NC-03)~; aucun secteur illisible.
    \item SRV-FIC01~: volume de données chiffré par BitLocker, clé chez le
          prestataire~; capture mémoire puis copie logique du volume
          déverrouillé, à chaud, le 18~mars (NC-07)~; la clé de récupération
          arrive le 22~mars et permet une acquisition à froid complémentaire.
\end{itemize}

\textbf{Analyse} (confiance élevée). Les trois choix sont justifiés par ce
que chaque méthode aurait fait perdre. L'acquisition à chaud de SRV-FIC01 a
modifié la machine (outil exécuté, fichiers temporaires)~; ces
modifications sont consignées heure par heure, ce qui permet de les écarter
de l'analyse. L'acquisition complémentaire du 22~mars confirme que la copie
logique du 18 contenait bien tous les fichiers du volume.
\end{boxscenario}

\begin{boxobjet}{Image du disque de SRV-CPT01, 15 mars 2025}
\textbf{Source~:} serveur SRV-CPT01, disque système de 500~Go.

\textbf{Pièce~:} image E01, scellé NC-03, copie de travail séparée.

\textbf{Méthode~:} arrêt contrôlé, bloqueur matériel vérifié le matin même,
acquisition de 16~h~00 à 19~h~25~; empreinte SHA-256 de la source et de
l'image identiques.

\textbf{Ce qu'elle peut prouver~:} les fichiers présents et supprimés, leurs
dates, les journaux Windows, les traces d'exécution de programmes.

\textbf{Ce qu'elle ne prouve pas~:} ce qui était en mémoire avant l'arrêt du
15~mars, ni ce qui s'est passé sur d'autres serveurs.
\end{boxobjet}

\begin{boxlecture}{Les acquisitions de NYANGA au 22 mars}
\begin{tableaucours}{L{3.4cm}ccc}{%
    \tableauHeader{Question} &
    \tableauHeader{SRV-AD01} & \tableauHeader{SRV-CPT01} &
    \tableauHeader{SRV-FIC01}
}
Mémoire capturée avant arrêt & \OUI{} & \NON{} & \OUI{} \\
Disque acquis sans altération & \OUI{} & \OUI{} & \OUI{} \\
Altérations éventuelles documentées & \OUI{} & \OUI{} & \OUI{} \\
Complétude de l'image vérifiée & \OUI{} & \OUI{} & \OUI{} \\
\end{tableaucours}

\medskip
\textbf{Lecture.} Une seule cote \NON{}~: la mémoire de SRV-CPT01 n'a pas
été capturée avant son arrêt. L'enquête ne pourra donc pas savoir si un
programme de l'attaquant y tournait encore le 15~mars. C'est une limite, et
le rapport la mentionne.
\end{boxlecture}

% ============================================================================
\section{Ce que l'Acquisition demande à chaque Métier}
\label{sec:c03-metiers}
% ============================================================================

\begin{boxmetier}{Analyste en réponse à incident}
Avoir, avant l'incident, une trousse prête~: bloqueur matériel vérifié,
disques neufs, outils de capture mémoire sur un support en lecture seule,
formulaires imprimés. On n'improvise pas une acquisition à trois heures du
matin.
\end{boxmetier}

\begin{boxmetier}{Application OPJ/Expert Judiciaire}
Quand l'acquisition a été faite par d'autres, l'expert ne la refait pas
nécessairement~; il la contrôle~: journal de l'outil, empreintes, bloqueur,
méthode. S'il refait une acquisition, il explique pourquoi.
\end{boxmetier}

\begin{boxmetier}{RSSI}
Conserver, hors des machines, les clés de récupération des volumes chiffrés,
dans un coffre accessible à l'équipe de crise. À NYANGA, quatre jours ont
été perdus à attendre une clé détenue par un prestataire.
\end{boxmetier}

\separateur

\begin{boxresume}
\textbf{Ce qu'il faut retenir de ce chapitre~:}
\begin{itemize}
    \item On recueille de la trace la plus volatile à la moins volatile~;
          les journaux soumis à rotation sont plus fragiles qu'ils ne
          paraissent.
    \item Le choix entre froid et chaud se décide d'après ce que l'on
          perdrait, et se justifie par écrit.
    \item Un disque s'acquiert derrière un bloqueur d'écriture vérifié,
          dans un format qui garde les métadonnées, et l'acquisition se
          vérifie.
    \item Une capture mémoire ne se répète pas~: on documente l'outil,
          l'heure et l'empreinte laissée sur la machine.
    \item À NYANGA, trois serveurs ont reçu trois méthodes différentes,
          chacune justifiée~; la mémoire de SRV-CPT01 manque.
\end{itemize}
\end{boxresume}

\begin{boxinfo}
\textbf{Points de contrôle.} Vous avez compris ce chapitre si vous savez~:
\begin{itemize}
    \item classer des traces selon leur volatilité~;
    \item choisir et justifier une méthode d'acquisition~;
    \item lire le journal d'un outil d'acquisition et y repérer ce qui
          compte.
\end{itemize}
\end{boxinfo}

% ============================================================================
\section*{Exercices}
% ============================================================================

\begin{boxexo}[title={\ding{45}\quad Exercice~3.1. Classer par volatilité \niveaubase}]
Classez de la plus volatile à la moins volatile~: (a) le journal de la
passerelle, conservé 60~jours~; (b) la table ARP du contrôleur de domaine~;
(c) le fichier d'échange de SRV-CPT01~; (d) les sauvegardes sur bande de
février~; (e) les connexions réseau ouvertes de SRV-FIC01~; (f) la capture
de la sonde, conservée 48~heures.
\end{boxexo}

\begin{boxexo}[title={\ding{45}\quad Exercice~3.2. Choisir la méthode \niveaumoyen}]
Pour chaque situation, proposez une méthode d'acquisition et justifiez-la
par ce que l'on perdrait avec l'autre~:
\begin{enumerate}[(a)]
    \item un portable d'agent, éteint, saisi dans un sac~;
    \item le serveur de l'interface de paiement, qui ne peut pas s'arrêter
          plus de dix minutes~;
    \item une machine virtuelle hébergée sur l'hyperviseur du siège~;
    \item la messagerie de l'établissement, hébergée par un fournisseur en
          nuage.
\end{enumerate}
\end{boxexo}

\begin{boxexo}[title={\ding{45}\quad Exercice~3.3. Lire un journal d'acquisition \niveaumoyen}]
Le journal d'une acquisition indique~: « Source~: 500\,107\,862\,016
octets. Secteurs lus avec erreur~: 384. Remplacés par des zéros. Empreinte
SHA-256 de l'image~: 9f3c\ldots{} Vérification~: empreinte de l'image
identique à l'empreinte calculée pendant l'acquisition. »
\begin{enumerate}[(a)]
    \item L'image est-elle intègre~? Est-elle complète~?
    \item Quelle taille de données est perdue, si les secteurs font 512
          octets~?
    \item Comment le rapport doit-il présenter ce fait~?
\end{enumerate}
\end{boxexo}

\begin{boxexo}[title={\ding{45}\quad Exercice~3.4. SRV-FIC01 à chaud \niveauexpert}]
\begin{enumerate}[(a)]
    \item Rédigez la justification écrite du choix d'une acquisition à chaud
          de SRV-FIC01 le 18~mars.
    \item Listez les traces que l'acquisition à chaud a pu laisser sur la
          machine, et dites comment les distinguer de celles de
          l'attaquant.
    \item La clé BitLocker arrive le 22~mars. Que permet l'acquisition à
          froid complémentaire, et quelle comparaison faites-vous entre les
          deux copies~?
\end{enumerate}
\end{boxexo}

\bigskip

\begin{boxinfo}
\textbf{Pour aller plus loin.}
\begin{itemize}
    \item La RFC~3227~\cite{rfc3227} et l'ISO/IEC~27037~\cite{iso27037_2012}.
    \item B.~Carrier, \textit{File System Forensic
          Analysis}~\cite{carrier2005}, chapitre sur l'acquisition.
\end{itemize}
\end{boxinfo}

\begin{boxinfo}
\textbf{Transition.} L'image du disque de SRV-CPT01 est sous scellé, sa
copie de travail est prête. Que contient-elle~? Le Chapitre~\ref{chap:disque}
apprend à lire un système de fichiers et ses artefacts.
\end{boxinfo}

        % ============================================================================
% Fichier  : partie2/P02_C04_analyse_disque.tex
% Titre    : L'Analyse du Disque
% Auteur   : MINKA MI NGUIDJOÏ Thierry Emmanuel
% Version  : 1.0 -- novembre 2026
% Corrigés : corriges/Corr_C04.tex
% ============================================================================
% Sources : Carrier (2005) pour les systèmes de fichiers ; documentation
% Microsoft des journaux de sécurité (identifiants 4624, 4625, 4688, 4720,
% 7045, 1102, présentés sans citation) ; The Sleuth Kit et Autopsy.
% Faits nouveaux (bible) : sur SRV-CPT01, session RDP de adm.kamga depuis
% SRV-AD01 le 12/03 à 22 h 39 ; raccourci api.ini.lnk créé à 22 h 44 ;
% archive cpt.7z créée puis supprimée à 22 h 52, récupérée ; prefetch
% désactivé (Windows Server) ; journal Amcache.
% ============================================================================

\chapter{L'Analyse du Disque}
\label{chap:disque}
\chaptermark{L'analyse du disque}

\epigraph{%
    « Supprimer un fichier, pour un système de fichiers,\\
    c'est d'abord oublier où il se trouvait. »%
}{Principe du cours}

\niveauChapitre{Intermédiaire}{%
    Maîtrise du Chapitre~\ref{chap:acquisition}
}

\begin{boxobjectifs}
\begin{itemize}
    \item Lire la structure d'un disque~: table de partitions, volumes,
          système de fichiers.
    \item Expliquer comment NTFS enregistre un fichier, ses dates et sa
          suppression.
    \item Récupérer un fichier supprimé par ses métadonnées ou par ses
          signatures.
    \item Exploiter les artefacts Windows qui prouvent une connexion, une
          exécution ou l'ouverture d'un fichier.
    \item Interpréter un horodatage avec son fuseau et ses limites.
\end{itemize}
\end{boxobjectifs}

\bigskip

% ============================================================================
\section*{Avant de commencer : la question posée au disque}
% ============================================================================

Yaoundé, laboratoire de SANKOFA EXPERTISE, 19 mars 2025. La copie de travail
de l'image de SRV-CPT01 est montée en lecture seule. Le disque compte un peu
plus de 300\,000 fichiers. Mme~MBARGA n'ouvre pas l'outil d'analyse au
hasard~: elle écrit d'abord trois questions. L'attaquant s'est-il connecté à
ce serveur, et quand~? A-t-il ouvert le fichier qui contient la clé de
l'interface de paiement~? A-t-il copié quelque chose et effacé ses traces~?

\begin{boxalert}
\textbf{Violation typique.} Parcourir l'arborescence d'une image en
espérant « tomber sur quelque chose », puis présenter ce que l'on a trouvé
comme le résultat d'un examen complet.

\smallskip
\textbf{Règle~:} on examine un disque pour répondre à des questions écrites
d'avance, en choisissant les artefacts capables d'y répondre, et l'on dit
pour chaque question quelles sources ont été lues.
\end{boxalert}

\separateur

% ============================================================================
\section{Du Disque au Fichier}
\label{sec:c04-structure}
% ============================================================================

Un disque se lit par couches. La \termecle{table de partitions} (MBR sur les
anciens systèmes, GPT sur les récents) découpe le disque en volumes. Chaque
volume porte un \termecle{système de fichiers}, qui organise les données en
fichiers et en répertoires~: NTFS sous Windows, ext4 sous la plupart des
Linux, APFS sous macOS. L'espace d'un volume se répartit entre les données
allouées, les métadonnées du système de fichiers, et l'espace non alloué,
qui contient souvent les restes de fichiers supprimés.

\begin{center}
\begin{tikzpicture}[
    couche/.style = {draw=CouleurBordDef, fill=CouleurDef, minimum width=10cm,
        minimum height=0.75cm, align=center, font=\small\color{CouleurPrimaire}}
]
\node[couche] at (0,0)    {Disque physique~: secteurs de 512 ou 4\,096 octets};
\node[couche] at (0,0.9)  {Table de partitions (MBR, GPT)~: où commence chaque volume};
\node[couche] at (0,1.8)  {Système de fichiers (NTFS, ext4)~: métadonnées et données};
\node[couche] at (0,2.7)  {Fichiers et répertoires~: ce que voit l'utilisateur};
\node[couche] at (0,3.6)  {Artefacts~: journaux, registre, raccourcis, traces d'exécution};
\end{tikzpicture}
\end{center}
\begin{center}{\small\itshape Figure~4.1. Les couches de lecture d'un disque}\end{center}

% ============================================================================
\section{NTFS : la Table MFT et les Dates}
\label{sec:c04-ntfs}
% ============================================================================

NTFS décrit chaque fichier dans un enregistrement de la \termecle{table
MFT} (\textit{Master File Table}). Un enregistrement contient des
attributs~: le nom, les dates, la taille et l'emplacement des données sur le
disque. Les petits fichiers, de quelques centaines d'octets, tiennent même
entièrement dans leur enregistrement.

\subsection{Deux jeux de quatre dates}

Chaque fichier porte quatre dates~: modification du contenu (M), dernier
accès (A), modification de l'enregistrement MFT (C), et création (B). On les
note MACB. NTFS les enregistre deux fois~: dans l'attribut
\texttt{\$STANDARD\_INFORMATION}, que les programmes modifient librement,
et dans l'attribut \texttt{\$FILE\_NAME}, que seul le système met à jour lors
de certaines opérations (création, renommage, déplacement). Un écart
incohérent entre les deux jeux est l'un des signes d'une falsification des
dates, étudiée au Chapitre~\ref{chap:antiforensique}.

\begin{boxinfo}
\textbf{La date d'accès ne vaut pas grand-chose.} Les versions récentes de
Windows ne mettent pas toujours à jour la date de dernier accès, pour des
raisons de performance, et un antivirus qui lit un fichier la modifie. On ne
fonde jamais une conclusion « l'attaquant a ouvert ce fichier » sur la seule
date d'accès.
\end{boxinfo}

\subsection{Les journaux de NTFS}

Le journal \texttt{\$UsnJrnl} enregistre les changements des fichiers
(création, renommage, suppression), avec leur heure, tant que sa taille le
permet. Il permet de savoir qu'un fichier a existé, sous quel nom, et quand
il a été supprimé, même si son enregistrement MFT a été réutilisé.

% ============================================================================
\section{Les Fichiers Supprimés}
\label{sec:c04-supprimes}
% ============================================================================

Supprimer un fichier, sous NTFS, revient à marquer son enregistrement MFT et
ses blocs comme libres. Les données restent sur le disque jusqu'à ce qu'un
autre fichier les écrase. Deux méthodes permettent de les retrouver.

\begin{tableaucours}{L{3.2cm}L{4.6cm}L{4.4cm}}{%
    \tableauHeader{Méthode} &
    \tableauHeader{Principe} &
    \tableauHeader{Limite}
}
Par les métadonnées & L'enregistrement MFT libre indique encore le nom, les
dates et l'emplacement des données & Échoue si l'enregistrement a été
réutilisé \\
Par les signatures (\textit{carving}) & On cherche dans l'espace non alloué
les octets qui marquent le début d'un format (\texttt{PK} pour une archive
ZIP, \texttt{7z} pour une archive 7-Zip, \texttt{\%PDF} pour un PDF) & Ni
nom ni date~; fichiers fragmentés mal reconstitués \\
\end{tableaucours}
\vspace{-0.8em}
\begin{center}{\small\itshape Tableau~4.1. Deux façons de récupérer un fichier}\end{center}

Sur un disque SSD, la commande TRIM réduit fortement les chances de
récupération (Chapitre~\ref{chap:acquisition}). Sur un serveur virtualisé, le
comportement dépend de l'hyperviseur.

% ============================================================================
\section{Les Artefacts Windows}
\label{sec:c04-artefacts}
% ============================================================================

Un \termecle{artefact} est une trace qu'un système laisse en fonctionnant
normalement, sans qu'on lui demande de garder une preuve~: une entrée de
journal, une clé de registre, un raccourci. L'enquêteur les choisit d'après
la question.

\begin{tableaucours}{L{3.0cm}L{4.4cm}L{4.8cm}}{%
    \tableauHeader{Question} &
    \tableauHeader{Artefact} &
    \tableauHeader{Ce qu'il dit}
}
Qui s'est connecté, quand, comment~? & Journal de sécurité~: événements
4624 (ouverture réussie), 4625 (échec), type d'ouverture 10 pour le bureau
à distance & Compte, heure, adresse d'origine \\
Quel programme a été exécuté~? & Événement 4688 si l'audit est activé~;
Amcache~; Prefetch sur les postes de travail & Nom, chemin, heure, parfois
empreinte \\
Quel service a été installé~? & Événement 7045 du journal système & Nom du
service, binaire, compte \\
Quel fichier a été ouvert~? & Raccourcis \texttt{.lnk} du dossier
\texttt{Recent}, listes de raccourcis, clés \texttt{RecentDocs} du
registre de l'utilisateur & Chemin du fichier, heure d'ouverture \\
Le journal a-t-il été effacé~? & Événement 1102 & Effacement du journal
de sécurité, compte qui l'a fait \\
Un compte a-t-il été créé~? & Événement 4720 & Nom du compte, créateur \\
\end{tableaucours}
\vspace{-0.8em}
\begin{center}{\small\itshape Tableau~4.2. Questions et artefacts Windows}\end{center}

Deux précautions accompagnent ce tableau. Les artefacts dépendent de la
configuration~: le Prefetch, utile sur un poste de travail, est désactivé
par défaut sur les versions serveur de Windows, et l'événement 4688 n'existe
que si l'audit de la création de processus a été activé. Et chaque artefact
dit moins qu'il ne paraît~: une entrée Amcache prouve qu'un programme était
présent et a été vu par le système, pas toujours qu'il a été exécuté à
l'heure indiquée. On croise donc les artefacts entre eux.

% ============================================================================
\section{Ce que dit le Disque de SRV-CPT01}
\label{sec:c04-nyanga}
% ============================================================================

\begin{boxscenario}{Opération NYANGA. Trois questions, trois réponses}
\textbf{Question 1~: connexion.} Le journal de sécurité contient, le
12~mars 2025 à 22~h~39, un événement 4624 de type 10 pour le compte
\texttt{adm.kamga}, depuis l'adresse interne de SRV-AD01. Aucune ouverture
de ce compte sur SRV-CPT01 n'existe dans les trente jours précédents.

\textbf{Question 2~: ouverture du fichier de clé.} Le dossier \texttt{Recent}
du profil \texttt{adm.kamga} contient un raccourci \texttt{api.ini.lnk},
créé le 12~mars à 22~h~44, qui pointe vers
\texttt{C:\textbackslash PaySuite\textbackslash conf\textbackslash api.ini}.
La clé \texttt{RecentDocs} du registre de l'utilisateur cite le même
fichier.

\textbf{Question 3~: copie et effacement.} Le journal \texttt{\$UsnJrnl}
enregistre la création, à 22~h~52, d'un fichier \texttt{cpt.7z} dans le
dossier temporaire de l'utilisateur, puis sa suppression à 23~h~06.
L'enregistrement MFT est libre mais intact~: l'archive est récupérée, 41~Mo,
protégée par un mot de passe. Sa liste de fichiers, lisible sans le mot de
passe, comprend \texttt{api.ini} et le dossier des exports comptables de
février.

\medskip
\textbf{Analyse} (confiance élevée). Le compte \texttt{adm.kamga} a été
utilisé depuis le contrôleur de domaine pour ouvrir le fichier de clé, puis
pour l'archiver avec des exports comptables, avant d'effacer l'archive. Le
journal de sécurité ne montre aucun effacement (pas d'événement 1102)~:
l'attaquant a supprimé son archive, pas les journaux.
\end{boxscenario}

\begin{boxobjet}{Raccourci api.ini.lnk, profil adm.kamga, SRV-CPT01}
\textbf{Source~:} image NC-03, volume système,
\texttt{Users\textbackslash adm.kamga\textbackslash AppData\textbackslash
Roaming\textbackslash Microsoft\textbackslash Windows\textbackslash
Recent}.

\textbf{Contenu~:} chemin de la cible, taille et dates de la cible au moment
de l'ouverture, nom et numéro de série du volume, dates propres du
raccourci (création le 12/03/2025 à 22~h~44).

\textbf{Ce qu'il prouve~:} qu'une session du profil \texttt{adm.kamga} a
ouvert \texttt{api.ini} à cette heure, par l'interface de Windows.

\textbf{Ce qu'il ne prouve pas~:} qui était assis devant l'écran~; il ne dit
pas non plus que la clé a été lue en entier ni recopiée.
\end{boxobjet}

\begin{boxlecture}{Les trois questions au disque de SRV-CPT01}
\begin{tableaucours}{L{6.6cm}cL{3.4cm}}{%
    \tableauHeader{Proposition} &
    \tableauHeader{Cote} &
    \tableauHeader{Sources}
}
Le compte \texttt{adm.kamga} s'est connecté le 12~mars à 22~h~39 & \OUI{} &
4624 type 10 \\
La connexion venait de SRV-AD01 & \OUI{} & 4624 (adresse source) \\
Le fichier \texttt{api.ini} a été ouvert dans cette session & \OUI{} &
\texttt{.lnk}, \texttt{RecentDocs} \\
La clé a été copiée hors du serveur & \NEK{} & Archive créée, aucune trace de
transfert sur ce disque \\
Les journaux de SRV-CPT01 ont été effacés & \NON{} & Pas de 1102, journal
continu \\
M.~KAMGA lui-même était aux commandes & \NEK{} & Aucune pièce sur la
personne \\
\end{tableaucours}

\medskip
\textbf{Lecture.} Trois \OUI{}, un \NON{}, deux \NEK{}. Le disque répond
aux questions de connexion et d'ouverture~; il ne répond pas à celle du
transfert, qui relève de l'analyse du réseau (Chapitre~\ref{chap:reseau}).
\end{boxlecture}

% ============================================================================
\section{Les Heures et leurs Pièges}
\label{sec:c04-heures}
% ============================================================================

NTFS enregistre ses dates en temps universel (UTC). L'outil d'analyse les
affiche souvent dans le fuseau de la station d'examen, et les journaux
Windows dans celui que l'on choisit. Le Cameroun est à UTC+1, sans heure
d'été. Une heure lue sans son fuseau peut donc être fausse d'une heure, et
une chronologie qui mélange les deux devient incohérente.

\begin{tableaucours}{L{4.0cm}L{8.2cm}}{%
    \tableauHeader{Piège} &
    \tableauHeader{Parade}
}
Fuseau implicite & Noter le fuseau de chaque source, convertir tout en UTC
dans les notes de travail \\
Horloge déréglée & Mesurer la dérive (heure système contre heure de
référence) au moment de l'acquisition \\
Date modifiée par l'attaquant & Comparer \texttt{\$STANDARD\_INFORMATION} et
\texttt{\$FILE\_NAME}, recouper avec \texttt{\$UsnJrnl} et les journaux \\
Résolution différente & Les journaux à la seconde, NTFS à la centaine de
nanosecondes~: on compare à la précision la plus faible \\
\end{tableaucours}
\vspace{-0.8em}
\begin{center}{\small\itshape Tableau~4.3. Les pièges de l'heure}\end{center}

À NYANGA CRÉDIT, les serveurs se synchronisaient sur un serveur de temps
interne, lui-même calé sur une source publique~; la dérive mesurée le
15~mars sur SRV-CPT01 était de 2~secondes. Les heures du chapitre sont
données en heure locale (UTC+1), comme partout dans le livre.

% ============================================================================
\section{Ce que l'Analyse du Disque demande à chaque Métier}
\label{sec:c04-metiers}
% ============================================================================

\begin{boxmetier}{Analyste en réponse à incident}
\textbf{Examiner un disque}
\begin{enumerate}
    \item Écrire les questions avant d'ouvrir l'image.
    \item Pour chaque question, choisir les artefacts et vérifier qu'ils
          existent dans cette configuration.
    \item Croiser au moins deux artefacts indépendants pour chaque
          conclusion.
    \item Noter le fuseau de chaque heure.
    \item Dire, pour chaque question, ce qui a été lu et ce qui ne l'a pas
          été.
\end{enumerate}
\end{boxmetier}

\begin{boxmetier}{Application OPJ/Expert Judiciaire}
L'expert sépare ce que le disque établit sur un compte de ce qu'il ne dit
pas sur une personne. « Le compte \texttt{adm.kamga} a ouvert le fichier » est
une constatation~; « M.~KAMGA a ouvert le fichier » n'en est pas une.
\end{boxmetier}

\begin{boxmetier}{RSSI}
Activer, avant tout incident, l'audit de la création de processus et
dimensionner les journaux pour qu'ils couvrent au moins trois mois.
Sans l'événement 4688, l'enquête de NYANGA a dû se contenter d'Amcache pour
dater les exécutions.
\end{boxmetier}

\separateur

\begin{boxresume}
\textbf{Ce qu'il faut retenir de ce chapitre~:}
\begin{itemize}
    \item Un disque se lit par couches~: partitions, système de fichiers,
          fichiers, artefacts.
    \item NTFS garde deux jeux de dates (\texttt{\$STANDARD\_INFORMATION} et
          \texttt{\$FILE\_NAME})~; la date d'accès est peu fiable.
    \item Un fichier supprimé se récupère par ses métadonnées tant que
          l'enregistrement n'est pas réutilisé, sinon par ses signatures.
    \item Les artefacts Windows répondent à des questions précises~;
          ils dépendent de la configuration et se croisent.
    \item À NYANGA, le disque établit la connexion de \texttt{adm.kamga}
          depuis SRV-AD01 et l'ouverture du fichier de clé~; le transfert
          reste \NEK{} sur ce seul disque.
\end{itemize}
\end{boxresume}

\begin{boxinfo}
\textbf{Points de contrôle.} Vous avez compris ce chapitre si vous savez~:
\begin{itemize}
    \item citer l'artefact qui répond à une question donnée~;
    \item expliquer pourquoi un fichier supprimé est souvent récupérable~;
    \item convertir une heure entre UTC et l'heure du Cameroun, et dire
          pourquoi c'est nécessaire.
\end{itemize}
\end{boxinfo}

% ============================================================================
\section*{Exercices}
% ============================================================================

\begin{boxexo}[title={\ding{45}\quad Exercice~4.1. Quel artefact~? \niveaubase}]
Pour chaque question, citez l'artefact Windows qui y répond le mieux~:
(a) un compte a-t-il été créé le 9~mars sur SRV-AD01~? (b) quelqu'un s'est-il
connecté par le bureau à distance~? (c) un service inconnu a-t-il été
installé~? (d) le journal de sécurité a-t-il été vidé~? (e) quel document
Word l'utilisateur a-t-il ouvert en dernier~?
\end{boxexo}

\begin{boxexo}[title={\ding{45}\quad Exercice~4.2. Récupérer cpt.7z \niveaumoyen}]
\begin{enumerate}[(a)]
    \item Expliquez pourquoi l'archive \texttt{cpt.7z} a pu être récupérée
          avec son nom et ses dates.
    \item Si l'enregistrement MFT avait été réutilisé, quelle méthode
          resterait-il, et qu'aurait-on perdu~?
    \item L'archive est chiffrée. Que peut-on en tirer sans le mot de passe,
          et que faire pour l'obtenir~?
\end{enumerate}
\end{boxexo}

\begin{boxexo}[title={\ding{45}\quad Exercice~4.3. Deux heures pour un fichier \niveaumoyen}]
L'outil d'analyse affiche, pour \texttt{api.ini.lnk}, une création à
« 21:44:07 », sans indication de fuseau. Le journal de sécurité, exporté en
heure locale, donne l'ouverture de session à 22:39:12.
\begin{enumerate}[(a)]
    \item Le raccourci a-t-il été créé avant l'ouverture de session~?
    \item Comment lever le doute~?
    \item Réécrivez les deux heures en UTC et en heure du Cameroun.
\end{enumerate}
\end{boxexo}

\begin{boxexo}[title={\ding{45}\quad Exercice~4.4. Du compte à la personne \niveauexpert}]
Le directeur général lit l'encadré « Les trois questions au disque de
SRV-CPT01 » et demande~: « Donc
M.~KAMGA est coupable~? »
\begin{enumerate}[(a)]
    \item Répondez-lui en séparant constatation et analyse.
    \item Quelles pièces, sur ce disque ou ailleurs, pourraient faire passer
          la dernière ligne de \NEK{} à \OUI{} ou à \NON{}~?
    \item Quelle pièce du Chapitre~\ref{chap:investigation} rend déjà peu
          probable une action de M.~KAMGA depuis le siège~?
\end{enumerate}
\end{boxexo}

\bigskip

\begin{boxinfo}
\textbf{Pour aller plus loin.}
\begin{itemize}
    \item B.~Carrier, \textit{File System Forensic
          Analysis}~\cite{carrier2005}, pour NTFS et ext.
    \item The Sleuth Kit et Autopsy~\cite{sleuthkit}, outils libres
          d'analyse de disque.
\end{itemize}
\end{boxinfo}

\begin{boxinfo}
\textbf{Transition.} La session de \texttt{adm.kamga} venait du contrôleur
de domaine. Comment l'attaquant y a-t-il obtenu ces identifiants~? La
réponse est dans la mémoire de SRV-AD01~: Chapitre~\ref{chap:memoire}.
\end{boxinfo}

        % ============================================================================
% Fichier  : partie2/P02_C05_analyse_memoire.tex
% Titre    : L'Analyse de la Mémoire
% Auteur   : MINKA MI NGUIDJOÏ Thierry Emmanuel
% Version  : 1.0 -- novembre 2026
% Corrigés : corriges/Corr_C05.tex
% ============================================================================
% Sources : Ligh et al. (2014) ; Volatility 3 (greffons windows.info,
% pslist, psscan, pstree, cmdline, netscan, malfind, svcscan) ; structures
% EPROCESS et VAD de Windows (présentées sans citation de la documentation).
% Faits nouveaux (bible) : sur SRV-AD01 (capture du 14/03 à 10 h 02),
% processus « svchost.exe » PID 4412 lancé depuis C:\ProgramData, parent
% services.exe, service WinUpdSvc créé le 09/03 à 1 h 35 ; connexion établie
% vers 203.0.113.47:443 ; fichier lsass.dmp créé le 09/03 à 1 h 37 sur le
% disque de SRV-AD01 par rundll32.exe (commande retrouvée dans les journaux).
% ============================================================================

\chapter{L'Analyse de la Mémoire}
\label{chap:memoire}
\chaptermark{L'analyse de la mémoire}

\epigraph{%
    « Un programme peut effacer son fichier.\\
    Il ne peut pas tourner sans être en mémoire. »%
}{Principe du cours}

\niveauChapitre{Intermédiaire}{%
    Maîtrise des Chapitres~\ref{chap:acquisition} et~\ref{chap:disque}
}

\begin{boxobjectifs}
\begin{itemize}
    \item Expliquer ce que la mémoire vive contient et que le disque ne
          garde pas.
    \item Identifier le profil d'une capture et en vérifier la cohérence.
    \item Lister les processus par deux méthodes et interpréter leurs
          écarts.
    \item Repérer un processus anormal par son parent, son chemin et sa
          ligne de commande.
    \item Relier un processus à ses connexions réseau et à sa persistance.
    \item Reconnaître les traces d'un vol d'identifiants, sans en
          reproduire la technique.
\end{itemize}
\end{boxobjectifs}

\bigskip

% ============================================================================
\section*{Avant de commencer : un svchost de trop}
% ============================================================================

Yaoundé, 22 mars 2025. M.~ABOUBAKAR, l'expert réquisitionné par le
procureur, a reçu le scellé NC-01, l'a ouvert devant témoin, a vérifié son
empreinte et travaille sur sa copie~: 16~Go de mémoire du contrôleur de
domaine, capturés le 14~mars à 10~h~02. Il compte, dans la liste des
processus, quatorze \texttt{svchost.exe}. Treize sont lancés depuis
\texttt{C:\textbackslash Windows\textbackslash System32}. Le quatorzième
tourne depuis \texttt{C:\textbackslash ProgramData}, et il est le seul à
entretenir une connexion vers l'extérieur.

\begin{boxalert}
\textbf{Violation typique.} Conclure qu'un processus est légitime parce que
son nom est celui d'un programme du système.

\smallskip
\textbf{Règle~:} un processus se juge sur son chemin, son parent, sa ligne de
commande, son heure de démarrage et son comportement réseau, comparés à ce
qu'on attend de ce programme. Le nom est le premier attribut qu'un attaquant
imite.
\end{boxalert}

\separateur

% ============================================================================
\section{Ce que la Mémoire Contient}
\label{sec:c05-contenu}
% ============================================================================

La mémoire vive d'un serveur Windows contient, à l'instant de la capture,
l'état complet du système~: les processus et leurs fils d'exécution, les
bibliothèques chargées, les connexions réseau, les objets ouverts (fichiers,
clés de registre), le cache des fichiers, une partie du registre, et les
données que les programmes manipulent, y compris des mots de passe, des clés
de chiffrement et des commandes saisies. Certains programmes malveillants
n'existent qu'en mémoire~: ils sont chargés depuis le réseau, ou injectés
dans un processus légitime, sans jamais être écrits sur le disque.

\begin{tableaucours}{L{4.4cm}L{3.6cm}L{4.2cm}}{%
    \tableauHeader{Question} &
    \tableauHeader{Mémoire} &
    \tableauHeader{Disque}
}
Quels programmes tournaient~? & Oui, avec leurs parents & Indirectement
(Amcache, journaux) \\
Quelles connexions étaient ouvertes~? & Oui, avec le processus & Non \\
Quelle ligne de commande~? & Oui & Seulement si l'audit l'enregistre \\
Code injecté dans un processus~? & Oui & Rarement \\
Ce qui s'est passé il y a un mois~? & Non & Oui, par les journaux et
les dates \\
\end{tableaucours}
\vspace{-0.8em}
\begin{center}{\small\itshape Tableau~5.1. Mémoire et disque~: deux sources complémentaires}\end{center}

% ============================================================================
\section{L'Outil et la Méthode}
\label{sec:c05-outil}
% ============================================================================

L'outil de référence, libre, est Volatility, dans sa version~3~\cite{volatility3}.
Il reconstruit les structures du système d'exploitation à partir de la
capture brute, grâce à des tables de symboles propres à chaque version du
noyau. On l'emploie greffon par greffon, chacun répondant à une question.

\begin{tableaucours}{L{3.6cm}L{8.6cm}}{%
    \tableauHeader{Greffon} &
    \tableauHeader{Question}
}
\texttt{windows.info} & La capture est-elle lisible, de quelle version de
Windows, à quelle heure système~? \\
\texttt{windows.pslist} & Quels processus figurent dans la liste tenue par
le noyau~? \\
\texttt{windows.psscan} & Quelles structures de processus trouve-t-on en
balayant toute la mémoire, y compris des processus terminés ou cachés~? \\
\texttt{windows.pstree} & Qui a lancé qui~? \\
\texttt{windows.cmdline} & Avec quelle ligne de commande~? \\
\texttt{windows.netscan} & Quelles connexions et quels ports, rattachés à
quel processus~? \\
\texttt{windows.malfind} & Quelles zones de mémoire exécutables et
modifiables, sans fichier associé, suggèrent une injection de code~? \\
\texttt{windows.svcscan} & Quels services sont enregistrés, avec quel
binaire~? \\
\end{tableaucours}
\vspace{-0.8em}
\begin{center}{\small\itshape Tableau~5.2. Les greffons de base de Volatility~3}\end{center}

\subsection{Deux listes de processus}

Windows tient la liste des processus actifs dans une liste chaînée de
structures \texttt{EPROCESS}. Le greffon \texttt{pslist} suit cette liste.
Le greffon \texttt{psscan}, lui, cherche ces structures partout dans la
capture, par leur forme. La comparaison des deux est instructive~: un
processus présent dans \texttt{psscan} et absent de \texttt{pslist} est soit
terminé (sa structure traîne encore en mémoire, avec une heure de fin), soit
caché par un programme qui l'a retiré de la liste. Le premier cas est
banal~; le second est un signe fort de compromission.

\subsection{Le processus normal}

Pour juger un processus, on compare ses attributs à ce que l'on sait des
processus du système. Quelques repères suffisent pour les plus imités.

\begin{tableaucours}{L{2.6cm}L{3.8cm}L{2.6cm}L{3.2cm}}{%
    \tableauHeader{Nom} &
    \tableauHeader{Chemin attendu} &
    \tableauHeader{Parent attendu} &
    \tableauHeader{Nombre}
}
\texttt{lsass.exe} & \texttt{System32} & \texttt{wininit.exe} & Un seul \\
\texttt{services.exe} & \texttt{System32} & \texttt{wininit.exe} & Un seul \\
\texttt{svchost.exe} & \texttt{System32} & \texttt{services.exe} &
Plusieurs \\
\texttt{explorer.exe} & \texttt{Windows} & Variable (lancé à la session) &
Un par session \\
\end{tableaucours}
\vspace{-0.8em}
\begin{center}{\small\itshape Tableau~5.3. Quelques processus et leurs attributs attendus}\end{center}

% ============================================================================
\section{Ce que la Mémoire de SRV-AD01 Révèle}
\label{sec:c05-nyanga}
% ============================================================================

\begin{boxscenario}{Opération NYANGA. Le processus 4412}
\textbf{Constatations} (copie de NC-01, Volatility~3).
\begin{itemize}
    \item \texttt{windows.info}~: Windows Server 2019, heure système de la
          capture 14/03/2025 09:02:11 UTC, soit 10:02:11 heure locale,
          cohérente avec le registre des actions.
    \item \texttt{pslist} et \texttt{psscan}~: quatorze processus
          \texttt{svchost.exe}~; le PID~4412 a pour chemin
          \texttt{C:\textbackslash ProgramData\textbackslash
          WinUpd\textbackslash svchost.exe}, pour parent
          \texttt{services.exe}, et a démarré le 09/03 à 1~h~35.
    \item \texttt{svcscan}~: un service \texttt{WinUpdSvc}, en démarrage
          automatique, a pour binaire ce même fichier.
    \item \texttt{netscan}~: le PID~4412 maintient une connexion établie vers
          203.0.113.47, port~443.
    \item \texttt{malfind}~: aucune zone suspecte dans \texttt{lsass.exe}~;
          deux zones exécutables sans fichier associé dans le PID~4412.
\end{itemize}

\textbf{Analyse} (confiance élevée). Le PID~4412 est un programme étranger
au système, déguisé sous le nom \texttt{svchost.exe}, installé comme service
pour redémarrer avec la machine, et en communication avec l'adresse vers
laquelle sont partis les 3,2~Go. C'est l'outil d'accès à distance de
l'attaquant sur le contrôleur de domaine.
\end{boxscenario}

Le parent \texttt{services.exe} ne disculpe pas le processus~: un service
est toujours lancé par \texttt{services.exe}, y compris un service créé par
un attaquant. C'est le chemin qui trahit le PID~4412, et le service qui
explique son parent.

\begin{boxobjet}{Processus 4412, capture mémoire de SRV-AD01}
\textbf{Source~:} scellé NC-01, copie de travail, Volatility~3.

\textbf{Attributs~:} nom \texttt{svchost.exe}~; chemin
\texttt{C:\textbackslash ProgramData\textbackslash WinUpd}~; parent
\texttt{services.exe}~; démarrage 09/03/2025 1~h~35~; service
\texttt{WinUpdSvc}~; connexion vers 203.0.113.47:443.

\textbf{Ce qu'il prouve~:} la présence, le 14~mars, d'un programme
persistant étranger au système, en communication avec l'adresse de
l'exfiltration.

\textbf{Ce qu'il ne prouve pas~:} comment le programme est arrivé sur la
machine, ni qui le pilotait.
\end{boxobjet}

% ============================================================================
\section{Le Vol d'Identifiants}
\label{sec:c05-identifiants}
% ============================================================================

Sous Windows, le processus \texttt{lsass.exe} gère les authentifications et
garde en mémoire des éléments d'identification des sessions ouvertes. Un
attaquant qui obtient les droits d'administration sur une machine peut
copier la mémoire de ce processus dans un fichier, puis l'emporter et
l'analyser chez lui pour en extraire des identifiants. Le cours ne décrit
pas la technique d'extraction~: elle relève de l'UE MSN-2S3-SEC03. Il
apprend à en reconnaître les traces.

\begin{tableaucours}{L{4.4cm}L{7.8cm}}{%
    \tableauHeader{Trace} &
    \tableauHeader{Où la chercher}
}
Fichier de vidage (\texttt{.dmp}) de \texttt{lsass.exe} & Disque~:
\texttt{\$UsnJrnl}, MFT, dossiers temporaires \\
Ligne de commande d'un outil de vidage & Mémoire (\texttt{cmdline}) si le
processus tourne encore~; événement 4688 si l'audit est actif \\
Accès au processus \texttt{lsass.exe} par un autre processus & Journaux de
surveillance avancée, s'ils existent \\
Réutilisation des identifiants & Ouvertures de session du compte volé depuis
des machines inhabituelles (4624) \\
\end{tableaucours}
\vspace{-0.8em}
\begin{center}{\small\itshape Tableau~5.4. Les traces d'un vol d'identifiants}\end{center}

\begin{boxscenario}{Opération NYANGA. Le fichier du 9 mars}
\textbf{Constatations.} Sur l'image du disque de SRV-AD01 (acquise le
16~mars), le \texttt{\$UsnJrnl} enregistre la création, le 9~mars à 1~h~37,
d'un fichier \texttt{lsass.dmp} de 54~Mo dans
\texttt{C:\textbackslash Windows\textbackslash Temp}, et sa suppression à
1~h~49. Le journal système enregistre, à 1~h~35, la création du service
\texttt{WinUpdSvc} (événement 7045) par le compte \texttt{svc\_backup}. À
partir du 10~mars, le compte \texttt{adm.kamga} ouvre des sessions depuis
SRV-AD01 vers les autres serveurs.

\textbf{Analyse} (confiance élevée). L'attaquant, entré sous
\texttt{svc\_backup}, a installé son outil le 9~mars, copié la mémoire de
\texttt{lsass.exe} deux minutes plus tard, puis utilisé les identifiants de
\texttt{adm.kamga}, présents en mémoire parce que l'administrateur s'était
connecté au contrôleur de domaine dans les jours précédents. Cette séquence
explique l'usage du compte sans la participation de son titulaire~; elle ne
l'exclut pas.
\end{boxscenario}

\begin{boxlecture}{La mémoire de SRV-AD01 et le disque associé}
\begin{tableaucours}{L{7.2cm}c}{%
    \tableauHeader{Proposition} &
    \tableauHeader{Cote}
}
Un programme étranger tournait sur SRV-AD01 le 14~mars & \OUI{} \\
Il communiquait avec 203.0.113.47 & \OUI{} \\
Il était installé comme service depuis le 9~mars & \OUI{} \\
La mémoire de \texttt{lsass.exe} a été copiée le 9~mars & \OUI{} \\
Les identifiants de \texttt{adm.kamga} ont été extraits de cette copie &
\NEK{} \\
Du code a été injecté dans \texttt{lsass.exe} & \NON{} \\
\end{tableaucours}

\medskip
\textbf{Lecture.} Quatre \OUI{}, un \NON{}, un \NEK{}. L'extraction des
identifiants depuis la copie reste \NEK{}~: elle a eu lieu hors du
système, chez l'attaquant, et aucune pièce ne peut la montrer. On l'infère
de la séquence, avec une confiance élevée, mais on ne la constate pas.
\end{boxlecture}

% ============================================================================
\section{Les Limites de l'Analyse Mémoire}
\label{sec:c05-limites}
% ============================================================================

\begin{itemize}
    \item Une capture est un instantané~: elle ne dit rien de ce qui a
          tourné puis s'est terminé, sauf traces résiduelles.
    \item Les tables de symboles doivent correspondre exactement à la
          version du noyau~; une version inconnue de l'outil donne des
          résultats faux ou vides.
    \item La capture elle-même peut être incohérente, parce que la mémoire a
          changé pendant qu'on la lisait~: des structures à moitié copiées
          apparaissent.
    \item Les zones signalées par \texttt{malfind} ne sont pas toutes
          malveillantes~: certains programmes légitimes (machines
          virtuelles Java, navigateurs) génèrent du code en mémoire.
\end{itemize}

% ============================================================================
\section{Ce que l'Analyse Mémoire demande à chaque Métier}
\label{sec:c05-metiers}
% ============================================================================

\begin{boxmetier}{Analyste en réponse à incident}
\textbf{Examiner une capture mémoire}
\begin{enumerate}
    \item Vérifier le profil et l'heure système avec \texttt{windows.info}.
    \item Comparer \texttt{pslist} et \texttt{psscan}.
    \item Examiner l'arbre des processus, les chemins et les lignes de
          commande.
    \item Rattacher chaque connexion à son processus.
    \item Chercher la persistance (services, tâches) du processus suspect.
    \item Confirmer chaque découverte par une autre source (disque, réseau).
\end{enumerate}
\end{boxmetier}

\begin{boxmetier}{Application OPJ/Expert Judiciaire}
L'expert indique la version de l'outil et des tables de symboles, la liste
des greffons et leur ordre, et conserve leurs sorties brutes en annexe~: un
autre expert doit pouvoir refaire l'examen sur la même copie et obtenir les
mêmes listes.
\end{boxmetier}

\begin{boxmetier}{RSSI}
Limiter les connexions des comptes d'administration du domaine aux seules
machines d'administration~: un identifiant qui ne se trouve pas en mémoire
d'un serveur ne peut pas y être volé.
\end{boxmetier}

\separateur

\begin{boxresume}
\textbf{Ce qu'il faut retenir de ce chapitre~:}
\begin{itemize}
    \item La mémoire montre ce qui tournait, avec quels parents, quelles
          commandes et quelles connexions~; le disque montre l'histoire.
    \item Volatility s'emploie greffon par greffon, chacun répondant à une
          question~; \texttt{pslist} et \texttt{psscan} se comparent.
    \item Un processus se juge sur son chemin, son parent, sa commande, son
          heure et son réseau, jamais sur son nom.
    \item Le vol d'identifiants se reconnaît à ses traces (fichier de
          vidage, réutilisation du compte)~; l'extraction elle-même reste
          souvent \NEK{}.
    \item À NYANGA, le PID~4412, faux \texttt{svchost.exe} installé comme
          service le 9~mars, communique avec 203.0.113.47.
\end{itemize}
\end{boxresume}

\begin{boxinfo}
\textbf{Points de contrôle.} Vous avez compris ce chapitre si vous savez~:
\begin{itemize}
    \item interpréter un écart entre \texttt{pslist} et \texttt{psscan}~;
    \item dire pourquoi un processus est suspect à partir de ses attributs~;
    \item distinguer ce qui est constaté de ce qui est inféré dans un vol
          d'identifiants.
\end{itemize}
\end{boxinfo}

% ============================================================================
\section*{Exercices}
% ============================================================================

\begin{boxexo}[title={\ding{45}\quad Exercice~5.1. Normal ou suspect \niveaubase}]
Pour chaque processus, dites s'il est normal ou suspect, et pourquoi~:
(a) \texttt{lsass.exe}, \texttt{System32}, parent \texttt{wininit.exe}~;
(b) \texttt{lsass.exe}, \texttt{System32}, parent \texttt{explorer.exe}~;
(c) \texttt{svchost.exe}, \texttt{System32}, parent \texttt{services.exe}~;
(d) \texttt{svchost.exe}, \texttt{Users\textbackslash Public}, parent
\texttt{cmd.exe}~; (e) deux processus \texttt{lsass.exe}.
\end{boxexo}

\begin{boxexo}[title={\ding{45}\quad Exercice~5.2. pslist et psscan \niveaumoyen}]
\texttt{psscan} trouve trois processus absents de \texttt{pslist}~: un
\texttt{rundll32.exe} terminé le 9~mars à 1~h~38~; un \texttt{cmd.exe}
terminé le 13~mars~; un \texttt{svch0st.exe} sans heure de fin.
\begin{enumerate}[(a)]
    \item Interprétez chacun.
    \item Lequel demande une vérification immédiate, et laquelle~?
    \item Que pourrait apporter la ligne de commande du \texttt{rundll32.exe},
          si elle est encore lisible~?
\end{enumerate}
\end{boxexo}

\begin{boxexo}[title={\ding{45}\quad Exercice~5.3. Relier les sources \niveaumoyen}]
Construisez un tableau qui relie, pour le PID~4412, chaque attribut constaté
en mémoire à une confirmation sur le disque de SRV-AD01 ou dans la capture
réseau. Indiquez pour chaque ligne si la confirmation existe, manque, ou
reste à chercher.
\end{boxexo}

\begin{boxexo}[title={\ding{45}\quad Exercice~5.4. Ce qu'on ne peut pas montrer \niveauexpert}]
L'avocat de M.~KAMGA soutient que la copie de la mémoire de
\texttt{lsass.exe} « ne prouve pas » le vol de ses identifiants.
\begin{enumerate}[(a)]
    \item A-t-il raison sur le plan de la constatation~? Sur celui de
          l'analyse~?
    \item Rédigez le passage du rapport d'expertise qui traite ce point.
    \item Quelle mesure, en place avant l'incident, aurait permis de
          constater l'usage des identifiants volés plutôt que de l'inférer~?
\end{enumerate}
\end{boxexo}

\bigskip

\begin{boxinfo}
\textbf{Pour aller plus loin.}
\begin{itemize}
    \item M.~H.~Ligh et al., \textit{The Art of Memory
          Forensics}~\cite{ligh2014memory}.
    \item La documentation de Volatility~3~\cite{volatility3}.
\end{itemize}
\end{boxinfo}

\begin{boxinfo}
\textbf{Transition.} Le PID~4412 parle à 203.0.113.47. Qu'a-t-il envoyé, et
quand~? La capture de la sonde réseau le dit~: Chapitre~\ref{chap:reseau}.
\end{boxinfo}

        % ============================================================================
% Fichier  : partie2/P02_C06_analyse_reseau.tex
% Titre    : L'Analyse du Réseau
% Auteur   : MINKA MI NGUIDJOÏ Thierry Emmanuel
% Version  : 1.0 -- novembre 2026
% Corrigés : corriges/Corr_C06.tex
% ============================================================================
% Sources : Bejtlich (2013) pour la supervision de la sécurité réseau ;
% Sanders (2017) pour l'analyse de paquets ; Wireshark, tshark, Zeek
% (présentés sans citation) ; RFC 5737 (adresses de documentation).
% Faits nouveaux (bible) : adresses internes SRV-AD01 10.20.1.10, SRV-FIC01
% 10.20.1.20, SRV-CPT01 10.20.1.30 ; lien Internet de 20 Mbit/s ; flux
% SRV-FIC01 -> 203.0.113.47:443 du 14/03 de 1 h 12 à 3 h 47, 3,2 Go montant,
% 11 Mo descendant, SNI cdn-update.example.net, certificat auto-signé du
% 01/03/2025 ; balise de SRV-AD01 toutes les 60 s (+/- 6 s) ; virements
% émis depuis 203.0.113.62, hors du réseau de NYANGA ; journal de l'outil de
% copie sur SRV-FIC01 : 18 412 fichiers du dossier Clients.
% ============================================================================

\chapter{L'Analyse du Réseau}
\label{chap:reseau}
\chaptermark{L'analyse du réseau}

\epigraph{%
    « Le chiffrement cache le contenu.\\
    Il ne cache ni la taille, ni l'heure, ni la destination. »%
}{Principe du cours}

\niveauChapitre{Intermédiaire}{%
    Maîtrise des Chapitres~\ref{chap:acquisition} et~\ref{chap:memoire}
}

\begin{boxobjectifs}
\begin{itemize}
    \item Distinguer les sources réseau~: capture complète, flux, journaux
          d'équipements, alertes de sonde.
    \item Lire une capture avec un analyseur de paquets et en extraire les
          conversations utiles.
    \item Caractériser une exfiltration par le volume, la durée, le débit
          et la destination.
    \item Reconnaître une balise de commande par sa régularité.
    \item Tirer d'un flux chiffré tout ce qu'il dit sans être déchiffré.
    \item Dire ce que l'absence d'un flux dans une capture établit, et ce
          qu'elle n'établit pas.
\end{itemize}
\end{boxobjectifs}

\bigskip

% ============================================================================
\section*{Avant de commencer : 48 heures de trafic}
% ============================================================================

Yaoundé, 24 mars 2025. La capture de la sonde de NYANGA CRÉDIT (scellé
NC-04) couvre le lien Internet du siège du 13~mars à 8~h au 14~mars à 8~h~:
37~Go de paquets. M.~FOTSO, analyste de SANKOFA EXPERTISE, a trois questions
de Mme~MBARGA~: qu'est-ce qui est sorti, vers où, et à quel moment~; le
programme du contrôleur de domaine parlait-il à la même adresse~; les
virements de la nuit sont-ils passés par ce lien~?

\begin{boxalert}
\textbf{Violation typique.} Conclure qu'un échange n'a pas eu lieu parce
qu'on ne le trouve pas dans la capture.

\smallskip
\textbf{Règle~:} une capture ne voit que ce qui passe par le point où elle
est placée, pendant la période qu'elle couvre. Avant d'écrire « aucun flux
vers telle adresse », on dit où était la sonde, ce qu'elle voyait, et si le
flux pouvait emprunter un autre chemin.
\end{boxalert}

\separateur

% ============================================================================
\section{Quatre Sources Réseau}
\label{sec:c06-sources}
% ============================================================================

\begin{tableaucours}{L{2.8cm}L{4.0cm}L{2.6cm}L{2.8cm}}{%
    \tableauHeader{Source} &
    \tableauHeader{Contenu} &
    \tableauHeader{Durée typique} &
    \tableauHeader{Limite}
}
Capture complète (PCAP) & Tous les paquets, en-têtes et contenu & Heures à
jours (volume) & Coût de stockage~; contenu souvent chiffré \\
Flux (NetFlow, IPFIX) & Qui a parlé à qui, quand, combien d'octets & Mois &
Pas de contenu \\
Journaux d'équipements & Pare-feu, proxy, DNS~: décisions et requêtes &
Semaines à mois & Ce que l'équipement a jugé utile de noter \\
Alertes de sonde & Événements jugés suspects par des règles & Mois & Ne
voit que ce que les règles décrivent \\
\end{tableaucours}
\vspace{-0.8em}
\begin{center}{\small\itshape Tableau~6.1. Les sources de l'analyse réseau}\end{center}

R.~Bejtlich, dans son ouvrage sur la supervision de la sécurité
réseau~\cite{bejtlich2013nsm}, montre que ces sources se complètent~: les
flux et les journaux couvrent une longue période sans contenu, la capture
complète couvre une période courte avec tout le détail. Une enquête commence
souvent par les flux, pour repérer la fenêtre utile, puis descend dans la
capture.

% ============================================================================
\section{Lire une Capture}
\label{sec:c06-lire}
% ============================================================================

L'analyseur de paquets Wireshark, et sa version en ligne de commande tshark,
décodent les protocoles et regroupent les paquets en
\termecle{conversations} (une paire d'adresses et de ports). Les outils de
supervision comme Zeek transforment une capture en journaux par protocole~:
connexions, requêtes DNS, échanges TLS, transferts HTTP. On travaille dans
cet ordre.

\begin{enumerate}
    \item Les statistiques de conversations~: quelles paires ont échangé le
          plus d'octets, dans quel sens.
    \item Les requêtes DNS~: quels noms ont été résolus, vers quelles
          adresses.
    \item Les échanges TLS~: nom de serveur annoncé (SNI), certificat
          présenté, versions.
    \item Le détail des paquets, seulement pour les conversations retenues.
\end{enumerate}

\begin{boxinfo}
\textbf{Toujours sur une copie, toujours avec un filtre écrit.} On ouvre la
copie de travail de la capture, jamais le scellé, et l'on consigne chaque
filtre appliqué (par exemple \texttt{ip.addr == 203.0.113.47}). Le filtre
fait partie de la méthode~: un autre analyste doit pouvoir le rejouer.
\end{boxinfo}

% ============================================================================
\section{Caractériser une Exfiltration}
\label{sec:c06-exfiltration}
% ============================================================================

Une exfiltration se reconnaît moins à son contenu, souvent chiffré, qu'à sa
forme~: un volume sortant important, un rapport entre octets envoyés et
reçus très déséquilibré, une destination inhabituelle, un horaire anormal,
un débit soutenu.

\begin{boxscenario}{Opération NYANGA. La conversation de 1 h 12}
\textbf{Constatations} (copie de NC-04).
\begin{itemize}
    \item Le 14~mars, de 1~h~12 à 3~h~47, une conversation TLS relie
          SRV-FIC01 (10.20.1.20) à 203.0.113.47, port~443.
    \item Octets envoyés par SRV-FIC01~: 3,2~Go~; octets reçus~: 11~Mo.
    \item Le client annonce le nom \texttt{cdn-update.example.net}~; le
          serveur présente un certificat auto-signé émis le 1\textsuperscript{er}~mars
          2025.
    \item Aucune autre conversation de SRV-FIC01 vers Internet sur les
          24~heures de la capture.
\end{itemize}

\textbf{Calcul.} Durée~: 2~h~35, soit 9\,300~s. Débit moyen~: $3{,}2 \times
10^{9} / 9\,300 \approx 344\,000$~octets par seconde, soit environ
2,75~Mbit/s, sur un lien de 20~Mbit/s. Rapport envoyé sur reçu~: environ
290.

\textbf{Analyse} (confiance élevée). Un serveur de fichiers n'envoie pas
3,2~Go en pleine nuit vers un serveur au certificat auto-signé récent qui
se fait passer pour un service de mise à jour. La forme du flux est celle
d'une exfiltration. Le débit, limité à un septième du lien, est compatible
avec un outil réglé pour ne pas saturer la ligne et passer inaperçu.
\end{boxscenario}

\subsection{Ce que le chiffrement laisse voir}

\begin{tableaucours}{L{4.2cm}L{8.0cm}}{%
    \tableauHeader{Élément visible} &
    \tableauHeader{Ce qu'il apprend}
}
Adresses, ports, heures & Qui, vers où, quand \\
Volumes dans chaque sens & Envoi ou réception, ordre de grandeur des
données \\
Nom annoncé (SNI) & Le service que le client prétend joindre \\
Certificat du serveur & Émetteur, date de création, auto-signature \\
Rythme des paquets & Transfert continu, balise régulière, session
interactive \\
\end{tableaucours}
\vspace{-0.8em}
\begin{center}{\small\itshape Tableau~6.2. Ce que laisse voir un flux chiffré}\end{center}

Le contenu, lui, reste inaccessible sans les clés. Pour savoir quels
fichiers sont partis, il faut une autre source. À NYANGA, elle se trouve sur
SRV-FIC01 lui-même~: l'outil de copie utilisé par l'attaquant a laissé un
journal dans le dossier de l'utilisateur, retrouvé dans l'acquisition du
18~mars (NC-07). Il liste 18\,412 fichiers du dossier \texttt{Clients},
pour un total de 3,19~Go, avec des heures de début et de fin qui encadrent
la conversation réseau.

% ============================================================================
\section{Reconnaître une Balise}
\label{sec:c06-balise}
% ============================================================================

\begin{boxdef}
\textbf{\termecle{Balise}} (\textit{beacon})~: connexion qu'un programme
malveillant établit à intervalles réguliers vers son serveur de commande,
pour signaler sa présence et recevoir des ordres. Sa régularité la trahit,
même quand le contenu est chiffré et le volume faible.
\end{boxdef}

\begin{boxscenario}{Opération NYANGA. Toutes les soixante secondes}
\textbf{Constatations.} Sur les 24~heures de la capture, SRV-AD01
(10.20.1.10) ouvre 1\,431 connexions courtes vers 203.0.113.47:443.
L'intervalle entre deux connexions est compris entre 54 et 66~secondes~;
chaque échange fait entre 1,2 et 2,8~Ko, sauf onze échanges de 40 à
900~Ko, tous situés entre 0~h~50 et 3~h~50 le 14~mars.

\textbf{Calcul.} 24~heures font 1\,440 minutes~: 1\,431 connexions
correspondent à une connexion par minute, à neuf près.

\textbf{Analyse} (confiance élevée). C'est la balise du PID~4412 vu en
mémoire (Chapitre~\ref{chap:memoire}), réglée sur 60~secondes avec une
variation de 10~\%. Les onze échanges plus volumineux, pendant la nuit,
correspondent à des ordres plus lourds~: ils précèdent et accompagnent
l'exfiltration depuis SRV-FIC01.
\end{boxscenario}

% ============================================================================
\section{Ce que la Capture ne Voit pas}
\label{sec:c06-absence}
% ============================================================================

La troisième question de Mme~MBARGA reçoit une réponse inattendue. La
capture ne contient aucun échange avec l'interface de paiement pendant la
nuit. Les journaux de l'interface (NC-05) indiquent que les vingt-trois
virements ont été émis depuis l'adresse 203.0.113.62, qui n'appartient pas à
NYANGA CRÉDIT. L'interface est hébergée par le prestataire, sur Internet~:
l'attaquant a utilisé la clé volée depuis sa propre machine, sans passer par
le réseau du siège.

\begin{boxlecture}{Les trois questions au réseau}
\begin{tableaucours}{L{7.0cm}cL{3.4cm}}{%
    \tableauHeader{Proposition} &
    \tableauHeader{Cote} &
    \tableauHeader{Sources}
}
3,2~Go sont sortis de SRV-FIC01 vers 203.0.113.47 la nuit du 14~mars &
\OUI{} & NC-04 \\
Les fichiers sortis sont ceux du dossier \texttt{Clients} & \OUI{} &
Journal de l'outil (NC-07) \\
Le programme de SRV-AD01 communiquait avec la même adresse & \OUI{} & NC-04,
NC-01 \\
Les virements sont passés par le lien du siège & \NON{} & NC-04, NC-05 \\
Aucune autre donnée n'est sortie par un autre chemin & \NEK{} & La sonde ne
voit que le lien Internet du siège \\
\end{tableaucours}

\medskip
\textbf{Lecture.} Trois \OUI{}, un \NON{}, un \NEK{}. Le \NON{} sur les
virements est solide parce que les deux sources se recoupent~: absence dans
la capture, et adresse d'émission externe dans les journaux de l'interface.
Le \NEK{} rappelle que les agences ont leur propre accès Internet, hors de la
vue de la sonde.
\end{boxlecture}

% ============================================================================
\section{Ce que l'Analyse Réseau demande à chaque Métier}
\label{sec:c06-metiers}
% ============================================================================

\begin{boxmetier}{Analyste en réponse à incident}
\textbf{Analyser une capture}
\begin{enumerate}
    \item Situer la sonde~: lien observé, période, pertes éventuelles.
    \item Partir des statistiques de conversations, pas des paquets.
    \item Pour chaque conversation suspecte~: volumes, durée, débit,
          rythme, nom annoncé, certificat.
    \item Rattacher chaque flux à un processus (mémoire) et à des fichiers
          (disque).
    \item Écrire les filtres utilisés.
\end{enumerate}
\end{boxmetier}

\begin{boxmetier}{RSSI}
Placer la sonde là où passent les flux qui comptent, et savoir ce qu'elle
ne voit pas. À NYANGA, les vingt-deux agences sortent sur Internet par leurs
propres accès~: elles sont hors du champ de la sonde.
\end{boxmetier}

\begin{boxmetier}{Application OPJ/Expert Judiciaire}
L'adresse 203.0.113.47 désigne une machine à un instant donné. Pour aller
vers une personne, l'enquête doit obtenir, par réquisition ou par entraide
internationale, les données du fournisseur qui l'exploitait~: le
Chapitre~\ref{chap:normes-droit} en donne les fondements.
\end{boxmetier}

\separateur

\begin{boxresume}
\textbf{Ce qu'il faut retenir de ce chapitre~:}
\begin{itemize}
    \item Capture complète, flux, journaux et alertes se complètent~; on
          part du large pour descendre vers le détail.
    \item Une exfiltration se caractérise par son volume, son déséquilibre,
          son débit, son horaire et sa destination~; une balise par sa
          régularité.
    \item Un flux chiffré laisse voir les adresses, les volumes, le nom
          annoncé, le certificat et le rythme~; son contenu se cherche
          ailleurs.
    \item L'absence d'un flux dans une capture ne vaut \NON{} que si l'on
          sait que le flux devait passer par la sonde.
    \item À NYANGA, 3,2~Go du dossier \texttt{Clients} sont partis de
          SRV-FIC01~; les virements, eux, ont été émis de l'extérieur.
\end{itemize}
\end{boxresume}

\begin{boxinfo}
\textbf{Points de contrôle.} Vous avez compris ce chapitre si vous savez~:
\begin{itemize}
    \item calculer le débit moyen d'une conversation~;
    \item reconnaître une balise dans une liste de connexions~;
    \item dire pourquoi l'absence d'un flux ne prouve pas toujours son
          inexistence.
\end{itemize}
\end{boxinfo}

% ============================================================================
\section*{Exercices}
% ============================================================================

\begin{boxexo}[title={\ding{45}\quad Exercice~6.1. Quelle source~? \niveaubase}]
Pour chaque question, dites quelle source réseau y répond le mieux~:
(a) SRV-FIC01 a-t-il parlé à 203.0.113.47 il y a trois mois~? (b) quel nom
de domaine le contrôleur de domaine a-t-il résolu avant sa première balise~?
(c) que contenait la requête HTTP non chiffrée d'un poste vers un site
inconnu~? (d) la sonde a-t-elle signalé quelque chose cette nuit-là~?
\end{boxexo}

\begin{boxexo}[title={\ding{45}\quad Exercice~6.2. Le calcul du débit \niveaumoyen}]
Une seconde conversation, découverte dans les flux d'une agence, montre
900~Mo envoyés en 1~h~15.
\begin{enumerate}[(a)]
    \item Calculez le débit moyen en octets par seconde et en Mbit/s.
    \item L'agence dispose d'un lien de 4~Mbit/s. Qu'en déduisez-vous sur le
          réglage de l'outil~?
    \item Cette conversation change-t-elle une cote de l'encadré « Les trois
          questions au réseau »~?
\end{enumerate}
\end{boxexo}

\begin{boxexo}[title={\ding{45}\quad Exercice~6.3. Une balise ou pas \niveaumoyen}]
Voici les heures de connexion d'un poste vers une même adresse~: 10:00:03,
10:05:01, 10:10:04, 10:15:02, 10:20:03, puis 11:47:30, 12:13:05.
\begin{enumerate}[(a)]
    \item Lesquelles évoquent une balise, et de quelle période~?
    \item Proposez une explication légitime de ce rythme.
    \item Quelles vérifications feriez-vous avant de conclure~?
\end{enumerate}
\end{boxexo}

\begin{boxexo}[title={\ding{45}\quad Exercice~6.4. Le NON des virements \niveauexpert}]
\begin{enumerate}[(a)]
    \item Expliquez pourquoi la seule absence des virements dans la capture
          ne suffirait pas à coter \NON{} la proposition « les virements
          sont passés par le lien du siège ».
    \item Quelle pièce indépendante permet ce \NON{}~? Quelle est sa propre
          faiblesse (Chapitre~\ref{chap:preuve})~?
    \item Proposez une démarche pour renforcer la cote, en vous plaçant du
          côté de l'enquête judiciaire.
\end{enumerate}
\end{boxexo}

\bigskip

\begin{boxinfo}
\textbf{Pour aller plus loin.}
\begin{itemize}
    \item R.~Bejtlich, \textit{The Practice of Network Security
          Monitoring}~\cite{bejtlich2013nsm}.
    \item C.~Sanders, \textit{Practical Packet Analysis}~\cite{sanders2017packet}.
\end{itemize}
\end{boxinfo}

\begin{boxinfo}
\textbf{Transition.} Le disque, la mémoire et le réseau ont chacun parlé,
chacun avec ses heures. Il reste à les mettre dans un même ordre, sur une
même échelle de temps~: Chapitre~\ref{chap:chronologie}.
\end{boxinfo}

        % ============================================================================
% Fichier  : partie2/P02_C07_chronologie_correlation.tex
% Titre    : Chronologie et Corrélation
% Auteur   : MINKA MI NGUIDJOÏ Thierry Emmanuel
% Version  : 1.0 -- novembre 2024
% Corrigés : corriges/Corr_C07.tex
% ============================================================================
% Sources : MITRE ATT&CK (tactiques) ; chaîne d'intrusion de Lockheed Martin
% (2011) ; outils de super-chronologie (log2timeline/Plaso, présentés sans
% citation).
% Faits nouveaux (bible) : temps de présence 11 j 7 h 42 min (02/03 23 h 48
% -> 14/03 7 h 30) ; trou du 3 au 8 mars ; dérive SRV-AD01 +2 s, passerelle
% -47 s (horloge non synchronisée).
% ============================================================================

\chapter{Chronologie et Corrélation}
\label{chap:chronologie}
\chaptermark{Chronologie et corrélation}

\epigraph{%
    « Une chronologie n'est pas une liste d'événements.\\
    C'est une liste d'événements dont on a vérifié les horloges. »%
}{Principe du cours}

\niveauChapitre{Intermédiaire}{%
    Maîtrise des Chapitres~\ref{chap:disque} à~\ref{chap:reseau}
}

\begin{boxobjectifs}
\begin{itemize}
    \item Construire une chronologie unifiée à partir de sources aux
          horloges et aux formats différents.
    \item Normaliser les heures et corriger une dérive mesurée.
    \item Corréler des événements par des pivots~: adresse, compte,
          empreinte, heure.
    \item Repérer les trous d'une chronologie et les coter.
    \item Calculer un temps de présence et le rattacher aux tactiques
          d'une attaque.
\end{itemize}
\end{boxobjectifs}

\bigskip

% ============================================================================
\section*{Avant de commencer : quarante-sept secondes}
% ============================================================================

Yaoundé, 26 mars 2025. M.~FOTSO assemble les événements établis par le
disque, la mémoire et le réseau. Une incohérence apparaît aussitôt. Selon le
journal de la passerelle, la session de \texttt{svc\_backup} qui précède
l'installation du service sur SRV-AD01, le 9~mars, s'ouvre à 1~h~35~min~41~s.
Selon le journal système de SRV-AD01, le service est créé à
1~h~35~min~09~s. L'attaquant aurait installé son outil trente-deux secondes
avant d'ouvrir sa session.

Mme~MBARGA ne conclut pas à une falsification. Elle demande la dérive de
chaque horloge. Les serveurs du siège se synchronisaient sur le serveur de
temps interne~; la passerelle, installée par un prestataire, ne le faisait
pas. Mesurée le 14~mars, son horloge avançait de 47~secondes.

\begin{boxalert}
\textbf{Violation typique.} Fusionner les heures de plusieurs sources sans
vérifier leurs fuseaux ni leurs dérives, puis tirer des conclusions de
l'ordre obtenu.

\smallskip
\textbf{Règle~:} chaque source reçoit, avant toute fusion, son fuseau, sa
dérive mesurée et sa précision. La chronologie se construit en UTC
corrigé~; un ordre entre deux événements plus proches que l'incertitude des
horloges n'est pas établi.
\end{boxalert}

\separateur

% ============================================================================
\section{Normaliser les Heures}
\label{sec:c07-normaliser}
% ============================================================================

\begin{tableaucours}{L{3.6cm}cccL{2.6cm}}{%
    \tableauHeader{Source} &
    \tableauHeader{Fuseau} &
    \tableauHeader{Dérive} &
    \tableauHeader{Précision} &
    \tableauHeader{Synchronisée}
}
Passerelle d'accès & UTC+1 & $+47$~s & 1~s & Non \\
SRV-AD01 (journaux) & UTC+1 & $+2$~s & 1~s & Oui \\
SRV-CPT01 (journaux, NTFS) & UTC+1, UTC & $+2$~s & 1~s, 100~ns & Oui \\
Sonde réseau & UTC & $0$~s & 1~µs & Oui (source publique) \\
Interface de paiement & UTC & Inconnue & 1~s & \NEK{} \\
\end{tableaucours}
\vspace{-0.8em}
\begin{center}{\small\itshape Tableau~7.1. Les horloges de l'affaire NYANGA}\end{center}

La dérive se mesure au moment de l'acquisition, en comparant l'heure de la
machine à une référence, et se suppose constante sur la période, ce qui est
une hypothèse à écrire. Corrigée de 47~secondes, l'ouverture de session sur
la passerelle passe à 1~h~34~min~54~s, quinze secondes avant la création du
service~: l'ordre redevient cohérent.

\begin{boxinfo}
\textbf{Préparer une source avant fusion}
\begin{enumerate}
    \item Identifier le fuseau dans lequel la source écrit ses heures.
    \item Mesurer ou rechercher la dérive de son horloge.
    \item Noter sa précision (seconde, milliseconde).
    \item Convertir en UTC et corriger la dérive, en gardant l'heure
          d'origine à côté.
    \item Noter les hypothèses (dérive constante, fuseau supposé).
\end{enumerate}
\end{boxinfo}

% ============================================================================
\section{Fusionner et Corréler}
\label{sec:c07-correler}
% ============================================================================

Des outils de super-chronologie, comme log2timeline et Plaso, extraient
automatiquement les horodatages de centaines de formats (journaux, registre,
système de fichiers, navigateurs) et les fusionnent dans une seule liste.
Une telle liste compte vite des millions de lignes. La corrélation consiste
à y relier des événements par un élément commun, appelé \termecle{pivot}.

\begin{tableaucours}{L{2.8cm}L{5.0cm}L{4.4cm}}{%
    \tableauHeader{Pivot} &
    \tableauHeader{Exemple à NYANGA} &
    \tableauHeader{Ce qu'il relie}
}
Adresse IP & 203.0.113.47 & Balise de SRV-AD01, exfiltration de SRV-FIC01 \\
Compte & \texttt{adm.kamga} & Sessions sur SRV-CPT01 et SRV-FIC01 depuis
SRV-AD01 \\
Empreinte de fichier & Binaire \texttt{WinUpd\textbackslash svchost.exe} &
Même fichier sur SRV-AD01 et SRV-FIC01 \\
Heure & 14~mars, 1~h~12 & Début de l'exfiltration, échanges volumineux de la
balise \\
Nom & \texttt{cdn-update.example.net} & Requêtes DNS, échanges TLS \\
\end{tableaucours}
\vspace{-0.8em}
\begin{center}{\small\itshape Tableau~7.2. Les pivots de corrélation}\end{center}

Un pivot relie, il ne prouve pas à lui seul un lien de cause. Deux machines
qui parlent à la même adresse peuvent obéir au même attaquant, ou utiliser
le même service légitime. Le lien se confirme par une seconde coïncidence
indépendante~: même adresse \emph{et} même binaire, même compte \emph{et}
même origine.

% ============================================================================
\section{La Chronologie de NYANGA}
\label{sec:c07-nyanga}
% ============================================================================

\begin{boxscenario}{Opération NYANGA. Du correctif au virement}
\begin{tableaucours}{L{2.8cm}L{6.0cm}L{2.2cm}c}{%
    \tableauHeader{Date et heure} &
    \tableauHeader{Événement} &
    \tableauHeader{Source} &
    \tableauHeader{Cote}
}
12/12/2024 & Bulletin de l'éditeur sur VULN-2024-4471 & Éditeur & \OUI{} \\
15/12/2024 & Correctif publié & Éditeur & \OUI{} \\
19/12/2024 & Inscription au catalogue KEV & CISA & \OUI{} \\
02/03/2025, 23~h~48 & Première session \texttt{svc\_backup} depuis
203.0.113.18 & Passerelle (corrigée) & \OUI{} \\
03 au 08/03 & Activité de l'attaquant & Aucune & \NEK{} \\
09/03, 1~h~34 & Session \texttt{svc\_backup} vers SRV-AD01 & Passerelle,
SRV-AD01 & \OUI{} \\
09/03, 1~h~35 & Service \texttt{WinUpdSvc} installé & SRV-AD01 (7045) &
\OUI{} \\
09/03, 1~h~37 & \texttt{lsass.dmp} créé & SRV-AD01 (\$UsnJrnl) & \OUI{} \\
10/03 et suivants & Sessions \texttt{adm.kamga} depuis SRV-AD01 & Serveurs
(4624) & \OUI{} \\
12/03, 22~h~39 à 23~h~06 & Ouverture de \texttt{api.ini}, archive,
suppression & SRV-CPT01 & \OUI{} \\
14/03, 1~h~12 à 3~h~47 & 3,2~Go de SRV-FIC01 vers 203.0.113.47 & Sonde,
NC-07 & \OUI{} \\
14/03, 2~h~05 à 3~h~31 & 23 virements depuis 203.0.113.62 & Interface &
\OUI{} \\
14/03, 7~h~30 & Détection par le service des paiements & Témoignage &
\OUI{} \\
14/03, 10~h~02 & Capture mémoire de SRV-AD01 & Registre des actions &
\OUI{} \\
\end{tableaucours}

\medskip
\textbf{Constatations.} Entre la première session et la détection
s'écoulent 11~jours, 7~heures et 42~minutes. Aucune source ne couvre
l'activité de l'attaquant du 3 au 8~mars~: la passerelle n'enregistre aucune
session de \texttt{svc\_backup}, et aucun serveur du siège ne garde trace de
lui.

\textbf{Analyse} (confiance moyenne). Le trou peut signifier que l'attaquant
a attendu, ou qu'il a agi par un autre chemin que les sources ne voient pas
(une agence, un autre compte). La seconde explication ne peut être écartée.
\end{boxscenario}

\begin{boxdef}
\textbf{\termecle{Temps de présence}} (\textit{dwell time})~: durée entre
la première action établie de l'attaquant dans le système et la détection
de l'incident. C'est une borne inférieure~: une action plus ancienne, non
détectée, l'allongerait.
\end{boxdef}

% ============================================================================
\section{Lire la Chronologie par Tactiques}
\label{sec:c07-tactiques}
% ============================================================================

La chaîne d'intrusion décrite par Lockheed Martin~\cite{killchain2011} et la
matrice ATT\&CK de MITRE~\cite{mitreattack} rangent les actions d'un
attaquant en étapes, ou tactiques. Placer chaque événement de la chronologie
dans une tactique permet de voir ce qui manque.

\begin{tableaucours}{L{3.6cm}L{5.4cm}c}{%
    \tableauHeader{Tactique (ATT\&CK)} &
    \tableauHeader{Événement de NYANGA} &
    \tableauHeader{Cote}
}
Accès initial & Session du 2~mars par la passerelle & \OUI{} \\
Persistance & Service \texttt{WinUpdSvc} & \OUI{} \\
Accès aux identifiants & Copie de la mémoire de \texttt{lsass.exe} & \OUI{} \\
Mouvement latéral & Sessions \texttt{adm.kamga} depuis SRV-AD01 & \OUI{} \\
Collecte & Archive \texttt{cpt.7z}, dossier \texttt{Clients} & \OUI{} \\
Commande et contrôle & Balise de 60~secondes vers 203.0.113.47 & \OUI{} \\
Exfiltration & 3,2~Go le 14~mars & \OUI{} \\
Impact & 23 virements frauduleux & \OUI{} \\
Découverte (reconnaissance interne) & Aucun événement retrouvé & \NEK{} \\
\end{tableaucours}
\vspace{-0.8em}
\begin{center}{\small\itshape Tableau~7.3. La chronologie de NYANGA par tactiques}\end{center}

La dernière ligne est instructive. Un attaquant qui trouve le fichier de clé
d'une application de paiement a forcément exploré le réseau avant. Aucune
trace de cette exploration n'a été trouvée~: elle s'est probablement
déroulée pendant le trou du 3 au 8~mars, avec des moyens que les sources du
siège n'enregistrent pas. La tactique est presque certaine, l'événement
reste \NEK{}.

\begin{boxlecture}{La chronologie au 26 mars}
\textbf{Décompte.} Sur les quatorze lignes de la chronologie, treize sont
cotées \OUI{} et une \NEK{}~; sur les neuf tactiques, huit \OUI{} et une
\NEK{}.

\medskip
\textbf{Lecture.} La chronologie est solide sur ce qu'elle montre, et ses
deux \NEK{} portent sur la même période~: du 3 au 8~mars. C'est là que
l'enquête doit chercher, en élargissant ses sources aux agences et aux
autres comptes, avant de conclure que le temps de présence est de onze
jours.
\end{boxlecture}

% ============================================================================
\section{Ce que la Chronologie demande à chaque Métier}
\label{sec:c07-metiers}
% ============================================================================

\begin{boxmetier}{Analyste en réponse à incident}
\textbf{Construire une chronologie}
\begin{enumerate}
    \item Dresser le tableau des horloges~: fuseau, dérive, précision.
    \item Normaliser en UTC corrigé, en gardant l'heure d'origine.
    \item Fusionner, puis filtrer par pivots.
    \item Confirmer chaque lien par deux coïncidences indépendantes.
    \item Coter chaque ligne et nommer les trous.
\end{enumerate}
\end{boxmetier}

\begin{boxmetier}{Application OPJ/Expert Judiciaire}
Devant un tribunal, une chronologie se présente avec ses incertitudes~: on
écrit « entre 1~h~34 et 1~h~35, à une minute près » plutôt qu'une seconde
précise qu'aucune horloge ne garantit.
\end{boxmetier}

\begin{boxmetier}{RSSI}
Synchroniser toutes les horloges, y compris celles des équipements installés
par des prestataires. La passerelle de NYANGA, seule hors synchronisation,
était aussi le point d'entrée.
\end{boxmetier}

\separateur

\begin{boxresume}
\textbf{Ce qu'il faut retenir de ce chapitre~:}
\begin{itemize}
    \item Avant toute fusion, chaque source reçoit son fuseau, sa dérive et
          sa précision~; la chronologie se construit en UTC corrigé.
    \item Un pivot relie des événements~; un lien de cause demande deux
          coïncidences indépendantes.
    \item Le temps de présence est une borne inférieure.
    \item Ranger la chronologie par tactiques montre ce qui manque.
    \item À NYANGA, le temps de présence établi est de 11~jours, 7~heures
          et 42~minutes~; la période du 3 au 8~mars reste \NEK{}.
\end{itemize}
\end{boxresume}

\begin{boxinfo}
\textbf{Points de contrôle.} Vous avez compris ce chapitre si vous savez~:
\begin{itemize}
    \item corriger une heure d'une dérive mesurée~;
    \item calculer un temps de présence~;
    \item ranger des événements par tactiques et repérer une tactique sans
          événement.
\end{itemize}
\end{boxinfo}

% ============================================================================
\section*{Exercices}
% ============================================================================

\begin{boxexo}[title={\ding{45}\quad Exercice~7.1. Corriger une heure \niveaubase}]
La passerelle avance de 47~secondes. Corrigez et convertissez en UTC les
heures suivantes, lues dans son journal en heure locale~: (a) 02/03/2025
23:48:35~; (b) 09/03/2025 01:35:41~; (c) 14/03/2025 00:58:12.
\end{boxexo}

\begin{boxexo}[title={\ding{45}\quad Exercice~7.2. Le temps de présence \niveaumoyen}]
\begin{enumerate}[(a)]
    \item Vérifiez le calcul du temps de présence de NYANGA.
    \item On découvre dans les journaux d'une agence une connexion de
          \texttt{svc\_backup} le 24~février à 21~h~10. Que devient le temps
          de présence~? Quelle vérification faites-vous d'abord~?
    \item Pourquoi parle-t-on de borne inférieure~?
\end{enumerate}
\end{boxexo}

\begin{boxexo}[title={\ding{45}\quad Exercice~7.3. Un pivot trompeur \niveaumoyen}]
Un poste de la direction a parlé à 203.0.113.47 le 11~mars, à 9~h~00, pendant
quatre secondes.
\begin{enumerate}[(a)]
    \item Ce poste est-il compromis~?
    \item Quelles secondes coïncidences chercheriez-vous~?
    \item Comment cotez-vous la proposition « le poste de la direction est
          compromis » avant et après ces vérifications~?
\end{enumerate}
\end{boxexo}

\begin{boxexo}[title={\ding{45}\quad Exercice~7.4. Combler le trou \niveauexpert}]
Proposez un plan d'investigation pour la période du 3 au 8~mars~: sources à
recueillir, questions, pivots, et cotes attendues selon ce que l'on
trouvera. Indiquez ce qui ferait passer la tactique « découverte » de \NEK{}
à \OUI{}, et ce qui ne permettra jamais de la passer à \NON{}.
\end{boxexo}

\bigskip

\begin{boxinfo}
\textbf{Pour aller plus loin.}
\begin{itemize}
    \item La matrice MITRE ATT\&CK~\cite{mitreattack}.
    \item E.~Hutchins, M.~Cloppert, R.~Amin, sur la chaîne
          d'intrusion~\cite{killchain2011}.
\end{itemize}
\end{boxinfo}

\begin{boxinfo}
\textbf{Transition.} La partie~II a montré comment acquérir et analyser.
Mais qui a le droit de le faire, selon quelles normes, et sous quelles
règles de droit~? C'est l'objet du Chapitre~\ref{chap:normes-droit}.
\end{boxinfo}


    % PARTIE 3 · CADRE NORMATIF ET JURIDIQUE (ISO/IEC 27037 à 27043 ; loi
    % n° 2010/012 ; Budapest ; Malabo)
    \part{Cadre Normatif et Juridique}
    \label{part:normatif}

        % ============================================================================
% Fichier  : partie3/P03_C08_normes_droit.tex
% Titre    : Normes et Droit de l'Investigation
% Auteur   : MINKA MI NGUIDJOÏ Thierry Emmanuel
% Version  : 1.0 -- novembre 2026
% Corrigés : corriges/Corr_C08.tex
% ============================================================================
% Textes cités dans leur version officielle (Normes/loi_2010-012, lus le
% 08/11/2026) : loi n° 2010/012, art. 26 al. 4, 29 al. 1 à 3, 30 al. 1,
% 52, 55, 56, 65 al. 2, 69, 72, 91 al. 3. Normes ISO/IEC paraphrasées.
% À vérifier par l'auteur : statut du Cameroun au regard de la Convention de
% Budapest et de la Convention de Malabo ; dispositions du Code de procédure
% pénale sur l'expertise.
% ============================================================================

\chapter{Normes et Droit de l'Investigation}
\label{chap:normes-droit}
\chaptermark{Normes et droit}

\epigraph{%
    « La compétence technique dit ce que l'on peut faire.\\
    Le droit dit ce que l'on a le droit de faire. »%
}{Principe du cours}

\niveauChapitre{Fondamental}{%
    Lecture des Chapitres~\ref{chap:investigation} et~\ref{chap:preuve}
}

\begin{boxobjectifs}
\begin{itemize}
    \item Situer les normes de l'investigation numérique (famille ISO/IEC
          27037 à 27043, ISO/IEC 27035, ISO/IEC 17025, NIST) et dire ce que
          chacune encadre.
    \item Citer les articles de la loi n°~2010/012 qui organisent l'enquête
          (articles~52 à~58).
    \item Qualifier des faits au regard des infractions de la loi
          n°~2010/012.
    \item Identifier les obligations de l'exploitant d'un système
          d'information qui pèsent sur la préparation d'une enquête.
    \item Distinguer les pouvoirs de l'OPJ, de l'expert requis et de
          l'enquêteur privé.
    \item Décrire la voie de l'entraide judiciaire internationale.
\end{itemize}
\end{boxobjectifs}

\bigskip

% ============================================================================
\section*{Avant de commencer : soixante jours au lieu de dix ans}
% ============================================================================

Bafoussam, 30 mars 2025. Me~ETOUNDI, la juriste de NYANGA CRÉDIT, relit avec
Mme~NDOUMBE la loi n°~2010/012 relative à la cybersécurité et à la
cybercriminalité au Cameroun. Elle cherche les infractions que la plainte
peut viser. Elle trouve aussi un article qu'elle n'attendait pas~: la loi
oblige les exploitants de systèmes d'information à conserver les données de
connexion et de trafic pendant dix ans. Les journaux de la passerelle de
l'établissement n'en gardaient que soixante jours.

\begin{boxalert}
\textbf{Violation typique.} Traiter le droit comme une question qui ne se
pose qu'au moment de porter plainte.

\smallskip
\textbf{Règle~:} le droit encadre l'enquête du premier geste au dernier~:
qui peut recueillir quoi, comment saisir, comment requérir un tiers, et ce
que l'organisation devait avoir préparé. L'enquêteur qui l'ignore produit
des pièces fragiles et expose son client.
\end{boxalert}

\separateur

% ============================================================================
\section{Les Normes de l'Investigation}
\label{sec:c08-normes}
% ============================================================================

Les normes ne disent pas le droit~; elles décrivent l'état de l'art. Un juge
ou un expert adverse qui conteste une méthode demandera si elle suit une
pratique reconnue. Les normes fournissent cette référence.

\begin{tableaucours}{L{2.8cm}L{5.4cm}L{4.0cm}}{%
    \tableauHeader{Texte} &
    \tableauHeader{Objet (paraphrasé)} &
    \tableauHeader{Chapitres du cours}
}
ISO/IEC 27037 & Identification, recueil, acquisition et préservation de la
preuve numérique & \ref{chap:preuve}, \ref{chap:acquisition} \\
ISO/IEC 27041 & Assurance que les méthodes d'examen sont adaptées et
validées & \ref{chap:disque} à \ref{chap:reseau} \\
ISO/IEC 27042 & Analyse et interprétation de la preuve numérique &
\ref{chap:disque} à \ref{chap:chronologie} \\
ISO/IEC 27043 & Principes et processus de l'investigation, de la
préparation à la clôture & \ref{chap:investigation} \\
ISO/IEC 27035 & Gestion des incidents de sécurité de l'information &
\ref{chap:investigation}, \ref{chap:enquete} \\
ISO/IEC 17025 & Compétence des laboratoires d'essais et d'étalonnage, base
de l'accréditation des laboratoires forensiques & \ref{chap:rapport} \\
NIST SP 800-86 & Intégration de la forensique à la réponse aux incidents &
\ref{chap:investigation} \\
NIST SP 800-61 & Gestion des incidents de sécurité & \ref{chap:investigation} \\
\end{tableaucours}
\vspace{-0.8em}
\begin{center}{\small\itshape Tableau~8.1. Les normes de l'investigation numérique}\end{center}

\begin{boxinfo}
\textbf{ISO/IEC 27041 et la validation des outils.} La norme
ISO/IEC~27041~\cite{iso27041_2015} demande que l'enquêteur puisse montrer
qu'une méthode fait ce qu'il prétend. En pratique, on teste un outil sur une
image de référence dont le contenu est connu, et l'on garde la trace du
test. Un outil non validé n'est pas interdit, mais son résultat porte une
confiance moindre, et l'enquêteur doit le dire.
\end{boxinfo}

% ============================================================================
\section{Le Cadre Camerounais de l'Enquête}
\label{sec:c08-enquete}
% ============================================================================

La loi n°~2010/012 du 21~décembre 2010 relative à la cybersécurité et à la
cybercriminalité~\cite{loi2010012} contient un chapitre consacré à la
procédure. Ses articles~52 à~58 organisent l'enquête.

\begin{boxjuris}
\textbf{\loicam{n°~2010/012}, article~52}~\cite{loi2010012}~:
\textit{« (1) En cas d'infraction cybernétique, les Officiers de Police
Judiciaire à compétence générale et les agents habilités de l'Agence,
procèdent aux enquêtes conformément aux dispositions du Code de Procédure
Pénale. [\ldots] (3) Les Officiers de Police Judiciaire et les agents
habilités de l'Agence peuvent, lors des investigations, accéder aux moyens
de transport, à tout local à usage professionnel, à l'exclusion des
domiciles privés, en vue de rechercher, de constater les infractions, de
demander la communication de tous les documents professionnels et en
prendre copie, recueillir, sur convocation ou sur place, les renseignements
et justifications. »}

\medskip
\textbf{Portée} (analyse). L'enquête pénale appartient aux OPJ et aux agents
habilités de l'Agence nationale des technologies de l'information et de la
communication (ANTIC). Leurs pouvoirs d'accès s'arrêtent au domicile privé,
qui relève des règles générales de la perquisition.
\end{boxjuris}

\begin{tableaucours}{cL{9.4cm}}{%
    \tableauHeader{Article} &
    \tableauHeader{Objet}
}
52 & Enquêteurs compétents (OPJ, agents habilités de l'ANTIC) et pouvoirs
d'accès \\
53 & Perquisition portant sur des données, copie en présence des
personnes, scellés (Chapitre~\ref{chap:preuve}) \\
54 & Perquisitions et saisies selon le Code de procédure pénale, en tenant
compte du dépérissement des preuves \\
55 & Réquisition d'une personne qualifiée pour obtenir la version en clair
de données transformées~; remise de la convention de déchiffrement \\
56 & Réquisition à expert, selon les règles de la commission d'expert \\
57 & Commissions rogatoires nationales et internationales \\
58 & Obligations des prestataires de cryptographie envers les enquêteurs \\
\end{tableaucours}
\vspace{-0.8em}
\begin{center}{\small\itshape Tableau~8.2. Les articles de procédure de la loi n°~2010/012}\end{center}

\begin{boxjuris}
\textbf{\loicam{n°~2010/012}, article~55}~\cite{loi2010012}~:
\textit{« (1) Lorsqu'il apparaît que les données saisies ou obtenues au
cours de l'enquête ou de l'instruction ont fait l'objet d'opérations de
transformation empêchant d'accéder en clair ou sont de nature à
compromettre les informations qu'elles contiennent, le Procureur de la
République, le Juge d'Instruction ou la juridiction de jugement peuvent
réquisitionner toute personne physique ou morale qualifiée, en vue
d'effectuer les opérations techniques permettant d'obtenir la version en
clair desdites données. (2) Lorsqu'un moyen de cryptographie a été utilisé,
les autorités judiciaires peuvent exiger la convention secrète de
déchiffrement du cryptogramme. »}

\medskip
\textbf{\loicam{n°~2010/012}, article~56}~\cite{loi2010012}~:
\textit{« La réquisition prévue à l'article 50 ci-dessus peut être faite à
tout expert. Dans ce cas, son exécution est faite conformément aux
dispositions du Code de procédure pénale relatives à la commission
d'expert. »}

\medskip
\textbf{Portée} (analyse). L'article~55 fonde la réquisition qui permettrait
d'ouvrir l'archive chiffrée \texttt{cpt.7z} (Chapitre~\ref{chap:disque}).
L'article~56 fonde la désignation de M.~ABOUBAKAR~: l'expert requis agit
selon les règles de l'expertise du Code de procédure pénale [dispositions à
vérifier par le lecteur sur le texte officiel], et non selon un mandat de
l'établissement.
\end{boxjuris}

% ============================================================================
\section{Qualifier les Faits}
\label{sec:c08-infractions}
% ============================================================================

L'enquête technique établit des faits~; la qualification les rattache à des
infractions. Elle appartient au procureur et au juge, mais l'enquêteur doit
savoir quels faits chaque infraction exige, pour ne pas oublier d'établir
l'un d'eux.

\begin{boxjuris}
\textbf{\loicam{n°~2010/012}, article~65, alinéa~2}~\cite{loi2010012}~:
\textit{« Est puni des peines prévues à l'alinéa 1 ci-dessus, tout accès non
autorisé, à l'ensemble ou à une partie d'un réseau de communications
électroniques ou d'un système d'information ou d'un équipement
terminal. »}

\medskip
\textbf{\loicam{n°~2010/012}, article~72}~\cite{loi2010012}~:
\textit{« Est puni des peines prévues par l'article 66 ci-dessus celui qui,
de quelque manière que ce soit, sans droit, introduit, altère, efface, ou
supprime, afin d'obtenir un bénéfice économique, les données électroniques,
de manière à causer un préjudice patrimonial à autrui. »}
\end{boxjuris}

\begin{boxscenario}{Opération NYANGA. Les faits et les textes}
\begin{tableaucours}{L{4.6cm}L{2.4cm}L{5.0cm}}{%
    \tableauHeader{Fait établi} &
    \tableauHeader{Texte possible} &
    \tableauHeader{Élément que l'enquête doit établir}
}
Session du 2~mars sur la passerelle, par contournement de
l'authentification & Art.~65 al.~2, art.~69 & Accès sans droit~;
violation des mesures de sécurité (art.~69) \\
Copie du dossier \texttt{Clients} & Art.~69 & Accès « afin d'obtenir des
informations ou des données » \\
Installation d'un service, copie de la mémoire de \texttt{lsass.exe} &
Art.~71 & Introduction de données en vue de modifier celles du système \\
23 virements avec la clé volée & Art.~72 & Introduction de données sans
droit, bénéfice économique, préjudice patrimonial \\
\end{tableaucours}

\medskip
\textbf{Analyse} (confiance moyenne, qualification réservée au parquet).
Les faits établis aux Chapitres~\ref{chap:disque} à~\ref{chap:chronologie}
couvrent les éléments matériels de ces infractions. L'identité des auteurs,
\NEK{}, conditionne les poursuites, pas la qualification des faits.
\end{boxscenario}

\begin{boxjuris}
\textbf{\loicam{n°~2010/012}, article~69}~\cite{loi2010012}~:
\textit{« Est puni d'un emprisonnement de cinq (05) à dix (10) ans et d'une
amende de 10.000.000 (dix millions) à 100.000.000 (cent millions) F CFA ou
de l'une de ces deux peines seulement, celui qui accède sans droit, et en
violation des mesures de sécurité, à l'ensemble ou à une partie d'un réseau
de communications électroniques, d'un système d'information ou d'un
équipement terminal, afin d'obtenir des informations ou des données, en
relation avec un système d'information connecté à un autre système
d'information. »}

\medskip
\textbf{Portée} (analyse). L'article~69 exige une violation des mesures de
sécurité. Le contournement de l'authentification de la passerelle en est
une~; l'enquête doit l'établir techniquement, ce que fera le
Chapitre~\ref{chap:artefact-vuln}. Le même article interdit à l'enquêteur
privé d'entrer sans autorisation dans un système, fût-ce pour suivre
l'attaquant~: remonter jusqu'au serveur 203.0.113.47 et s'y connecter
serait un accès sans droit.
\end{boxjuris}

% ============================================================================
\section{Les Obligations de l'Exploitant}
\label{sec:c08-obligations}
% ============================================================================

La loi pèse aussi sur la victime, au titre de ses obligations d'exploitant de
système d'information. Trois dispositions intéressent directement la
préparation d'une enquête.

\begin{boxjuris}
\textbf{\loicam{n°~2010/012}, article~29}~\cite{loi2010012}~:
\textit{« (1) Les exploitants des systèmes d'information ont l'obligation
de conserver les données de connexion et de trafic de leurs systèmes
d'information pendant une période de dix (10) ans. (2) Les exploitants des
systèmes d'information sont tenus d'installer des mécanismes de
surveillance de contrôle d'accès aux données de leurs systèmes
d'information. Les données conservées peuvent être accessibles lors des
investigations judiciaires. (3) Les installations des exploitants des
systèmes d'information peuvent faire l'objet de perquisition ou de saisie
sur ordre d'une autorité judiciaire dans les conditions prévues par les lois
et règlements en vigueur. »}

\medskip
\textbf{\loicam{n°~2010/012}, article~26, alinéa~4}~\cite{loi2010012}~:
\textit{« Les plates-formes des systèmes d'information font l'objet de
protection contre d'éventuels rayonnements et des intrusions qui pourraient
compromettre l'intégrité des données transmises et contre toute autre
attaque externe notamment par un système de détection d'intrusions. »}

\medskip
\textbf{\loicam{n°~2010/012}, article~30, alinéa~1}~\cite{loi2010012}~:
\textit{« Les exploitants des systèmes d'information évaluent, révisent leurs
systèmes de sécurité et introduisent en cas de nécessité les modifications
appropriées dans leurs pratiques, mesures et techniques de sécurité en
fonction de l'évolution des technologies. »}
\end{boxjuris}

\begin{boxlecture}{NYANGA CRÉDIT face à ses obligations d'exploitant}
\begin{tableaucours}{L{6.6cm}cL{3.6cm}}{%
    \tableauHeader{Obligation} &
    \tableauHeader{Cote} &
    \tableauHeader{Pièce}
}
Conserver dix ans les données de connexion et de trafic (art.~29 al.~1) &
\NON{} & Passerelle~: 60~jours \\
Surveiller le contrôle d'accès aux données (art.~29 al.~2) & \OUI{} &
Journaux d'authentification du domaine \\
Disposer d'un système de détection d'intrusions (art.~26 al.~4) & \OUI{} &
Sonde installée en 2024 \\
Réviser les mesures selon l'évolution des technologies (art.~30 al.~1) &
\NON{} & Correctif non appliqué en 89~jours \\
Les mécanismes ont reçu le visa de l'Agence (art.~26 al.~3) & \NEK{} &
Aucun document produit \\
\end{tableaucours}

\medskip
\textbf{Lecture} (analyse, confiance moyenne). Deux \OUI{}, deux \NON{}, un
\NEK{}. Les deux \NON{} pèsent sur la victime~: la durée de conservation a
failli effacer le premier accès, et l'absence de correction a ouvert la
porte. Un avocat de la défense, ou un client lésé par la fuite, pourra s'en
servir. Le rapport d'enquête ne les cache pas~; il les sépare des faits
reprochés aux attaquants.
\end{boxlecture}

% ============================================================================
\section{Trois Statuts, Trois Pouvoirs}
\label{sec:c08-statuts}
% ============================================================================

\begin{tableaucours}{L{2.6cm}L{3.2cm}L{3.2cm}L{3.2cm}}{%
    \tableauHeader{} &
    \tableauHeader{OPJ, agent de l'ANTIC} &
    \tableauHeader{Expert requis} &
    \tableauHeader{Enquêteur privé}
}
Fondement & Art.~52 & Art.~56, Code de procédure pénale & Contrat avec
l'organisation \\
Mandant & Procureur, juge & L'autorité qui le requiert & Le client \\
Accès aux locaux & Locaux professionnels (art.~52 al.~3) & Selon sa
mission & Ceux du client seulement \\
Saisie et scellés & Oui (art.~53) & Non, il reçoit des scellés & Non~;
il conserve sous chaîne de possession \\
Requérir un tiers & Oui, par réquisition & Non & Non~; il demande, sans
contrainte \\
Destinataire du rapport & Le parquet & L'autorité requérante & Le client \\
\end{tableaucours}
\vspace{-0.8em}
\begin{center}{\small\itshape Tableau~8.3. Les pouvoirs selon le statut}\end{center}

L'enquêteur privé n'a aucun pouvoir de contrainte. Il travaille sur les
systèmes que son client possède ou dont il peut autoriser l'examen, et il
remet ses pièces à l'OPJ quand une plainte est déposée. C'est pourquoi la
qualité de ses premières acquisitions compte autant~: elles deviendront les
pièces de l'enquête pénale.

% ============================================================================
\section{L'Entraide Internationale}
\label{sec:c08-entraide}
% ============================================================================

Les adresses 203.0.113.18, .47 et .62 de l'affaire, si elles étaient
réelles, appartiendraient très probablement à des hébergeurs étrangers.
Obtenir leurs données suppose l'entraide judiciaire internationale.

\begin{boxjuris}
\textbf{\loicam{n°~2010/012}, article~91, alinéa~3}~\cite{loi2010012}~:
\textit{« En cas d'urgence, les demandes d'entraide demandées par les
autorités camerounaises ou étrangères peuvent être transmises directement
aux autorités de l'Etat requis pour leur exécution. Le renvoi des pièces
d'exécution aux autorités compétentes de l'Etat requérant est effectué selon
les mêmes modalités. »}

\medskip
\textbf{Portée} (analyse). La voie ordinaire passe par le ministère chargé
des relations extérieures (alinéa~1)~; l'urgence autorise la transmission
directe. Le dépérissement des données des hébergeurs justifie souvent
l'urgence.
\end{boxjuris}

Deux conventions internationales structurent la coopération. La Convention
du Conseil de l'Europe sur la cybercriminalité, dite de
Budapest~\cite{budapest2001}, organise la conservation rapide des données et
l'entraide entre ses parties. La Convention de l'Union africaine sur la
cybersécurité et la protection des données à caractère personnel, dite de
Malabo~\cite{ua2014malabo}, entrée en vigueur en 2023, fixe un cadre
africain. Le statut du Cameroun au regard de chacune [à vérifier par le
lecteur sur les sites des deux organisations] commande la voie utilisable~:
sans convention applicable, l'entraide repose sur la réciprocité et les
accords bilatéraux.

% ============================================================================
\section{Ce que le Droit demande à chaque Métier}
\label{sec:c08-metiers}
% ============================================================================

\begin{boxmetier}{Analyste en réponse à incident}
Ne jamais suivre l'attaquant hors des systèmes du client~; demander par
écrit la conservation des données aux tiers dès le premier jour~; préparer
ses pièces comme si elles devaient être remises à un OPJ, ce qui arrive
souvent.
\end{boxmetier}

\begin{boxmetier}{Application OPJ/Expert Judiciaire}
L'OPJ requiert les tiers sans délai, en invoquant le dépérissement des
preuves (art.~54). L'expert s'en tient aux questions de sa mission et signale
à l'autorité tout fait découvert hors de celle-ci.
\end{boxmetier}

\begin{boxmetier}{RSSI}
Mettre la conservation des journaux en conformité avec l'article~29~: dix
ans pour les données de connexion et de trafic. La mesure se chiffre en
stockage, et elle conditionne toute enquête future.
\end{boxmetier}

\separateur

\begin{boxresume}
\textbf{Ce qu'il faut retenir de ce chapitre~:}
\begin{itemize}
    \item Les normes ISO/IEC~27037 à~27043 décrivent l'état de l'art de
          l'investigation~; l'ISO/IEC~27041 exige des méthodes validées.
    \item La loi n°~2010/012 confie l'enquête aux OPJ et aux agents de
          l'ANTIC (art.~52), encadre la saisie de données (art.~53, 54), le
          déchiffrement (art.~55) et l'expertise (art.~56).
    \item Les faits de NYANGA relèvent notamment des articles~65 al.~2,
          69, 71 et~72~; la qualification appartient au parquet.
    \item L'exploitant doit conserver dix ans ses données de connexion et
          de trafic (art.~29 al.~1)~: NYANGA ne les gardait que 60~jours.
    \item L'enquêteur privé n'a aucun pouvoir de contrainte~; l'entraide
          internationale passe par les autorités judiciaires.
\end{itemize}
\end{boxresume}

\begin{boxinfo}
\textbf{Points de contrôle.} Vous avez compris ce chapitre si vous savez~:
\begin{itemize}
    \item citer l'article qui fonde la saisie par copie de données~;
    \item rattacher un fait établi à une infraction et nommer l'élément à
          prouver~;
    \item dire ce qu'un enquêteur privé peut et ne peut pas faire.
\end{itemize}
\end{boxinfo}

% ============================================================================
\section*{Exercices}
% ============================================================================

\begin{boxexo}[title={\ding{45}\quad Exercice~8.1. Quel texte~? \niveaubase}]
Indiquez l'article de la loi n°~2010/012 qui répond à chaque question~:
(a) qui conduit l'enquête pénale en cas d'infraction cybernétique~?
(b) peut-on saisir une copie des données plutôt que le serveur~? (c) comment
obtenir la version en clair d'une archive chiffrée~? (d) combien de temps
l'exploitant doit-il conserver ses données de trafic~? (e) comment désigner
un expert~?
\end{boxexo}

\begin{boxexo}[title={\ding{45}\quad Exercice~8.2. Ce que l'enquêteur privé ne fait pas \niveaumoyen}]
Mme~MBARGA reçoit trois demandes de son client. Dites, pour chacune, si elle
peut y répondre et pourquoi~: (a) se connecter au serveur 203.0.113.47 avec
les identifiants trouvés dans l'outil de l'attaquant~; (b) examiner le
portable personnel de M.~KAMGA, qu'il apporte de lui-même~; (c) obtenir de
l'opérateur de paiement mobile l'identité des titulaires des onze
portefeuilles.
\end{boxexo}

\begin{boxexo}[title={\ding{45}\quad Exercice~8.3. Qualifier \niveaumoyen}]
Pour chacun des faits suivants, proposez un article de la loi n°~2010/012
et nommez l'élément que l'enquête doit établir~: (a) la création du compte
\texttt{svc\_backup} par l'attaquant, s'il était établi~; (b) la balise de
soixante secondes~; (c) l'effacement de l'archive \texttt{cpt.7z}~;
(d) l'envoi de 3,2~Go vers l'extérieur.
\end{boxexo}

\begin{boxexo}[title={\ding{45}\quad Exercice~8.4. Les obligations de la victime \niveauexpert}]
Un client dont le dossier figure parmi les fichiers exfiltrés assigne NYANGA
CRÉDIT en réparation.
\begin{enumerate}[(a)]
    \item Quels manquements aux articles~26, 29 et~30 peut-il invoquer, et
          sur quelles pièces de l'enquête~?
    \item NYANGA répond qu'elle est victime d'une infraction. Montrez que
          les deux positions ne s'excluent pas.
    \item Rédigez la recommandation du rapport d'enquête qui traite la
          conservation des journaux, en citant le texte.
\end{enumerate}
\end{boxexo}

\bigskip

\begin{boxinfo}
\textbf{Pour aller plus loin.}
\begin{itemize}
    \item La loi n°~2010/012~\cite{loi2010012}, dans son texte officiel.
    \item Les normes ISO/IEC~27037~\cite{iso27037_2012},
          27041~\cite{iso27041_2015}, 27042~\cite{iso27042_2015},
          27043~\cite{iso27043_2015}.
    \item La Convention de Budapest~\cite{budapest2001} et la Convention de
          Malabo~\cite{ua2014malabo}.
\end{itemize}
\end{boxinfo}

\begin{boxinfo}
\textbf{Transition.} L'enquête a montré par où l'attaquant est entré~: une
faille connue de la passerelle. La partie~IV change de point de vue~:
qu'est-ce qu'une vulnérabilité, et comment la mesurer~?
Chapitre~\ref{chap:vulnerabilite}.
\end{boxinfo}


    % PARTIE 4 · VULNÉRABILITÉS (CVE, CWE, CVSS v4.0, EPSS, KEV ; NIST SP 800-40)
    \part{Vulnérabilités : Identifier, Coter, Corriger}
    \label{part:vulnerabilites}

        % ============================================================================
% Fichier  : partie4/P04_C09_vulnerabilite.tex
% Titre    : Qu'est-ce qu'une Vulnérabilité
% Auteur   : MINKA MI NGUIDJOÏ Thierry Emmanuel
% Version  : 1.0 -- novembre 2026
% Corrigés : corriges/Corr_C09.tex
% ============================================================================
% Sources : programme CVE (MITRE), CWE, NVD (NIST) ; CVSS v4.0 (FIRST, 2023)
% et v3.1 (2019) ; EPSS (FIRST) ; catalogue KEV (CISA) ; SSVC (CISA, 2020).
% Les scores CVSS v4.0 donnés (9,3 et 10,0) correspondent aux vecteurs cités
% (calculateur du FIRST) ; les autres valeurs sont fictives et marquées
% comme telles. Identifiants VULN-2025-xxxx fictifs.
% Faits nouveaux (bible) : vecteur de VULN-2024-4471 ; quatre
% vulnérabilités d'exemple chez SOLITRANS (V-A à V-D).
% ============================================================================

\chapter{Qu'est-ce qu'une Vulnérabilité}
\label{chap:vulnerabilite}
\chaptermark{Qu'est-ce qu'une vulnérabilité}

\epigraph{%
    « Un score dit la gravité d'une faille dans l'absolu.\\
    Il ne dit pas qu'elle sera exploitée chez vous. »%
}{Principe du cours}

\niveauChapitre{Fondamental}{%
    Lecture du Chapitre~\ref{chap:prologue}
}

\begin{boxobjectifs}
\begin{itemize}
    \item Définir une vulnérabilité et la distinguer de la faiblesse, de la
          menace et du risque.
    \item Lire un identifiant CVE et le rattacher à une faiblesse CWE.
    \item Lire un vecteur CVSS v4.0 et en expliquer chaque métrique.
    \item Interpréter un score EPSS et une inscription au catalogue KEV.
    \item Expliquer pourquoi la gravité technique ne suffit pas à décider
          d'une priorité.
\end{itemize}
\end{boxobjectifs}

\bigskip

% ============================================================================
\section*{Avant de commencer : 412 critiques}
% ============================================================================

Douala, mai 2025. SOLITRANS LOGISTIQUE\footnote{SOLITRANS LOGISTIQUE, ses
salariés, ses sites et ses chiffres sont fictifs~; les référentiels cités
(CVE, CWE, CVSS, EPSS, KEV, SSVC) sont exacts.}, transitaire et
logisticien du port, vient de réaliser son premier balayage complet de
vulnérabilités sur ses 1\,380 actifs. Le rapport de l'outil compte 3\,812
constats, dont 412 qualifiés de « critiques » par leur score. Mme~EBODE,
responsable de la gestion des vulnérabilités depuis trois semaines, doit
dire au comité de direction ce qu'on corrige d'abord. Le RSSI, M.~NJIKE, a
une réponse simple~: « les 412 critiques, dans l'ordre du score ». Mme~EBODE
sait que l'équipe système corrige environ quarante failles par semaine. Il
faudrait dix semaines pour les seules critiques.

\begin{boxalert}
\textbf{Violation typique.} Confondre la gravité technique d'une
vulnérabilité, mesurée par un score, avec le risque qu'elle fait courir à
l'organisation.

\smallskip
\textbf{Règle~:} le score de base décrit la faille telle qu'elle est partout.
La priorité de correction dépend aussi de ce que l'on sait de son
exploitation et de la place de l'actif dans l'organisation.
\end{boxalert}

\separateur

% ============================================================================
\section{Faiblesse, Vulnérabilité, Menace, Risque}
\label{sec:c09-definitions}
% ============================================================================

\begin{boxdef}
\textbf{\termecle{Faiblesse}}~: type d'erreur de conception ou de
programmation susceptible de créer des vulnérabilités, par exemple une
authentification mal vérifiée ou une entrée non contrôlée. Le catalogue CWE
(\textit{Common Weakness Enumeration}) les recense~\cite{mitrecwe}.

\smallskip
\textbf{\termecle{Vulnérabilité}}~: défaut présent dans une version précise
d'un produit, qui peut être exploité pour porter atteinte à la
confidentialité, à l'intégrité ou à la disponibilité. Le programme CVE
(\textit{Common Vulnerabilities and Exposures}) lui attribue un identifiant
unique~\cite{mitrecve}.

\smallskip
\textbf{\termecle{Menace}}~: source potentielle d'atteinte, par exemple un
groupe qui exploite une famille de failles.

\smallskip
\textbf{\termecle{Risque}}~: combinaison de la vraisemblance qu'une menace
exploite une vulnérabilité et de la gravité de l'atteinte qui en résulte
pour l'organisation.
\end{boxdef}

La faille de la passerelle de NYANGA CRÉDIT illustre la chaîne. La
faiblesse est un contrôle d'authentification incomplet (CWE-287,
authentification incorrecte). La vulnérabilité est son instance dans le
micrologiciel 7.2.1 de la passerelle, désignée dans ce livre par
l'identifiant fictif VULN-2024-4471. La menace, ce sont les groupes qui
l'exploitaient dès décembre 2024. Le risque, pour NYANGA, combinait une
passerelle exposée sur Internet, non corrigée, et un accès direct au
contrôleur de domaine.

\subsection{Le cycle de vie d'un identifiant CVE}

Un chercheur, un éditeur ou un utilisateur découvre une faille. Une
\termecle{autorité de numérotation} (CNA), souvent l'éditeur lui-même,
réserve un identifiant de la forme \texttt{CVE-AAAA-NNNNN}, puis le publie
avec une description quand le correctif ou l'avis est prêt. La base
nationale américaine NVD~\cite{nistnvd} enrichit ensuite l'enregistrement~:
produits et versions concernés, score, faiblesse CWE. Entre la publication
et l'enrichissement, plusieurs jours peuvent passer~; l'analyste qui attend
le score de la NVD perd ce temps.

% ============================================================================
\section{Le Score CVSS}
\label{sec:c09-cvss}
% ============================================================================

Le \termecle{CVSS} (\textit{Common Vulnerability Scoring System}), tenu par
le FIRST, mesure la gravité d'une vulnérabilité sur une échelle de 0 à 10.
Sa version~4.0, publiée en 2023~\cite{first2023cvss4}, remplace
progressivement la version~3.1~\cite{first2019cvss31}.

\subsection{Quatre groupes de métriques}

\begin{tableaucours}{L{2.8cm}L{5.4cm}L{4.0cm}}{%
    \tableauHeader{Groupe} &
    \tableauHeader{Ce qu'il mesure} &
    \tableauHeader{Qui le renseigne}
}
Base & Les propriétés de la faille, identiques partout & L'éditeur, la
NVD \\
Menace (\textit{Threat}) & La maturité de l'exploitation connue & L'analyste,
à partir du renseignement \\
Environnement & L'importance de l'actif et les mesures en place & L'organisation \\
Supplémentaire & Des informations sans effet sur le score (automatisable,
effort de reprise) & L'éditeur, l'analyste \\
\end{tableaucours}
\vspace{-0.8em}
\begin{center}{\small\itshape Tableau~9.1. Les groupes de métriques de CVSS v4.0}\end{center}

La version~4.0 exige que l'on dise quels groupes ont servi~: CVSS-B pour la
seule base, CVSS-BT avec la menace, CVSS-BE avec l'environnement, CVSS-BTE
avec les trois. Un « 9,3 » sans cette mention est ambigu.

\subsection{Les métriques de base}

\begin{tableaucours}{L{1.6cm}L{4.4cm}L{6.2cm}}{%
    \tableauHeader{Code} &
    \tableauHeader{Métrique} &
    \tableauHeader{Valeurs}
}
AV & Vecteur d'attaque & Réseau (N), adjacent (A), local (L), physique (P) \\
AC & Complexité de l'attaque & Faible (L), élevée (H) \\
AT & Conditions préalables (nouvelle en 4.0) & Aucune (N), présentes (P) \\
PR & Privilèges requis & Aucun (N), faibles (L), élevés (H) \\
UI & Interaction de l'utilisateur & Aucune (N), passive (P), active (A) \\
VC, VI, VA & Impact sur le système vulnérable (confidentialité, intégrité,
disponibilité) & Élevé (H), faible (L), nul (N) \\
SC, SI, SA & Impact sur les systèmes suivants & Élevé (H), faible (L), nul
(N) \\
\end{tableaucours}
\vspace{-0.8em}
\begin{center}{\small\itshape Tableau~9.2. Les métriques de base de CVSS v4.0}\end{center}

La distinction entre système vulnérable et systèmes suivants remplace la
notion de « portée » de la version~3.1. Elle permet de dire qu'une faille de
la passerelle n'atteint pas seulement la passerelle, mais aussi les systèmes
auxquels elle donne accès.

\begin{boxobjet}{Vecteur CVSS v4.0 de VULN-2024-4471}
\textbf{Vecteur publié par l'éditeur (fictif, construit pour le cours)~:}
\texttt{CVSS:4.0/AV:N/AC:L/AT:N/PR:N/UI:N/VC:H/VI:H/VA:H/SC:N/SI:N/SA:N}.

\textbf{Score CVSS-B~:} 9,3, critique.

\textbf{Lecture~:} exploitable depuis Internet, sans condition, sans
compte, sans action d'un utilisateur, avec un impact élevé sur la
passerelle elle-même.

\textbf{Ce que l'éditeur ne pouvait pas savoir~:} qu'à NYANGA, la passerelle
ouvrait un accès direct au contrôleur de domaine. L'analyste de NYANGA, en
renseignant SC, SI et SA à « élevé » dans le groupe environnemental, aurait
obtenu un score de 10,0 pour son propre contexte.
\end{boxobjet}

\subsection{Les seuils de gravité}

\begin{tableaucours}{L{3.4cm}L{3.4cm}}{%
    \tableauHeader{Score} &
    \tableauHeader{Gravité}
}
0,0 & Aucune \\
0,1 à 3,9 & Faible \\
4,0 à 6,9 & Moyenne \\
7,0 à 8,9 & Élevée \\
9,0 à 10,0 & Critique \\
\end{tableaucours}
\vspace{-0.8em}
\begin{center}{\small\itshape Tableau~9.3. Les seuils de gravité du CVSS}\end{center}

\begin{boxinfo}
\textbf{Un score ne se calcule pas de tête.} En version~4.0, le score se
déduit du vecteur par une table de correspondance publiée par le FIRST, et
non par une formule simple comme en version~3.1. On utilise le calculateur
officiel, et l'on note toujours le vecteur à côté du score~: c'est le
vecteur qui permet de discuter le score.
\end{boxinfo}

% ============================================================================
\section{EPSS : la Probabilité d'Exploitation}
\label{sec:c09-epss}
% ============================================================================

L'\termecle{EPSS} (\textit{Exploit Prediction Scoring System}), lui aussi
tenu par le FIRST~\cite{firstepss}, estime chaque jour, pour chaque CVE
publiée, la probabilité qu'elle fasse l'objet de tentatives d'exploitation
observées dans les trente jours suivants. Il s'appuie sur un modèle
statistique alimenté par des données d'exploitation réelles. Il donne une
probabilité, entre 0 et 1, et un centile, qui situe la CVE parmi toutes les
autres.

Le CVSS et l'EPSS répondent à deux questions différentes. Le premier dit
« si cette faille est exploitée, à quel point est-ce grave~? »~; le second
dit « quelle est la chance qu'elle soit exploitée bientôt~? ». Les données
publiées par le projet EPSS~\cite{firstepss} montrent qu'une faible part
des CVE publiées fait l'objet d'une exploitation observée~: un score
critique n'annonce donc pas une exploitation, et un score moyen ne
l'exclut pas.

% ============================================================================
\section{KEV : ce qui est Exploité}
\label{sec:c09-kev}
% ============================================================================

Le catalogue KEV (\textit{Known Exploited Vulnerabilities}), tenu par
l'agence américaine CISA~\cite{cisakev}, recense les vulnérabilités dont
l'exploitation est constatée. Une inscription n'est pas une prévision~:
c'est un fait. Le catalogue fixe une échéance de correction pour les
agences fédérales américaines, mais toute organisation peut s'en servir
comme d'une liste de failles à traiter sans attendre.

\begin{tableaucours}{L{2.4cm}L{4.2cm}L{3.0cm}L{2.6cm}}{%
    \tableauHeader{Indicateur} &
    \tableauHeader{Question} &
    \tableauHeader{Nature} &
    \tableauHeader{Fréquence}
}
CVSS-B & Quelle gravité si exploitée~? & Mesure technique & À la
publication \\
EPSS & Quelle chance d'exploitation dans 30 jours~? & Probabilité estimée
& Quotidienne \\
KEV & Est-elle déjà exploitée~? & Fait constaté & À l'inscription \\
\end{tableaucours}
\vspace{-0.8em}
\begin{center}{\small\itshape Tableau~9.4. Trois indicateurs, trois questions}\end{center}

% ============================================================================
\section{SSVC : Décider plutôt que Noter}
\label{sec:c09-ssvc}
% ============================================================================

La CISA propose aussi une méthode de décision, le SSVC
(\textit{Stakeholder-Specific Vulnerability Categorization})~\cite{cisa2022ssvc}.
Au lieu d'un score, elle conduit, par un arbre de questions, à l'une de
quatre décisions~: suivre (\textit{Track}), suivre de près
(\textit{Track*}), traiter rapidement (\textit{Attend}), agir sans délai
(\textit{Act}). Les questions portent sur l'état de l'exploitation (aucune,
preuve de concept, active), la possibilité d'automatiser l'attaque,
l'impact technique, et l'importance de l'actif pour la mission.

Le SSVC rend explicite ce que le score cache~: la décision dépend de
l'organisation. La même faille peut appeler « agir » sur une passerelle
exposée et « suivre » sur un serveur de test isolé.

% ============================================================================
\section{Quatre Failles chez SOLITRANS}
\label{sec:c09-solitrans}
% ============================================================================

\begin{boxscenario}{Opération SOLITRANS. Quatre lignes du rapport de balayage}
\begin{tableaucours}{L{1.0cm}L{3.8cm}cccL{1.8cm}}{%
    \tableauHeader{Réf.} &
    \tableauHeader{Actif} &
    \tableauHeader{CVSS-B} &
    \tableauHeader{EPSS} &
    \tableauHeader{KEV} &
    \tableauHeader{Exposé}
}
V-A & Serveur d'impression, réseau interne & 9,8 & 0,002 & Non & Non \\
V-B & Portail des transitaires, Internet & 7,5 & 0,640 & Oui & Oui \\
V-C & Base de suivi des conteneurs, interne & 8,8 & 0,031 & Non & Non \\
V-D & Pare-feu du site de Kribi, Internet & 6,5 & 0,410 & Oui & Oui \\
\end{tableaucours}

\smallskip
{\footnotesize Valeurs fictives, construites pour l'exercice. EPSS~:
probabilité estimée d'exploitation dans les trente jours.}

\medskip
\textbf{Constatations.} Rangées par score de base, les failles donnent
l'ordre V-A, V-C, V-B, V-D. Rangées par exploitation constatée puis par
probabilité, elles donnent V-B, V-D, V-C, V-A.

\textbf{Analyse} (confiance élevée). Les deux failles exploitées, sur des
actifs exposés à Internet, passent avant les deux autres, malgré leurs
scores inférieurs. V-A, la plus « grave » sur le papier, touche un serveur
d'impression interne et n'est pas exploitée~: elle reste à corriger, mais
pas cette semaine.
\end{boxscenario}

\begin{boxlecture}{Que sait-on de chaque faille~?}
\begin{tableaucours}{L{5.6cm}cccc}{%
    \tableauHeader{Question} &
    \tableauHeader{V-A} & \tableauHeader{V-B} & \tableauHeader{V-C} &
    \tableauHeader{V-D}
}
La faille est exploitée dans le monde & \NEK{} & \OUI{} & \NEK{} & \OUI{} \\
L'actif est joignable depuis Internet & \NON{} & \OUI{} & \NON{} & \OUI{} \\
La faille est présente (constat vérifié) & \NEK{} & \OUI{} & \OUI{} &
\OUI{} \\
\end{tableaucours}

\medskip
\textbf{Lecture.} Une absence d'inscription au catalogue KEV n'est pas une
preuve de non-exploitation~: V-A et V-C sont \NEK{}, pas \NON{}. Et V-A n'a
pas encore été vérifiée~: l'outil l'a déduite de la version annoncée par le
service, ce qui peut produire un faux positif (Chapitre~\ref{chap:identifier}).
\end{boxlecture}

% ============================================================================
\section{Ce que la Notion demande à chaque Métier}
\label{sec:c09-metiers}
% ============================================================================

\begin{boxmetier}{Analyste en gestion des vulnérabilités}
\textbf{Lire une vulnérabilité}
\begin{enumerate}
    \item Identifier la CVE, la faiblesse CWE et les versions touchées.
    \item Lire le vecteur CVSS, pas seulement le score, et noter la version
          du CVSS.
    \item Relever l'EPSS du jour et l'inscription éventuelle au catalogue
          KEV.
    \item Situer l'actif~: exposition, importance, mesures en place.
    \item En déduire une décision, de type SSVC, et la justifier.
\end{enumerate}
\end{boxmetier}

\begin{boxmetier}{RSSI}
Fixer la règle de priorité par écrit, en combinant exploitation, exposition
et importance de l'actif, et la faire approuver. Une règle fondée sur le
seul score conduit l'équipe à corriger le mauvais serveur en premier.
\end{boxmetier}

\begin{boxmetier}{Analyste en réponse à incident}
Pendant une enquête, le catalogue KEV et l'EPSS orientent la recherche du
point d'entrée~: une faille exploitée, présente sur un actif exposé, est la
première hypothèse à tester.
\end{boxmetier}

\separateur

\begin{boxresume}
\textbf{Ce qu'il faut retenir de ce chapitre~:}
\begin{itemize}
    \item La faiblesse (CWE) est un type d'erreur~; la vulnérabilité (CVE)
          son instance dans un produit~; le risque dépend de la menace et
          du contexte.
    \item Le CVSS v4.0 a quatre groupes de métriques~; on précise CVSS-B,
          -BT, -BE ou -BTE, et l'on donne toujours le vecteur.
    \item L'EPSS estime la probabilité d'exploitation dans les trente
          jours~; le catalogue KEV constate l'exploitation.
    \item Le SSVC remplace le score par une décision qui dépend de
          l'organisation.
    \item Chez SOLITRANS, l'ordre par score et l'ordre par exploitation et
          exposition s'inversent.
\end{itemize}
\end{boxresume}

\begin{boxinfo}
\textbf{Points de contrôle.} Vous avez compris ce chapitre si vous savez~:
\begin{itemize}
    \item expliquer chaque élément d'un vecteur CVSS v4.0~;
    \item dire ce que l'EPSS et le catalogue KEV ajoutent au score~;
    \item justifier un ordre de correction différent de l'ordre des scores.
\end{itemize}
\end{boxinfo}

% ============================================================================
\section*{Exercices}
% ============================================================================

\begin{boxexo}[title={\ding{45}\quad Exercice~9.1. Faiblesse ou vulnérabilité \niveaubase}]
Classez chaque élément comme faiblesse, vulnérabilité, menace ou risque~:
(a) l'injection SQL~; (b) l'injection SQL de la page de connexion du portail
des transitaires, version 3.4~; (c) un groupe qui vend des accès à des
réseaux d'entreprise~; (d) la perte possible des données de 4\,000 clients
si le portail est compromis.
\end{boxexo}

\begin{boxexo}[title={\ding{45}\quad Exercice~9.2. Lire un vecteur \niveaumoyen}]
Expliquez en français le vecteur suivant, métrique par métrique~:
\texttt{CVSS:4.0/AV:L/AC:L/AT:N/PR:L/UI:N/VC:H/VI:H/VA:H/SC:N/SI:N/SA:N}.
Cette faille peut-elle servir de point d'entrée depuis Internet~? À quoi
sert-elle typiquement dans une attaque~?
\end{boxexo}

\begin{boxexo}[title={\ding{45}\quad Exercice~9.3. Trois indicateurs \niveaumoyen}]
Une faille a un score CVSS-B de 5,3, un EPSS de 0,92 et figure au catalogue
KEV. Une autre a un score de 9,8, un EPSS de 0,001 et n'y figure pas.
\begin{enumerate}[(a)]
    \item Laquelle corrigez-vous d'abord sur un serveur exposé à Internet~?
    \item Et si les deux étaient sur un serveur interne isolé~?
    \item Que manque-t-il pour décider avec certitude~?
\end{enumerate}
\end{boxexo}

\begin{boxexo}[title={\ding{45}\quad Exercice~9.4. Le contexte de NYANGA \niveauexpert}]
\begin{enumerate}[(a)]
    \item Expliquez pourquoi l'éditeur a coté SC, SI et SA à « nul » pour
          VULN-2024-4471, et pourquoi NYANGA aurait dû les coter
          autrement.
    \item Le 19~décembre 2024, quelle décision SSVC NYANGA aurait-elle dû
          prendre~? Justifiez chaque branche de l'arbre.
    \item Le jour de l'inscription au catalogue, l'EPSS de la faille
          dépassait 0,9 (valeur fictive). Ce chiffre ajoutait-il quelque
          chose à l'inscription~? Discutez.
\end{enumerate}
\end{boxexo}

\bigskip

\begin{boxinfo}
\textbf{Pour aller plus loin.}
\begin{itemize}
    \item La spécification CVSS~4.0~\cite{first2023cvss4}.
    \item Le site de l'EPSS~\cite{firstepss} et le catalogue
          KEV~\cite{cisakev}.
    \item Le guide SSVC de la CISA~\cite{cisa2022ssvc}.
\end{itemize}
\end{boxinfo}

\begin{boxinfo}
\textbf{Transition.} Coter suppose d'avoir trouvé. Comment savoir quelles
vulnérabilités on a, sans attaquer ses propres systèmes~?
Chapitre~\ref{chap:identifier}.
\end{boxinfo}

        % ============================================================================
% Fichier  : partie4/P04_C10_identifier.tex
% Titre    : Identifier les Vulnérabilités sans les Exploiter
% Auteur   : MINKA MI NGUIDJOÏ Thierry Emmanuel
% Version  : 1.0 -- novembre 2024
% Corrigés : corriges/Corr_C10.tex
% ============================================================================
% Sources : NIST SP 800-115 (techniques d'évaluation) ; NIST SP 800-40 r4
% (gestion des correctifs) ; CIS Benchmarks ; loi n° 2010/012 (art. 65 al.
% 2 et 69, rappel).
% Faits nouveaux (bible) : SOLITRANS, inventaire CMDB 1 380 actifs, 1 512
% découverts (132 inconnus) ; balayage authentifié réussi sur 78 % des
% serveurs ; 3 812 constats bruts -> 2 947 après dédoublonnage ; échantillon
% de 60 constats vérifiés : 9 faux positifs ; incident du 18/05/2025 sur un
% terminal de pesage à Bonabéri. NYANGA : liste de cibles tirée de la CMDB,
% passerelle absente.
% ============================================================================

\chapter{Identifier les Vulnérabilités sans les Exploiter}
\label{chap:identifier}
\chaptermark{Identifier sans exploiter}

\epigraph{%
    « Un outil de balayage ne trouve que ce qu'on lui montre,\\
    et il ne montre que ce qu'il a su lire. »%
}{Principe du cours}

\niveauChapitre{Intermédiaire}{%
    Maîtrise du Chapitre~\ref{chap:vulnerabilite}
}

\begin{boxobjectifs}
\begin{itemize}
    \item Construire et tenir l'inventaire des actifs, condition de toute
          identification.
    \item Choisir entre balayage non authentifié, authentifié, agent et
          analyse de configuration.
    \item Reconnaître les faux positifs et les faux négatifs, et en
          connaître les causes.
    \item Encadrer un balayage~: autorisation, périmètre, fenêtre, effets
          de bord.
    \item Mesurer la couverture d'un programme d'identification.
\end{itemize}
\end{boxobjectifs}

\bigskip

% ============================================================================
\section*{Avant de commencer : la liste des cibles}
% ============================================================================

Bafoussam, avril 2025. Mme~NDOUMBE cherche pourquoi l'outil de balayage de
NYANGA CRÉDIT n'a jamais signalé la faille de la passerelle. L'outil était
acheté, configuré et lancé chaque mois. Elle ouvre sa configuration~: la
liste des adresses à analyser est importée de la base de gestion des
configurations, tenue par le DSI. La passerelle n'y figure pas. Elle a été
installée en 2023 par un prestataire, dans une zone réseau à part, et
personne ne l'a jamais déclarée. L'outil a fait exactement ce qu'on lui
demandait.

\begin{boxalert}
\textbf{Violation typique.} Mesurer un programme de vulnérabilités au
nombre de failles trouvées, sans mesurer la part du parc réellement
analysée.

\smallskip
\textbf{Règle~:} l'identification commence par l'inventaire et se juge à sa
couverture. Un actif absent de l'inventaire, ou analysé sans succès, est un
\NEK{}~: on ne sait rien de ses failles.
\end{boxalert}

\separateur

% ============================================================================
\section{L'Inventaire d'Abord}
\label{sec:c10-inventaire}
% ============================================================================

On ne corrige que ce que l'on connaît. L'inventaire des actifs (matériels,
systèmes, logiciels et services) est le premier contrôle de tout programme,
et le plus souvent défaillant. Il se construit en croisant plusieurs
sources, parce qu'aucune n'est complète.

\begin{tableaucours}{L{3.6cm}L{4.4cm}L{4.2cm}}{%
    \tableauHeader{Source} &
    \tableauHeader{Ce qu'elle voit} &
    \tableauHeader{Ce qu'elle manque}
}
Base de gestion des configurations & Ce qui a été déclaré & Ce qui a été
installé sans déclaration \\
Annuaire (Active Directory) & Machines Windows du domaine & Équipements
réseau, Linux, appareils hors domaine \\
Serveur DHCP & Appareils qui demandent une adresse & Adresses fixes \\
Découverte réseau (balayage des plages) & Tout ce qui répond & Ce qui est
éteint, filtré ou hors des plages \\
Console des fournisseurs en nuage & Ressources en nuage & Ce qui est
ailleurs \\
\end{tableaucours}
\vspace{-0.8em}
\begin{center}{\small\itshape Tableau~10.1. Les sources de l'inventaire}\end{center}

\begin{boxscenario}{Opération SOLITRANS. Les 132 inconnus}
\textbf{Constatations.} La base de gestion des configurations déclare
1\,380 actifs. La découverte réseau des plages du siège, de Bonabéri et de
Kribi, croisée avec l'annuaire et le serveur DHCP, en trouve 1\,512. Les
132 actifs non déclarés comprennent 41 caméras et enregistreurs vidéo, 28
terminaux de pesage et de lecture de conteneurs, 19 postes de prestataires,
17 imprimantes, 14 commutateurs et 13 appareils non identifiés.

\textbf{Analyse} (confiance élevée). Un actif sur onze échappait à toute
analyse. Les catégories manquantes sont celles qu'installent les
prestataires ou les services métier, comme la passerelle de NYANGA.
\end{boxscenario}

% ============================================================================
\section{Quatre Façons d'Identifier}
\label{sec:c10-methodes}
% ============================================================================

\begin{tableaucours}{L{2.8cm}L{3.8cm}L{2.8cm}L{2.8cm}}{%
    \tableauHeader{Méthode} &
    \tableauHeader{Principe} &
    \tableauHeader{Force} &
    \tableauHeader{Faiblesse}
}
Balayage non authentifié & Interroger les services ouverts depuis le
réseau, lire bannières et réponses & Voit ce que voit un attaquant &
Déduit les versions, faux positifs nombreux \\
Balayage authentifié & Se connecter avec un compte en lecture et lire les
paquets installés & Précis sur les versions et correctifs & Exige un
compte~; échecs de connexion silencieux \\
Agent & Programme installé sur la machine, qui remonte son état & Couvre les
portables hors réseau & À déployer et maintenir \\
Analyse de configuration & Comparer la configuration à une référence (CIS
Benchmarks) & Trouve les réglages dangereux, sans CVE & Ne voit pas les
failles logicielles \\
\end{tableaucours}
\vspace{-0.8em}
\begin{center}{\small\itshape Tableau~10.2. Quatre méthodes d'identification}\end{center}

Aucune de ces méthodes n'exploite la faille. Elles constatent qu'un logiciel
vulnérable est présent, ou qu'un réglage est dangereux. Démontrer qu'un
attaquant peut s'en servir relève du test d'intrusion, conduit dans un
autre cadre (UE MSN-2S3-SEC03). Le guide NIST SP~800-115~\cite{nist800115}
décrit les deux familles et leurs précautions.

\begin{boxinfo}
\textbf{Les composants logiciels.} Une application peut embarquer une
bibliothèque vulnérable sans que le système d'exploitation le sache. Les
outils d'analyse de composants lisent les fichiers de l'application, ou une
liste de ses composants fournie par l'éditeur, pour retrouver ces
bibliothèques. Un balayage classique ne les voit pas.
\end{boxinfo}

% ============================================================================
\section{Faux Positifs, Faux Négatifs}
\label{sec:c10-faux}
% ============================================================================

\begin{boxdef}
\textbf{\termecle{Faux positif}}~: vulnérabilité signalée alors qu'elle est
absente. \textbf{\termecle{Faux négatif}}~: vulnérabilité présente et non
signalée.
\end{boxdef}

Les faux positifs coûtent du temps et de la crédibilité~: une équipe système
qui corrige trois failles inexistantes cesse de croire le rapport. Les faux
négatifs coûtent plus cher, parce qu'ils ne se voient pas.

\begin{tableaucours}{L{6.0cm}L{6.2cm}}{%
    \tableauHeader{Cause de faux positif} &
    \tableauHeader{Cause de faux négatif}
}
Correctif rétroporté~: la distribution corrige sans changer le numéro de
version & Actif absent de la liste des cibles \\
Bannière de service modifiée ou générique & Authentification échouée,
analyse réduite au non authentifié \\
Composant présent mais non utilisé & Service filtré par un pare-feu entre
l'outil et la cible \\
Règle de détection trop large & Règle de détection absente pour la faille
ou le produit \\
\end{tableaucours}
\vspace{-0.8em}
\begin{center}{\small\itshape Tableau~10.3. Les causes d'erreur d'un balayage}\end{center}

\begin{boxinfo}
\textbf{Vérifier un constat avant de le transmettre}
\begin{enumerate}
    \item Relire la preuve donnée par l'outil~: version lue, réponse du
          service, paquet installé.
    \item Confirmer la version par une autre voie (connexion à la machine,
          console d'administration).
    \item Vérifier les bulletins de la distribution pour un éventuel
          correctif rétroporté.
    \item Confirmer que le composant vulnérable est actif.
    \item Coter le constat~: confirmé, faux positif, ou non vérifiable.
\end{enumerate}
\end{boxinfo}

\begin{boxscenario}{Opération SOLITRANS. Soixante constats vérifiés}
\textbf{Constatations.} Après dédoublonnage (une même faille signalée par
plusieurs règles, ou sur plusieurs ports d'un même actif), les 3\,812
constats bruts deviennent 2\,947. Mme~EBODE fait vérifier un échantillon de
60 constats tirés au hasard~: 46 sont confirmés, 9 sont des faux positifs
(dont 7 correctifs rétroportés sur des serveurs Linux), 5 ne sont pas
vérifiables faute d'accès à la machine. Par ailleurs, l'authentification a
réussi sur 78~\% des serveurs~; sur les autres, l'outil s'est contenté d'un
balayage non authentifié sans le signaler dans le résumé.

\textbf{Analyse} (confiance moyenne, l'échantillon étant petit). Environ
15~\% des constats seraient des faux positifs. Et 22~\% des serveurs ont
été analysés de façon superficielle~: le nombre de failles réelles sur ces
serveurs est \NEK{}, sans doute supérieur à ce que dit le rapport.
\end{boxscenario}

% ============================================================================
\section{Encadrer un Balayage}
\label{sec:c10-encadrer}
% ============================================================================

Un balayage interroge des machines, parfois des milliers de fois. Il peut
ralentir un service, saturer un équipement fragile ou déclencher des
alertes. Et il touche au droit~: analyser un système sans autorisation peut
constituer un accès non autorisé (loi n°~2010/012, art.~65 al.~2, voir
Chapitre~\ref{chap:normes-droit}).

\begin{tableaucours}{L{3.4cm}L{8.8cm}}{%
    \tableauHeader{Élément} &
    \tableauHeader{Contenu}
}
Autorisation écrite & Signée par le propriétaire des actifs, et par
l'hébergeur ou le prestataire quand l'actif n'appartient pas à
l'organisation \\
Périmètre & Plages d'adresses, exclusions motivées, actifs fragiles \\
Fenêtre & Jours et heures, compatibles avec l'activité \\
Profil & Intensité, tests exclus (déni de service), comptes utilisés \\
Contacts & Personne à joindre en cas d'effet de bord, de part et d'autre \\
Information & Équipe de supervision prévenue, pour ne pas confondre le
balayage avec une attaque \\
\end{tableaucours}
\vspace{-0.8em}
\begin{center}{\small\itshape Tableau~10.4. Le cadre d'un balayage}\end{center}

\begin{boxscenario}{Opération SOLITRANS. Le terminal de pesage}
\textbf{Constatations.} Le 18~mai 2025 à 10~h~20, pendant le balayage du
site de Bonabéri, un terminal de pesage des conteneurs cesse de répondre. La
file des camions s'allonge pendant quarante minutes, jusqu'au redémarrage
manuel. Le terminal, non déclaré dans la base, n'était pas exclu du
périmètre.

\textbf{Analyse} (confiance élevée). Les équipements industriels et
portuaires supportent mal les balayages agressifs. Ils doivent être
inventoriés, exclus du balayage actif ou analysés avec un profil léger, en
dehors des heures d'exploitation.
\end{boxscenario}

% ============================================================================
\section{Mesurer la Couverture}
\label{sec:c10-couverture}
% ============================================================================

\begin{boxlecture}{La couverture de SOLITRANS en mai 2025}
\begin{tableaucours}{L{7.4cm}cc}{%
    \tableauHeader{Question} &
    \tableauHeader{Valeur} &
    \tableauHeader{Cote}
}
Tous les actifs sont inventoriés & 1\,380 sur 1\,512 & \NON{} \\
Tous les actifs inventoriés ont été analysés & 1\,342 sur 1\,380 & \NON{} \\
Les serveurs ont été analysés avec authentification & 78~\% & \NON{} \\
Les constats sont vérifiés & Échantillon de 60 & \NEK{} \\
Les actifs fragiles sont protégés des effets de bord & Incident du 18~mai &
\NON{} \\
\end{tableaucours}

\medskip
\textbf{Lecture.} Quatre \NON{} et un \NEK{}~: le premier balayage dit
beaucoup sur les actifs analysés et presque rien sur la part du parc qui
lui a échappé. Mme~EBODE présente au comité le nombre de failles
\emph{et} ces cinq lignes. Le premier chantier n'est pas la correction,
c'est l'inventaire.
\end{boxlecture}

% ============================================================================
\section{Ce que l'Identification demande à chaque Métier}
\label{sec:c10-metiers}
% ============================================================================

\begin{boxmetier}{Analyste en gestion des vulnérabilités}
\textbf{Identifier sans exploiter}
\begin{enumerate}
    \item Construire l'inventaire en croisant au moins trois sources.
    \item Préférer le balayage authentifié, et suivre le taux de réussite
          de l'authentification.
    \item Dédoublonner, puis vérifier un échantillon pour estimer les faux
          positifs.
    \item Exclure ou ménager les actifs fragiles, avec l'accord des métiers.
    \item Publier la couverture avec le nombre de failles.
\end{enumerate}
\end{boxmetier}

\begin{boxmetier}{RSSI}
Exiger de chaque prestataire qui installe un équipement sa déclaration à
l'inventaire, dans le contrat même. La passerelle de NYANGA et les 132
inconnus de SOLITRANS ont la même origine.
\end{boxmetier}

\begin{boxmetier}{Analyste en réponse à incident}
L'inventaire et l'historique des balayages servent aussi l'enquête~: ils
datent l'apparition d'une version vulnérable et montrent depuis quand une
faille était visible, ou invisible.
\end{boxmetier}

\separateur

\begin{boxresume}
\textbf{Ce qu'il faut retenir de ce chapitre~:}
\begin{itemize}
    \item L'identification commence par l'inventaire, construit en
          croisant plusieurs sources.
    \item Balayage non authentifié, authentifié, agent et analyse de
          configuration se complètent~; aucun n'exploite la faille.
    \item Les faux positifs se vérifient sur échantillon~; les faux négatifs
          se combattent par la couverture.
    \item Un balayage s'encadre~: autorisation écrite, périmètre, fenêtre,
          profil, contacts.
    \item Chez SOLITRANS, un actif sur onze était inconnu et 22~\% des
          serveurs analysés sans authentification.
\end{itemize}
\end{boxresume}

\begin{boxinfo}
\textbf{Points de contrôle.} Vous avez compris ce chapitre si vous savez~:
\begin{itemize}
    \item dire pourquoi un actif hors inventaire est un \NEK{}~;
    \item expliquer un correctif rétroporté et ses effets sur un rapport~;
    \item rédiger le cadre d'un balayage.
\end{itemize}
\end{boxinfo}

% ============================================================================
\section*{Exercices}
% ============================================================================

\begin{boxexo}[title={\ding{45}\quad Exercice~10.1. Quelle méthode~? \niveaubase}]
Choisissez une méthode d'identification pour~: (a) les portables des
commerciaux, rarement au bureau~; (b) les serveurs Linux du siège~; (c) le
portail des transitaires vu d'Internet~; (d) la configuration des
commutateurs~; (e) une application métier qui embarque des bibliothèques
open source.
\end{boxexo}

\begin{boxexo}[title={\ding{45}\quad Exercice~10.2. Faux positif ou pas \niveaumoyen}]
L'outil signale sur un serveur Debian une faille d'OpenSSH corrigée dans la
version 9.6, alors que le serveur annonce la version 9.2.
\begin{enumerate}[(a)]
    \item Pourquoi ce constat peut-il être un faux positif~?
    \item Comment le vérifier sans exploiter la faille~?
    \item Comment cotez-vous le constat avant et après vérification~?
\end{enumerate}
\end{boxexo}

\begin{boxexo}[title={\ding{45}\quad Exercice~10.3. La couverture \niveaumoyen}]
À partir de l'encadré « La couverture de SOLITRANS en mai 2025 »~:
\begin{enumerate}[(a)]
    \item Calculez le taux de couverture global (actifs analysés sur actifs
          existants).
    \item Si les 15~\% de faux positifs se confirment, combien des 2\,947
          constats seraient réels~?
    \item Pourquoi ce second chiffre est-il probablement trop bas~?
\end{enumerate}
\end{boxexo}

\begin{boxexo}[title={\ding{45}\quad Exercice~10.4. La passerelle invisible \niveauexpert}]
\begin{enumerate}[(a)]
    \item Décrivez trois contrôles qui auraient fait entrer la passerelle de
          NYANGA dans le périmètre du balayage.
    \item Une analyse non authentifiée de la passerelle, depuis Internet,
          aurait-elle détecté VULN-2024-4471~? À quelles conditions~?
    \item Rédigez le cadre (Tableau~10.4) d'un balayage mensuel de la
          zone réseau de la passerelle, installée et maintenue par un
          prestataire.
\end{enumerate}
\end{boxexo}

\bigskip

\begin{boxinfo}
\textbf{Pour aller plus loin.}
\begin{itemize}
    \item Le NIST SP~800-115~\cite{nist800115}, sur les techniques
          d'évaluation.
    \item Les CIS Benchmarks~\cite{cis2026benchmarks}, pour les références
          de configuration.
\end{itemize}
\end{boxinfo}

\begin{boxinfo}
\textbf{Transition.} SOLITRANS connaît désormais près de trois mille
failles, et sait qu'elle ne les connaît pas toutes. Lesquelles corriger
d'abord, dans quels délais, et comment suivre~? Chapitre~\ref{chap:gerer}.
\end{boxinfo}

        % ============================================================================
% Fichier  : partie4/P04_C11_gerer_vulnerabilites.tex
% Titre    : Gérer les Vulnérabilités
% Auteur   : MINKA MI NGUIDJOÏ Thierry Emmanuel
% Version  : 1.0 -- novembre 2026
% Corrigés : corriges/Corr_C11.tex
% ============================================================================
% Sources : NIST SP 800-40 r4 (gestion des correctifs comme maintenance
% préventive) ; ISO/IEC 30111 et 29147 (traitement et divulgation côté
% éditeur) ; ISO/IEC 27002:2022, mesure 8.8 (gestion des vulnérabilités
% techniques, paraphrasée) ; catalogue KEV ; EPSS.
% Faits nouveaux (bible) : politique SOLITRANS de juin 2025, classes P1 à P4
% et délais 7/15/30/90 jours ; suivi T3 2025 à T1 2026 ; exceptions E-01 à
% E-04 ; NYANGA adopte la méthode en septembre 2025.
% ============================================================================

\chapter{Gérer les Vulnérabilités}
\label{chap:gerer}
\chaptermark{Gérer les vulnérabilités}

\epigraph{%
    « Une vulnérabilité corrigée est une vulnérabilité\\
    dont on a vérifié la correction. »%
}{Principe du cours}

\niveauChapitre{Intermédiaire}{%
    Maîtrise des Chapitres~\ref{chap:vulnerabilite} et~\ref{chap:identifier}
}

\begin{boxobjectifs}
\begin{itemize}
    \item Décrire le cycle de gestion des vulnérabilités, de
          l'identification à la vérification.
    \item Rédiger une politique de priorité qui combine exploitation,
          exposition et importance de l'actif, avec des délais.
    \item Choisir entre correctif, changement de configuration et mesure
          compensatoire.
    \item Gérer une exception~: propriétaire, durée, mesure compensatoire,
          revue.
    \item Construire et lire les indicateurs d'un programme.
\end{itemize}
\end{boxobjectifs}

\bigskip

% ============================================================================
\section*{Avant de commencer : quarante failles par semaine}
% ============================================================================

Douala, juin 2025. Mme~EBODE présente sa politique au comité de direction
de SOLITRANS. Elle ne propose pas de corriger plus~: l'équipe système de
M.~FOUDA traite environ quarante failles par semaine, et rien ne permettra
de doubler ce rythme avant l'an prochain. Elle propose de corriger dans le
bon ordre, avec des délais écrits, et de vérifier ce qui est fait. Le
directeur financier demande combien de failles resteront ouvertes à la fin
de l'année. Elle répond qu'il en restera beaucoup, et qu'elle saura dire
lesquelles, pourquoi, et qui en a accepté le risque.

\begin{boxalert}
\textbf{Violation typique.} Déclarer une faille corrigée sur la foi du
ticket fermé par l'équipe système, sans nouvelle analyse de l'actif.

\smallskip
\textbf{Règle~:} une correction se vérifie par un nouveau balayage, ou par
une lecture directe de la version installée. Le ticket fermé est une
déclaration~; le balayage qui ne trouve plus la faille est une constatation.
\end{boxalert}

\separateur

% ============================================================================
\section{Le Cycle de Gestion}
\label{sec:c11-cycle}
% ============================================================================

Le guide NIST SP~800-40~\cite{nist80040} présente la gestion des correctifs
comme une maintenance préventive, planifiée, et non comme une suite
d'urgences. La mesure 8.8 de l'ISO/IEC~27002:2022~\cite{iso27002_2022},
consacrée à la gestion des vulnérabilités techniques, demande d'obtenir
l'information sur les vulnérabilités, d'évaluer l'exposition de
l'organisation et de prendre les mesures appropriées (paraphrase).

\begin{center}
\begin{tikzpicture}[
    etape/.style = {draw=CouleurBordDef, fill=CouleurDef, rounded corners=3pt,
        minimum width=2.4cm, minimum height=0.9cm, align=center,
        font=\small\bfseries\color{CouleurPrimaire}},
    arrow/.style = {-{Stealth[length=4pt]}, thick, color=CouleurSecondaire}
]
\node[etape] (a) at (0,0)     {Inventorier};
\node[etape] (b) at (3.0,0)   {Identifier};
\node[etape] (c) at (6.0,0)   {Prioriser};
\node[etape] (d) at (9.0,0)   {Traiter};
\node[etape] (e) at (9.0,-1.6) {Vérifier};
\node[etape] (f) at (4.5,-1.6) {Rendre compte};
\draw[arrow] (a) -- (b); \draw[arrow] (b) -- (c); \draw[arrow] (c) -- (d);
\draw[arrow] (d) -- (e); \draw[arrow] (e) -- (f); \draw[arrow] (f) -| (a);
\end{tikzpicture}
\end{center}
\begin{center}{\small\itshape Figure~11.1. Le cycle de gestion des vulnérabilités}\end{center}

Traiter ne veut pas toujours dire corriger. Quatre voies existent.

\begin{tableaucours}{L{2.8cm}L{5.0cm}L{4.4cm}}{%
    \tableauHeader{Voie} &
    \tableauHeader{Exemple} &
    \tableauHeader{Condition}
}
Correctif & Mise à jour de l'éditeur & Testé, appliqué dans une fenêtre,
retour arrière prévu \\
Configuration & Désactiver le service vulnérable, changer un réglage &
Fonction non utilisée, ou réglage documenté par l'éditeur \\
Mesure compensatoire & Filtrage réseau, restriction d'accès, règle de
pare-feu applicatif & Temporaire, en attendant le correctif \\
Acceptation & Exception formelle & Propriétaire, durée, revue (section~11.4) \\
\end{tableaucours}
\vspace{-0.8em}
\begin{center}{\small\itshape Tableau~11.1. Les quatre voies de traitement}\end{center}

% ============================================================================
\section{La Politique de Priorité}
\label{sec:c11-politique}
% ============================================================================

Une politique de priorité traduit en règles écrites ce que le
Chapitre~\ref{chap:vulnerabilite} a montré~: l'exploitation et l'exposition
comptent davantage que le score seul.

\begin{boxscenario}{Opération SOLITRANS. La politique de juin 2025}
\begin{tableaucours}{cL{6.6cm}cL{2.6cm}}{%
    \tableauHeader{Classe} &
    \tableauHeader{Règle d'entrée} &
    \tableauHeader{Délai} &
    \tableauHeader{Décideur d'exception}
}
P1 & Inscrite au catalogue KEV, ou EPSS $\geq 0{,}5$, sur un actif exposé à
Internet ou critique & 7 jours & Directeur général \\
P2 & Inscrite au catalogue KEV sur un actif interne, ou CVSS-B $\geq 9{,}0$
sur un actif exposé & 15 jours & RSSI \\
P3 & CVSS-B $\geq 7{,}0$, hors P1 et P2 & 30 jours & RSSI \\
P4 & Autres constats confirmés & 90 jours & Responsable des
vulnérabilités \\
\end{tableaucours}

\medskip
\textbf{Constatations.} Appliquée aux 2\,947 constats de mai, la politique
place 38 constats en P1, 211 en P2, 904 en P3 et 1\,794 en P4. À quarante
corrections par semaine, les 249 constats P1 et P2 demandent un peu plus de
six semaines, bien au-delà du délai de 15~jours fixé pour les P2.

\textbf{Analyse} (confiance élevée). La politique est tenable pour les P1,
pas encore pour le stock initial de P2~: Mme~EBODE demande un renfort
temporaire de deux semaines, et la fixation du délai à partir de la date
d'entrée en vigueur pour le stock existant.
\end{boxscenario}

\begin{boxinfo}
\textbf{Les délais courent depuis quand~?} Le délai d'une faille se compte à
partir du moment où l'organisation pouvait la connaître~: la publication du
correctif, ou la première détection sur son parc, selon la règle choisie.
Une politique qui ne le précise pas permet de repousser indéfiniment
l'échéance.
\end{boxinfo}

% ============================================================================
\section{Corriger en Production}
\label{sec:c11-production}
% ============================================================================

Un correctif peut casser une application. C'est la raison la plus souvent
donnée pour ne pas l'appliquer, et elle est parfois exacte. La réponse n'est
pas de renoncer, mais d'organiser~: un environnement de test qui reproduit
la production pour les applications critiques, des fenêtres de maintenance
connues des métiers, une procédure de retour arrière éprouvée, et un ordre
de déploiement par vagues (d'abord quelques machines représentatives, puis
le reste).

\begin{boxinfo}
\textbf{Déployer un correctif de classe P1}
\begin{enumerate}
    \item Lire le bulletin~: versions touchées, prérequis, effets connus.
    \item Appliquer une mesure compensatoire immédiate si le correctif
          demande plus de 24~heures.
    \item Tester sur une machine représentative, avec le propriétaire
          métier.
    \item Déployer par vagues, en surveillant les services.
    \item Vérifier par un nouveau balayage authentifié de chaque actif.
    \item Lever la mesure compensatoire, et clore le constat avec la
          preuve.
\end{enumerate}
\end{boxinfo}

% ============================================================================
\section{Les Exceptions}
\label{sec:c11-exceptions}
% ============================================================================

Certaines failles ne peuvent pas être corrigées dans le délai~: le
correctif n'existe pas, l'éditeur ne maintient plus le produit, ou
l'application métier ne fonctionne pas avec la version corrigée. L'exception
transforme ce fait en décision.

\begin{boxdef}
\textbf{\termecle{Exception}}~: décision écrite par laquelle un responsable
habilité accepte, pour une durée limitée, qu'une vulnérabilité reste
ouverte au-delà du délai, avec une mesure compensatoire et une date de
revue.
\end{boxdef}

\begin{boxscenario}{Opération SOLITRANS. Quatre exceptions}
\begin{tableaucours}{cL{3.6cm}L{3.4cm}cL{1.8cm}}{%
    \tableauHeader{Réf.} &
    \tableauHeader{Faille et actif} &
    \tableauHeader{Mesure compensatoire} &
    \tableauHeader{Durée} &
    \tableauHeader{Décideur}
}
E-01 & P2, logiciel de dédouanement en fin de support & Accès limité à
trois postes, réseau isolé & 6 mois & RSSI \\
E-02 & P3, terminaux de pesage, aucun correctif & Filtrage strict, aucun
accès Internet & 12 mois & RSSI \\
E-03 & P1, portail des transitaires, correctif incompatible & Règle de
pare-feu applicatif & 30 jours & Directeur général \\
E-04 & P3, serveur d'archives & Aucune & Sans limite & Non signée \\
\end{tableaucours}

\medskip
\textbf{Analyse} (confiance élevée). E-01 à E-03 sont des exceptions. E-04
n'en est pas une~: sans mesure, sans durée et sans signature, c'est une
faille ouverte que personne n'assume. Elle redevient un constat en retard.
\end{boxscenario}

% ============================================================================
\section{Mesurer le Programme}
\label{sec:c11-mesurer}
% ============================================================================

\begin{tableaucours}{L{4.0cm}L{5.2cm}L{3.0cm}}{%
    \tableauHeader{Indicateur} &
    \tableauHeader{Définition} &
    \tableauHeader{Piège}
}
Délai de correction & Durée médiane entre détection et correction
vérifiée, par classe & La moyenne masque les cas extrêmes \\
Respect des délais & Part des constats corrigés dans le délai de leur classe
& À calculer sur les constats échus \\
Couverture & Part des actifs existants analysés avec authentification &
Souvent omise \\
KEV ouvertes & Nombre de failles du catalogue encore ouvertes & Le seul
indicateur à suivre chaque semaine \\
Âge du stock & Âge des constats ouverts, par tranche & Le stock vieillit
sans bruit \\
\end{tableaucours}
\vspace{-0.8em}
\begin{center}{\small\itshape Tableau~11.2. Les indicateurs d'un programme}\end{center}

\begin{boxlecture}{SOLITRANS, trois trimestres}
\begin{tableaucours}{L{5.6cm}ccc}{%
    \tableauHeader{Indicateur} &
    \tableauHeader{T3 2025} &
    \tableauHeader{T4 2025} &
    \tableauHeader{T1 2026}
}
Couverture authentifiée & 71~\% & 84~\% & 93~\% \\
Délai médian P1 (jours) & 11 & 6 & 5 \\
P1 corrigées dans le délai & 58~\% & 81~\% & 92~\% \\
P2 corrigées dans le délai & 34~\% & 63~\% & 77~\% \\
KEV ouvertes en fin de trimestre & 14 & 5 & 2 \\
Exceptions en cours, toutes signées & \NON{} & \OUI{} & \OUI{} \\
\end{tableaucours}

\medskip
\textbf{Lecture.} Le programme progresse sur toutes les lignes. La
couverture passe de 71 à 93~\%~: une partie des nouvelles failles du T4
vient simplement de ce que l'on regarde mieux, et la hausse apparente du
stock ne traduit pas une dégradation. Les deux KEV ouvertes au T1 2026 sont
les terminaux de pesage, sous l'exception E-02.
\end{boxlecture}

\begin{boxscenario}{Opération NYANGA. Septembre 2025}
\textbf{Constatations.} Six mois après l'incident, NYANGA CRÉDIT adopte la
méthode de SOLITRANS~: inventaire croisé (la passerelle et quatorze autres
équipements y entrent), balayage authentifié mensuel, politique en quatre
classes, suivi hebdomadaire des KEV. Mme~NDOUMBE devient aussi responsable
du programme.

\textbf{Analyse} (confiance moyenne). Avec cette politique, VULN-2024-4471
aurait été classée P1 le 19~décembre 2024, avec une échéance au
26~décembre, soit plus de deux mois avant le premier accès.
\end{boxscenario}

% ============================================================================
\section{Ce que la Gestion demande à chaque Métier}
\label{sec:c11-metiers}
% ============================================================================

\begin{boxmetier}{Analyste en gestion des vulnérabilités}
\textbf{Tenir le programme}
\begin{enumerate}
    \item Classer chaque constat confirmé selon la politique écrite.
    \item Suivre chaque semaine les failles du catalogue KEV ouvertes.
    \item Vérifier chaque correction par un nouveau balayage.
    \item Instruire les exceptions et rappeler leurs dates de revue.
    \item Publier les indicateurs avec la couverture.
\end{enumerate}
\end{boxmetier}

\begin{boxmetier}{RSSI}
Faire approuver la politique et les délais par la direction, et signer les
exceptions de son niveau. Refuser une exception sans mesure compensatoire ni
durée~: elle ne protège personne, sauf de la responsabilité.
\end{boxmetier}

\begin{boxmetier}{Analyste en réponse à incident}
Après un incident, vérifier que la faille exploitée est corrigée partout où
elle existe, et pas seulement sur la machine compromise.
\end{boxmetier}

\separateur

\begin{boxresume}
\textbf{Ce qu'il faut retenir de ce chapitre~:}
\begin{itemize}
    \item Le cycle va de l'inventaire à la vérification~; traiter peut
          vouloir dire corriger, reconfigurer, compenser ou accepter.
    \item Une politique écrite classe les failles par exploitation,
          exposition et importance, avec des délais datés.
    \item Une exception a un décideur, une mesure compensatoire, une durée
          et une revue~; sinon, c'est un retard.
    \item Une correction se vérifie par un nouveau balayage~; un ticket
          fermé ne suffit pas.
    \item Les indicateurs se lisent avec la couverture~: une hausse des
          failles peut traduire un meilleur regard.
\end{itemize}
\end{boxresume}

\begin{boxinfo}
\textbf{Points de contrôle.} Vous avez compris ce chapitre si vous savez~:
\begin{itemize}
    \item classer une faille dans une politique en quatre classes~;
    \item distinguer une exception d'un retard~;
    \item interpréter une hausse du nombre de failles quand la couverture
          augmente.
\end{itemize}
\end{boxinfo}

% ============================================================================
\section*{Exercices}
% ============================================================================

\begin{boxexo}[title={\ding{45}\quad Exercice~11.1. Classer \niveaubase}]
Avec la politique de SOLITRANS, classez~: (a) faille inscrite au catalogue
KEV sur le portail des transitaires~; (b) faille de score 9,4 sur un serveur
interne, non exploitée~; (c) faille de score 7,2 sur un poste~; (d) faille
inscrite au catalogue KEV sur un serveur interne~; (e) faille de score 5,0
sur une imprimante.
\end{boxexo}

\begin{boxexo}[title={\ding{45}\quad Exercice~11.2. Le stock de P2 \niveaumoyen}]
\begin{enumerate}[(a)]
    \item Vérifiez le calcul de l'encadré « La politique de juin 2025 »~:
          combien de semaines pour absorber les P1 et les P2 à quarante
          corrections par semaine~?
    \item Combien de corrections par semaine faudrait-il pour traiter les
          211 P2 en 15 jours, après les 38 P1~?
    \item Proposez deux solutions autres que le renfort.
\end{enumerate}
\end{boxexo}

\begin{boxexo}[title={\ding{45}\quad Exercice~11.3. Instruire une exception \niveaumoyen}]
Le portail des transitaires présente une faille P1~; le correctif casse le
module de facturation, et l'éditeur annonce une version compatible dans six
semaines. Rédigez la demande d'exception complète, et dites qui doit la
signer.
\end{boxexo}

\begin{boxexo}[title={\ding{45}\quad Exercice~11.4. Lire les indicateurs \niveauexpert}]
Au T4 2025, le nombre total de failles ouvertes de SOLITRANS augmente de
12~\% alors que tous les indicateurs de délai s'améliorent.
\begin{enumerate}[(a)]
    \item Proposez deux explications compatibles avec l'encadré « SOLITRANS,
          trois trimestres ».
    \item Comment les départager~?
    \item Rédigez la phrase du rapport trimestriel qui présente ce chiffre au
          comité sans l'alarmer à tort ni le rassurer à tort.
\end{enumerate}
\end{boxexo}

\bigskip

\begin{boxinfo}
\textbf{Pour aller plus loin.}
\begin{itemize}
    \item Le NIST SP~800-40~\cite{nist80040}.
    \item L'ISO/IEC~30111~\cite{iso30111} et l'ISO/IEC~29147~\cite{iso29147},
          pour le traitement et la divulgation du côté des éditeurs.
\end{itemize}
\end{boxinfo}

\begin{boxinfo}
\textbf{Transition.} Le programme de SOLITRANS empêche les incidents~;
l'enquête de NYANGA en explique un. Le Chapitre~\ref{chap:artefact-vuln}
relie les deux~: comment remonter d'une trace à la faille qui l'a permise.
\end{boxinfo}

        % ============================================================================
% Fichier  : partie4/P04_C12_artefact_vulnerabilite.tex
% Titre    : De l'Artefact à la Vulnérabilité
% Auteur   : MINKA MI NGUIDJOÏ Thierry Emmanuel
% Version  : 1.0 -- novembre 2026
% Corrigés : corriges/Corr_C12.tex
% ============================================================================
% Sources : ISO/IEC 27035 (enseignements tirés d'un incident, paraphrasée) ;
% NIST SP 800-61 r2 (activité post-incident) ; analyse des causes (arbre des
% causes, cinq pourquoi : méthodes générales, présentées sans citation).
% Faits nouveaux (bible) : export de configuration de la passerelle,
% micrologiciel 7.2.1 ; 4 requêtes du 02/03 entre 23 h 46 et 23 h 48 vers le
% chemin décrit comme indicateur par le bulletin de l'éditeur ; svc_backup
% membre du groupe des administrateurs du domaine ; aucune authentification
% à deux facteurs sur la passerelle ; clé de l'interface en clair dans
% api.ini ; 3 autres équipements de bordure non corrigés trouvés en
% septembre 2025.
% ============================================================================

\chapter{De l'Artefact à la Vulnérabilité}
\label{chap:artefact-vuln}
\chaptermark{De l'artefact à la vulnérabilité}

\epigraph{%
    « Trouver par où l'attaquant est entré, c'est la moitié de l'enquête.\\
    Comprendre pourquoi la porte était ouverte, c'est l'autre moitié. »%
}{Principe du cours}

\niveauChapitre{Intermédiaire à expert}{%
    Maîtrise des Chapitres~\ref{chap:chronologie} et~\ref{chap:gerer}
}

\begin{boxobjectifs}
\begin{itemize}
    \item Remonter d'un artefact au mécanisme qui l'a produit, et de ce
          mécanisme à une vulnérabilité identifiée.
    \item Confronter des traces aux indicateurs publiés par un éditeur.
    \item Examiner un équipement de bordure dont le système n'est pas
          accessible, et en dire les limites.
    \item Distinguer cause immédiate, causes contributives et cause
          profonde.
    \item Rechercher les variantes d'une faille dans le reste du parc.
    \item Transformer les enseignements d'une enquête en règles du programme
          de vulnérabilités.
\end{itemize}
\end{boxobjectifs}

\bigskip

% ============================================================================
\section*{Avant de commencer : une session sans mot de passe}
% ============================================================================

Yaoundé, 6 avril 2025. Trois semaines d'enquête ont établi la séquence de
l'attaque. Il reste un point que M.~ABOUBAKAR ne veut pas laisser à
l'inférence~: la première session du compte \texttt{svc\_backup}, le 2~mars,
s'est ouverte sans aucun échec d'authentification. Soit l'attaquant
connaissait le mot de passe, soit il n'en a pas eu besoin. La réponse
détermine la cause de l'incident, et elle pèse sur la question encore
ouverte de la participation d'un employé.

\begin{boxalert}
\textbf{Violation typique.} Désigner la faille la plus connue du moment
comme cause de l'incident, parce qu'elle touche le produit en cause, sans
confronter les traces à ses indicateurs.

\smallskip
\textbf{Règle~:} une vulnérabilité n'est la cause d'un incident que si les
traces portent sa signature, décrite par l'éditeur ou par un tiers
compétent, et si les autres explications sont écartées.
\end{boxalert}

\separateur

% ============================================================================
\section{Trois Explications pour une Session}
\label{sec:c12-explications}
% ============================================================================

\begin{tableaucours}{L{4.2cm}L{4.2cm}L{3.8cm}}{%
    \tableauHeader{Explication} &
    \tableauHeader{Traces attendues} &
    \tableauHeader{Traces observées}
}
E1~: mot de passe connu (fuite, employé, devinette) & Session directe, sans
requête inhabituelle avant & Requêtes inhabituelles avant la session \\
E2~: exploitation de VULN-2024-4471 & Requêtes vers le chemin décrit dans le
bulletin, puis session sans authentification & Quatre requêtes de ce type,
puis la session \\
E3~: vol de jeton d'une session légitime & Session légitime de
\texttt{svc\_backup} juste avant & Aucune session légitime depuis 9~jours \\
\end{tableaucours}
\vspace{-0.8em}
\begin{center}{\small\itshape Tableau~12.1. Trois explications de la session du 2~mars}\end{center}

Le bulletin publié par l'éditeur le 12~décembre 2024 décrit la faille, les
versions touchées (jusqu'à 7.2.4) et des \termecle{indicateurs de
compromission}~: la forme particulière des requêtes HTTP envoyées à la
passerelle lors de l'exploitation, et leur trace dans son journal. C'est la
signature à rechercher.

\begin{boxscenario}{Opération NYANGA. La signature du 2 mars}
\textbf{Constatations.}
\begin{itemize}
    \item L'export de configuration de la passerelle, réalisé le 14~mars
          (scellé NC-02), indique le micrologiciel 7.2.1, compris dans les
          versions touchées.
    \item Le journal de la passerelle enregistre, le 2~mars entre 23~h~46 et
          23~h~48 (heures corrigées), quatre requêtes depuis 203.0.113.18
          vers le chemin désigné par le bulletin comme indicateur
          d'exploitation, suivies de la session de \texttt{svc\_backup}.
    \item Aucune requête de ce type n'apparaît ailleurs dans les 60~jours
          du journal.
    \item Le mot de passe de \texttt{svc\_backup}, long de 24 caractères, n'a
          pas changé depuis 2022 et n'apparaît dans aucune fuite connue
          consultée.
\end{itemize}

\textbf{Analyse} (confiance élevée). Les traces portent la signature de
VULN-2024-4471 et précèdent immédiatement la session. E2 explique la session
sans authentification~; E1 ne l'explique pas sans supposer, en plus, que
l'attaquant connaissait un mot de passe qu'il n'a pas eu à saisir~; E3 est
écartée.
\end{boxscenario}

\begin{boxlecture}{Les trois explications au 6 avril}
\begin{tableaucours}{L{7.0cm}c}{%
    \tableauHeader{Proposition} &
    \tableauHeader{Cote}
}
E2~: la session du 2~mars résulte de l'exploitation de VULN-2024-4471 &
\OUI{} \\
E3~: la session résulte d'un vol de jeton d'une session légitime & \NON{} \\
E1~: l'attaquant connaissait le mot de passe de \texttt{svc\_backup} &
\NEK{} \\
Un employé a participé à l'attaque & \NEK{} \\
\end{tableaucours}

\medskip
\textbf{Lecture.} E1 reste \NEK{}~: qu'on n'ait pas eu besoin du mot de
passe ne prouve pas qu'on l'ignorait. Mais l'hypothèse perd son utilité~:
l'entrée s'explique sans elle. La participation d'un employé reste \NEK{},
et elle le restera tant que l'identité des auteurs l'est~; l'enquête ne
dispose d'aucune pièce qui la soutienne.
\end{boxlecture}

% ============================================================================
\section{Examiner un Équipement de Bordure}
\label{sec:c12-bordure}
% ============================================================================

Les passerelles d'accès, les pare-feu et les routeurs sont des appliances~:
leur système est fermé, souvent sans accès au disque ni à la mémoire.
L'enquêteur dispose de ce que l'équipement veut bien exporter.

\begin{tableaucours}{L{4.0cm}L{4.2cm}L{4.0cm}}{%
    \tableauHeader{Pièce} &
    \tableauHeader{Ce qu'elle apporte} &
    \tableauHeader{Limite}
}
Export de configuration & Version, comptes, règles, réglages & Instantané
au moment de l'export \\
Journaux exportés & Sessions, requêtes, erreurs & Rotation, niveau de détail
choisi \\
Outil d'intégrité de l'éditeur & Comparaison des fichiers système avec les
fichiers d'origine & Disponible pour certains produits seulement \\
Image du système, avec l'aide de l'éditeur & Examen complet & Rare, lent,
dépend de l'éditeur \\
\end{tableaucours}
\vspace{-0.8em}
\begin{center}{\small\itshape Tableau~12.2. Les pièces d'un équipement de bordure}\end{center}

Un attaquant qui contrôle une passerelle peut y avoir laissé un mécanisme de
retour qui survit à la mise à jour. Sans outil d'intégrité ni image du
système, la proposition « la passerelle ne contient plus de code de
l'attaquant » reste \NEK{}. À NYANGA, l'éditeur ne fournissait pas d'outil
d'intégrité pour la version 7.2~: la passerelle a été remplacée par un
équipement neuf, et l'ancienne conservée sous scellé.

% ============================================================================
\section{Cause Immédiate, Causes Contributives, Cause Profonde}
\label{sec:c12-causes}
% ============================================================================

\begin{boxdef}
\textbf{\termecle{Cause immédiate}}~: l'événement qui a directement permis
l'incident. \textbf{\termecle{Cause contributive}}~: condition qui a
aggravé l'incident ou en a facilité une étape. \textbf{\termecle{Cause
profonde}}~: défaut d'organisation qui a laissé subsister la cause
immédiate et les causes contributives~; la corriger empêche des incidents
semblables.
\end{boxdef}

La méthode des cinq pourquoi aide à descendre de la cause immédiate vers la
cause profonde~: on demande « pourquoi~? » à chaque réponse, jusqu'à
atteindre une décision ou une absence de décision de l'organisation.

\begin{boxscenario}{Opération NYANGA. Cinq pourquoi}
\begin{enumerate}
    \item Pourquoi l'attaquant est-il entré~? Parce que la passerelle
          présentait VULN-2024-4471.
    \item Pourquoi la faille était-elle présente~? Parce que le correctif du
          15~décembre n'avait pas été appliqué.
    \item Pourquoi n'avait-il pas été appliqué~? Parce que personne ne savait
          que la passerelle était concernée~: elle n'apparaissait dans aucun
          rapport de balayage.
    \item Pourquoi n'apparaissait-elle pas~? Parce qu'elle était absente de
          la base de gestion des configurations, qui fournissait la liste des
          cibles.
    \item Pourquoi était-elle absente~? Parce qu'aucune règle n'obligeait un
          prestataire à déclarer les équipements qu'il installe, ni personne
          à vérifier l'inventaire contre le réseau réel.
\end{enumerate}
\end{boxscenario}

\begin{boxlecture}{L'arbre des causes de NYANGA}
\begin{tableaucours}{L{6.2cm}L{2.4cm}c}{%
    \tableauHeader{Cause} &
    \tableauHeader{Type} &
    \tableauHeader{Établie}
}
VULN-2024-4471 non corrigée sur la passerelle & Immédiate & \OUI{} \\
\texttt{svc\_backup} membre des administrateurs du domaine & Contributive &
\OUI{} \\
Aucune authentification à deux facteurs sur la passerelle & Contributive &
\OUI{} \\
Identifiants d'administration présents en mémoire du contrôleur & Contributive
& \OUI{} \\
Clé de l'interface de paiement en clair dans \texttt{api.ini} & Contributive
& \OUI{} \\
Journaux de la passerelle conservés 60~jours & Contributive (enquête) &
\OUI{} \\
Inventaire non confronté au réseau réel, prestataires non tenus de
déclarer & Profonde & \OUI{} \\
\end{tableaucours}

\medskip
\textbf{Lecture.} Sept causes, toutes établies par des pièces. Corriger la
seule cause immédiate (appliquer le correctif) laisserait intactes les six
autres~: la prochaine faille d'un équipement non déclaré aurait le même
effet. La recommandation principale du rapport porte sur la cause profonde.
\end{boxlecture}

\begin{boxinfo}
\textbf{L'authentification à deux facteurs n'aurait pas suffi.}
VULN-2024-4471 contourne l'authentification~: un second facteur, vérifié par
le même mécanisme défaillant, aurait pu être contourné aussi. On la classe
en cause contributive parce qu'elle aurait gêné les sessions ultérieures sur
d'autres comptes, pas parce qu'elle aurait empêché l'entrée. Le rapport doit
le dire, pour ne pas vendre une mesure pour ce qu'elle n'est pas.
\end{boxinfo}

% ============================================================================
\section{Chercher les Variantes}
\label{sec:c12-variantes}
% ============================================================================

Une faille exploitée chez soi se cherche ensuite partout où elle peut
exister. C'est l'\termecle{analyse de variantes}~: la même vulnérabilité sur
d'autres équipements, la même faiblesse dans d'autres produits, la même
cause profonde ailleurs.

\begin{boxscenario}{Opération NYANGA. Septembre 2025}
\textbf{Constatations.} L'inventaire croisé de septembre 2025
(Chapitre~\ref{chap:gerer}) fait entrer quinze équipements dans la base.
Parmi eux, trois équipements de bordure (deux routeurs d'agence et un
point d'accès de télémaintenance) présentent des failles inscrites au
catalogue KEV, non corrigées depuis plus de six mois. Une recherche dans les
fichiers des serveurs trouve deux autres clés d'interface en clair, pour la
messagerie d'alerte et pour un service de change.

\textbf{Analyse} (confiance élevée). La cause profonde produisait ses effets
ailleurs~: trois autres portes ouvertes, deux autres clés exposées.
L'enquête sur un incident en a évité au moins un autre.
\end{boxscenario}

% ============================================================================
\section{Des Enseignements aux Règles}
\label{sec:c12-regles}
% ============================================================================

L'ISO/IEC~27035~\cite{iso27035} et le guide NIST SP~800-61~\cite{nist80061r2}
font des enseignements tirés d'un incident une étape à part entière. Ils ne
valent que s'ils deviennent des règles vérifiables.

\begin{tableaucours}{L{4.6cm}L{7.6cm}}{%
    \tableauHeader{Enseignement} &
    \tableauHeader{Règle du programme de NYANGA (septembre 2025)}
}
Un équipement non déclaré n'est pas analysé & Découverte réseau mensuelle de
toutes les plages, écarts avec la base traités sous 15~jours \\
Les équipements de bordure sont la porte d'entrée & Toute faille inscrite au
catalogue KEV sur un équipement exposé est classée P1, délai 72~heures \\
Un prestataire installe sans déclarer & Clause contractuelle de déclaration
et d'accès au balayage \\
Les journaux de la passerelle ont failli disparaître & Conservation de dix
ans des données de connexion (loi n°~2010/012, art.~29) \\
Une clé en clair a permis les virements & Coffre de secrets pour toutes les
clés d'interface, recherche trimestrielle de secrets en clair \\
\end{tableaucours}
\vspace{-0.8em}
\begin{center}{\small\itshape Tableau~12.3. Des enseignements aux règles}\end{center}

% ============================================================================
\section{Ce que la Remontée demande à chaque Métier}
\label{sec:c12-metiers}
% ============================================================================

\begin{boxmetier}{Analyste en réponse à incident}
\textbf{Remonter à la vulnérabilité}
\begin{enumerate}
    \item Lister toutes les explications du premier accès.
    \item Pour chacune, écrire les traces attendues.
    \item Confronter les traces aux indicateurs publiés par l'éditeur.
    \item Coter chaque explication, et écrire ce qui reste \NEK{}.
    \item Construire l'arbre des causes jusqu'à la cause profonde.
    \item Lancer la recherche de variantes dans tout le parc.
\end{enumerate}
\end{boxmetier}

\begin{boxmetier}{Analyste en gestion des vulnérabilités}
Recevoir chaque rapport d'incident comme une mesure de son programme~: la
faille exploitée était-elle connue, classée, dans le délai~? Chaque réponse
« non » appelle une règle nouvelle.
\end{boxmetier}

\begin{boxmetier}{Application OPJ/Expert Judiciaire}
L'expert établit le mécanisme d'entrée parce qu'il répond à l'article~69
de la loi n°~2010/012 (violation des mesures de sécurité). Il ne juge pas
la diligence de la victime, sauf si sa mission le lui demande.
\end{boxmetier}

\begin{boxmetier}{RSSI}
Présenter à la direction l'arbre des causes complet, et faire décider la
correction de la cause profonde. Un correctif appliqué en une heure ne
clôt pas l'incident~; une règle d'inventaire tenue chaque mois le clôt.
\end{boxmetier}

\separateur

\begin{boxresume}
\textbf{Ce qu'il faut retenir de ce chapitre~:}
\begin{itemize}
    \item Une vulnérabilité est la cause d'un incident quand les traces
          portent sa signature et que les autres explications sont
          écartées.
    \item Un équipement de bordure s'examine par ses exports~; ce qu'on ne
          peut pas lire reste \NEK{}.
    \item On distingue cause immédiate, causes contributives et cause
          profonde~; la recommandation principale vise la cause profonde.
    \item L'analyse de variantes cherche la même faille et la même cause
          ailleurs.
    \item Un enseignement ne vaut que transformé en règle vérifiable.
\end{itemize}
\end{boxresume}

\begin{boxinfo}
\textbf{Points de contrôle.} Vous avez compris ce chapitre si vous savez~:
\begin{itemize}
    \item confronter des traces aux indicateurs d'un bulletin~;
    \item construire un arbre des causes à trois niveaux~;
    \item rédiger une règle de programme à partir d'un enseignement.
\end{itemize}
\end{boxinfo}

% ============================================================================
\section*{Exercices}
% ============================================================================

\begin{boxexo}[title={\ding{45}\quad Exercice~12.1. Immédiate, contributive, profonde \niveaubase}]
Classez~: (a) le correctif non appliqué~; (b) l'absence de règle de
déclaration des équipements~; (c) la clé en clair dans \texttt{api.ini}~;
(d) le compte \texttt{svc\_backup} trop privilégié~; (e) la durée de
conservation des journaux.
\end{boxexo}

\begin{boxexo}[title={\ding{45}\quad Exercice~12.2. Une autre signature \niveaumoyen}]
Supposez que le journal de la passerelle n'ait montré aucune requête vers le
chemin du bulletin, mais trois échecs d'authentification de
\texttt{svc\_backup} depuis 203.0.113.18, puis une session réussie.
\begin{enumerate}[(a)]
    \item Comment cotez-vous E1 et E2~?
    \item Quelle cause immédiate retenez-vous alors, et quelles causes
          contributives changent~?
    \item Que devient la question de la participation d'un employé~?
\end{enumerate}
\end{boxexo}

\begin{boxexo}[title={\ding{45}\quad Exercice~12.3. Les cinq pourquoi de SOLITRANS \niveaumoyen}]
Le terminal de pesage de Bonabéri a cessé de répondre pendant un balayage
(Chapitre~\ref{chap:identifier}). Appliquez la méthode des cinq pourquoi et
identifiez la cause profonde. Proposez la règle qui en découle.
\end{boxexo}

\begin{boxexo}[title={\ding{45}\quad Exercice~12.4. La passerelle sous scellé \niveauexpert}]
\begin{enumerate}[(a)]
    \item Pourquoi a-t-on remplacé la passerelle au lieu de simplement la
          mettre à jour~?
    \item Quelles questions l'ancienne passerelle, conservée sous scellé,
          pourrait-elle encore trancher, et avec l'aide de qui~?
    \item Rédigez le paragraphe du rapport d'expertise qui établit la
          violation des mesures de sécurité au sens de l'article~69, en
          séparant constatation et analyse.
\end{enumerate}
\end{boxexo}

\bigskip

\begin{boxinfo}
\textbf{Pour aller plus loin.}
\begin{itemize}
    \item L'ISO/IEC~27035~\cite{iso27035}, sur la gestion des incidents et
          leurs enseignements.
    \item Le NIST SP~800-61~\cite{nist80061r2}, section sur l'activité
          post-incident.
\end{itemize}
\end{boxinfo}

\begin{boxinfo}
\textbf{Transition.} L'attaquant de NYANGA n'a presque rien effacé. Un
attaquant plus soigneux maquille ses dates, vide les journaux, ou travaille
dans des systèmes que l'enquêteur ne touche pas~: Chapitre~\ref{chap:antiforensique}.
\end{boxinfo}


    % PARTIE 5 · DÉFIS AVANCÉS
    \part{Défis Avancés}
    \label{part:avance}

        % ============================================================================
% Fichier  : partie5/P05_C13_antiforensique_nuage.tex
% Titre    : Anti-forensique, Nuage et Conteneurs
% Auteur   : MINKA MI NGUIDJOÏ Thierry Emmanuel
% Version  : 1.0 -- novembre 2024
% Corrigés : corriges/Corr_C13.tex
% ============================================================================
% Sources : Carrier (2005) pour les attributs NTFS ; ISO/IEC 27037 et 27042
% (paraphrasées) ; modèle de responsabilité partagée de l'informatique en
% nuage (présenté sans citation) ; journaux d'audit de Kubernetes (présentés
% sans citation).
% Faits nouveaux (bible) : sur SRV-FIC01, outil upd.exe daté de 2017 dans
% $STANDARD_INFORMATION, du 13/03/2025 dans $FILE_NAME et $UsnJrnl ; journal
% de sécurité de SRV-FIC01 vidé le 14/03 à 3 h 52 (1102). SOLITRANS,
% février 2026 : instance du portail en nuage compromise, journaux d'audit
% du fournisseur conservés 90 jours, conteneur supprimé par la reprise
% automatique.
% ============================================================================

\chapter{Anti-forensique, Nuage et Conteneurs}
\label{chap:antiforensique}
\chaptermark{Anti-forensique, nuage, conteneurs}

\epigraph{%
    « Effacer une trace en laisse une autre.\\
    Il faut savoir où la chercher. »%
}{Principe du cours}

\niveauChapitre{Expert}{%
    Maîtrise de la partie~II
}

\begin{boxobjectifs}
\begin{itemize}
    \item Reconnaître les principales techniques d'anti-forensique et les
          traces qu'elles laissent.
    \item Détecter une falsification de dates par la confrontation des
          attributs NTFS et des journaux.
    \item Exploiter les traces d'un journal effacé.
    \item Conduire une acquisition dans un environnement en nuage, selon le
          modèle de responsabilité partagée.
    \item Préserver les traces d'un conteneur avant sa disparition.
    \item Coter \NEK{} ce que l'environnement rend inaccessible.
\end{itemize}
\end{boxobjectifs}

\bigskip

% ============================================================================
\section*{Avant de commencer : un outil de 2017}
% ============================================================================

Yaoundé, 28 mars 2025. Dans l'acquisition à froid de SRV-FIC01, M.~FOTSO
trouve l'outil de copie de l'attaquant~: \texttt{upd.exe}, rangé dans un
dossier de mises à jour. L'explorateur affiche une date de création du
4~juin 2017. Si elle est vraie, l'outil était sur le serveur bien avant
l'attaque, et peut-être installé par quelqu'un de la maison. M.~FOTSO ouvre
l'enregistrement MFT du fichier.

\begin{boxalert}
\textbf{Violation typique.} Prendre pour acquise une date affichée par
l'explorateur ou par un outil, sans la confronter aux autres sources qui
datent le même fichier.

\smallskip
\textbf{Règle~:} une date se recoupe. Quand plusieurs sources indépendantes
datent un même objet, leur désaccord est une constatation, et souvent la
trace d'une manipulation.
\end{boxalert}

\separateur

% ============================================================================
\section{Les Techniques d'Anti-forensique}
\label{sec:c13-techniques}
% ============================================================================

\begin{boxdef}
\textbf{\termecle{Anti-forensique}}~: ensemble des techniques qui visent à
empêcher, retarder ou tromper l'investigation~: effacer, masquer, falsifier
ou rendre illisibles les traces.
\end{boxdef}

\begin{tableaucours}{L{3.2cm}L{4.4cm}L{4.6cm}}{%
    \tableauHeader{Technique} &
    \tableauHeader{Ce qu'elle fait} &
    \tableauHeader{Trace qu'elle laisse}
}
Effacement des journaux & Vide le journal de sécurité ou système &
Événement 1102 ou 104, trou dans la numérotation, sources externes
intactes \\
Falsification de dates & Réécrit les dates d'un fichier & Désaccord entre
\texttt{\$STANDARD\_INFORMATION}, \texttt{\$FILE\_NAME} et
\texttt{\$UsnJrnl} \\
Effacement sécurisé & Écrase le contenu d'un fichier & Nom et dates dans
la MFT et le \texttt{\$UsnJrnl}, outil d'effacement lui-même \\
Chiffrement & Rend le contenu illisible & Le conteneur chiffré, l'outil,
parfois la clé en mémoire \\
Exécution en mémoire, outils du système & N'écrit rien sur le disque,
utilise \texttt{powershell}, \texttt{rundll32} & Mémoire, lignes de
commande, journaux de script s'ils sont activés \\
Dissimulation & Cache des fichiers ou des processus & Écarts entre deux
méthodes de lecture (\texttt{pslist} et \texttt{psscan}) \\
\end{tableaucours}
\vspace{-0.8em}
\begin{center}{\small\itshape Tableau~13.1. Techniques d'anti-forensique et traces résiduelles}\end{center}

Le tableau montre une constante~: chaque technique déplace la trace plus
qu'elle ne la supprime. L'enquêteur cherche donc là où l'attaquant ne
pense pas ou ne peut pas agir~: les journaux centralisés hors de la machine,
la mémoire, les structures internes du système de fichiers, les équipements
réseau.

% ============================================================================
\section{Détecter une Falsification de Dates}
\label{sec:c13-dates}
% ============================================================================

Les outils de falsification réécrivent en général les dates de
\texttt{\$STANDARD\_INFORMATION}, que tout programme peut modifier. Les
dates de \texttt{\$FILE\_NAME} et les entrées du \texttt{\$UsnJrnl} sont plus
difficiles à atteindre. Trois indices orientent l'examen~: une date de
création antérieure à l'installation du système, des fractions de seconde
toutes nulles (les outils écrivent souvent des secondes rondes), et un
désaccord entre les attributs.

\begin{boxscenario}{Opération NYANGA. L'outil upd.exe}
\begin{tableaucours}{L{4.6cm}L{4.0cm}L{3.0cm}}{%
    \tableauHeader{Source} &
    \tableauHeader{Date de création} &
    \tableauHeader{Fractions}
}
\texttt{\$STANDARD\_INFORMATION} & 04/06/2017 09:12:00 & ,0000000 \\
\texttt{\$FILE\_NAME} & 13/03/2025 23:58:41 & ,3318274 \\
\texttt{\$UsnJrnl}, création du fichier & 13/03/2025 23:58:41 & \\
Installation du système (registre) & 17/01/2020 & \\
\end{tableaucours}

\medskip
\textbf{Constatations.} La date de \texttt{\$STANDARD\_INFORMATION} précède de
deux ans et demi l'installation du système, et ses fractions sont nulles.
\texttt{\$FILE\_NAME} et le \texttt{\$UsnJrnl} concordent sur le 13~mars
2025 à 23~h~58.

\textbf{Analyse} (confiance élevée). La date de 2017 a été falsifiée.
L'outil a été déposé le 13~mars à 23~h~58, soit soixante-quatorze minutes
avant le début de l'exfiltration. La piste d'un outil ancien, installé par
un employé, tombe.
\end{boxscenario}

% ============================================================================
\section{Un Journal Vidé}
\label{sec:c13-journal}
% ============================================================================

\begin{boxscenario}{Opération NYANGA. Le journal de SRV-FIC01}
\textbf{Constatations.} Le journal de sécurité de SRV-FIC01 commence par un
événement 1102, le 14~mars à 3~h~52, enregistré sous le compte
\texttt{adm.kamga}~: le journal a été vidé, cinq minutes après la fin de
l'exfiltration. Les journaux du contrôleur de domaine, eux, conservent les
authentifications de \texttt{adm.kamga} vers SRV-FIC01 le 13 à 23~h~55 et
le 14 à 1~h~09.

\textbf{Analyse} (confiance élevée). L'attaquant a effacé les traces locales
de ses sessions, mais pas celles que le contrôleur de domaine garde sur les
authentifications. L'effacement lui-même est une constatation, datée et
attribuée à un compte.
\end{boxscenario}

\begin{boxinfo}
\textbf{La centralisation des journaux.} Un journal copié en temps réel sur
un serveur de collecte, hors de portée des comptes d'administration
ordinaires, survit à l'effacement local. C'est la mesure la plus simple
contre l'effacement des journaux, et NYANGA ne l'avait pas.
\end{boxinfo}

% ============================================================================
\section{L'Enquête dans le Nuage}
\label{sec:c13-nuage}
% ============================================================================

Dans l'informatique en nuage, l'organisation ne possède ni les machines ni
le réseau. Le \termecle{modèle de responsabilité partagée} répartit la
sécurité entre le fournisseur et le client~: plus le service est géré par
le fournisseur, moins le client a de prise, et moins l'enquêteur a de
traces directement accessibles.

\begin{tableaucours}{L{2.4cm}L{3.6cm}L{3.2cm}L{3.0cm}}{%
    \tableauHeader{Modèle} &
    \tableauHeader{Ce que gère le client} &
    \tableauHeader{Traces accessibles} &
    \tableauHeader{Traces chez le fournisseur}
}
Infrastructure (IaaS) & Système, applications, données & Disques des
machines (instantanés), journaux du système & Hyperviseur, réseau physique \\
Plateforme (PaaS) & Applications, données & Journaux de l'application et du
service & Système sous-jacent \\
Logiciel (SaaS) & Données, comptes & Journaux d'audit exposés par le
service & Presque tout le reste \\
\end{tableaucours}
\vspace{-0.8em}
\begin{center}{\small\itshape Tableau~13.2. Traces et responsabilité partagée}\end{center}

L'acquisition y change de forme. On ne débranche pas un serveur~: on prend
un instantané de ses disques par l'interface du fournisseur, on exporte les
journaux d'audit de la console, on copie les journaux du réseau virtuel s'ils
sont activés. La mémoire d'une machine virtuelle ne se capture que par un
outil exécuté à l'intérieur, ou par un service spécialisé du fournisseur.
Les durées de conservation par défaut des journaux d'audit sont souvent
courtes, et les données peuvent se trouver dans un autre pays que
l'organisation, ce qui fait intervenir l'entraide internationale
(Chapitre~\ref{chap:normes-droit}).

\begin{boxscenario}{Opération SOLITRANS. Le portail en février 2026}
\textbf{Constatations.} Le 9~février 2026, le portail des transitaires,
hébergé chez un fournisseur en nuage sous forme de conteneurs, présente des
pages modifiées. L'équipe de M.~NJIKE prend un instantané des disques des
machines virtuelles et exporte les journaux d'audit de la console, conservés
90~jours. Le conteneur compromis, lui, a déjà été remplacé~: le système de
reprise automatique l'a détruit et relancé à partir de l'image d'origine
quand il a cessé de répondre, à 4~h~10.

\textbf{Analyse} (confiance élevée). Le contenu du conteneur compromis est
perdu~; seuls restent les journaux que l'orchestrateur et le fournisseur
ont gardés. La reprise automatique, qui sert la disponibilité, a détruit la
preuve.
\end{boxscenario}

% ============================================================================
\section{Les Conteneurs}
\label{sec:c13-conteneurs}
% ============================================================================

Un conteneur est éphémère par conception~: il naît d'une image, vit
quelques minutes ou quelques jours, et disparaît avec ses fichiers et ses
journaux s'ils n'ont pas été envoyés ailleurs. L'enquête doit donc préparer
la préservation avant l'incident.

\begin{tableaucours}{L{4.0cm}L{8.2cm}}{%
    \tableauHeader{Trace} &
    \tableauHeader{Comment la préserver}
}
Image d'origine & Conserver son empreinte (\textit{digest}) et l'image
elle-même dans le registre \\
Modifications du conteneur & Copier la couche modifiable avant de détruire
le conteneur, ou le suspendre au lieu de le redémarrer \\
Journaux de l'application & Les envoyer en continu vers un collecteur
externe \\
Journaux de l'orchestrateur & Activer et conserver le journal d'audit
(créations, suppressions, accès aux secrets) \\
Réseau & Journaux des flux du réseau virtuel \\
\end{tableaucours}
\vspace{-0.8em}
\begin{center}{\small\itshape Tableau~13.3. Préserver les traces d'un conteneur}\end{center}

\begin{boxlecture}{Ce que l'enquête peut encore savoir, NYANGA et SOLITRANS}
\begin{tableaucours}{L{7.4cm}c}{%
    \tableauHeader{Proposition} &
    \tableauHeader{Cote}
}
\texttt{upd.exe} a été déposé sur SRV-FIC01 le 13~mars 2025 & \OUI{} \\
\texttt{upd.exe} était présent en 2017 & \NON{} \\
Le journal de SRV-FIC01 a été vidé par le compte \texttt{adm.kamga} &
\OUI{} \\
Les sessions de l'attaquant sur SRV-FIC01 sont datées & \OUI{} \\
Le conteneur du portail de SOLITRANS contenait un outil de l'attaquant &
\NEK{} \\
Le point d'entrée sur le portail de SOLITRANS est identifié & \NEK{} \\
\end{tableaucours}

\medskip
\textbf{Lecture.} Sur NYANGA, l'anti-forensique a échoué~: chaque trace
effacée ou falsifiée a été recoupée ailleurs. Sur SOLITRANS, deux \NEK{}
tiennent à l'architecture elle-même~: ce que la reprise automatique a
détruit ne se reconstitue pas.
\end{boxlecture}

% ============================================================================
\section{Ce que ces Défis demandent à chaque Métier}
\label{sec:c13-metiers}
% ============================================================================

\begin{boxmetier}{Analyste en réponse à incident}
\textbf{Face à l'anti-forensique}
\begin{enumerate}
    \item Recouper chaque date par au moins deux sources indépendantes.
    \item Chercher les journaux hors de la machine (contrôleur, collecteur,
          réseau).
    \item Traiter chaque effacement comme une constatation datée.
    \item Dans le nuage, exporter d'abord les journaux d'audit, dont la
          conservation est courte.
    \item Suspendre un conteneur suspect au lieu de le laisser redémarrer.
\end{enumerate}
\end{boxmetier}

\begin{boxmetier}{RSSI}
Centraliser les journaux hors de portée des comptes d'administration~;
négocier avec les fournisseurs en nuage la durée de conservation des
journaux d'audit~; régler la reprise automatique des conteneurs pour
qu'elle isole avant de détruire.
\end{boxmetier}

\begin{boxmetier}{Application OPJ/Expert Judiciaire}
Pour les données détenues par un fournisseur à l'étranger, l'expert indique
dans son rapport ce qu'il n'a pas pu examiner et pourquoi, et l'OPJ engage
sans délai la demande de conservation puis l'entraide.
\end{boxmetier}

\separateur

\begin{boxresume}
\textbf{Ce qu'il faut retenir de ce chapitre~:}
\begin{itemize}
    \item L'anti-forensique déplace les traces plus qu'elle ne les
          supprime~; on les cherche là où l'attaquant n'agit pas.
    \item Une falsification de dates se détecte par le désaccord entre
          \texttt{\$STANDARD\_INFORMATION}, \texttt{\$FILE\_NAME} et
          \texttt{\$UsnJrnl}.
    \item Un journal vidé laisse un événement daté et attribué~; les
          journaux centralisés survivent.
    \item Dans le nuage, l'acquisition passe par les instantanés et les
          journaux d'audit, aux durées courtes.
    \item Un conteneur se préserve avant l'incident~; la reprise automatique
          peut détruire la preuve.
\end{itemize}
\end{boxresume}

\begin{boxinfo}
\textbf{Points de contrôle.} Vous avez compris ce chapitre si vous savez~:
\begin{itemize}
    \item démontrer une falsification de dates à partir de trois sources~;
    \item dire ce qu'un client peut acquérir selon le modèle de service~;
    \item préparer la préservation d'un environnement de conteneurs.
\end{itemize}
\end{boxinfo}

% ============================================================================
\section*{Exercices}
% ============================================================================

\begin{boxexo}[title={\ding{45}\quad Exercice~13.1. Trace résiduelle \niveaubase}]
Pour chaque technique, citez une trace résiduelle à chercher~: (a) un
journal de sécurité vidé~; (b) un fichier écrasé par un outil d'effacement~;
(c) un script exécuté en mémoire par \texttt{powershell}~; (d) une archive
chiffrée.
\end{boxexo}

\begin{boxexo}[title={\ding{45}\quad Exercice~13.2. Vraie ou fausse date \niveaumoyen}]
Un fichier \texttt{rapport.xlsx} présente une création au 12/03/2025
14:05:33,8812094 dans \texttt{\$STANDARD\_INFORMATION} et au 12/03/2025
14:05:33,8812094 dans \texttt{\$FILE\_NAME}~; le \texttt{\$UsnJrnl} ne
contient plus d'entrée aussi ancienne.
\begin{enumerate}[(a)]
    \item La date est-elle falsifiée~?
    \item Pourquoi le \texttt{\$UsnJrnl} ne permet-il pas de trancher~?
    \item Comment cotez-vous la proposition « le fichier a été créé le
          12~mars »~?
\end{enumerate}
\end{boxexo}

\begin{boxexo}[title={\ding{45}\quad Exercice~13.3. Acquérir dans le nuage \niveaumoyen}]
Une machine virtuelle de SOLITRANS, hébergée en IaaS, est suspectée.
Décrivez, dans l'ordre, les acquisitions à faire, et dites pour chacune ce
qui serait perdu si on la faisait trop tard.
\end{boxexo}

\begin{boxexo}[title={\ding{45}\quad Exercice~13.4. Préparer les conteneurs \niveauexpert}]
\begin{enumerate}[(a)]
    \item Rédigez cinq règles de préparation qui auraient permis d'examiner
          le conteneur compromis du portail de SOLITRANS.
    \item La reprise automatique sert la disponibilité. Proposez un réglage
          qui concilie disponibilité et préservation.
    \item Reprenez l'encadré « Ce que l'enquête peut encore savoir » et dites
          lesquelles des deux cotes \NEK{} de SOLITRANS vos règles auraient
          pu changer.
\end{enumerate}
\end{boxexo}

\bigskip

\begin{boxinfo}
\textbf{Pour aller plus loin.}
\begin{itemize}
    \item B.~Carrier, \textit{File System Forensic
          Analysis}~\cite{carrier2005}, sur les attributs NTFS.
    \item L'ISO/IEC~27042~\cite{iso27042_2015}, sur l'analyse et
          l'interprétation.
\end{itemize}
\end{boxinfo}

\begin{boxinfo}
\textbf{Transition.} Tous les outils sont en place~: acquérir, analyser,
dater, remonter à la faille, contourner l'anti-forensique. Le
Chapitre~\ref{chap:enquete} les rassemble dans une enquête complète, du
premier appel au dernier scellé.
\end{boxinfo}

        % ============================================================================
% Fichier  : partie5/P05_C14_enquete_integratrice.tex
% Titre    : L'Enquête Intégratrice
% Auteur   : MINKA MI NGUIDJOÏ Thierry Emmanuel
% Version  : 1.0 -- novembre 2026
% Corrigés : corriges/Corr_C14.tex
% ============================================================================
% Chapitre de synthèse : reprend les Chapitres 1 à 13 dans une méthode en dix
% étapes, appliquée à NYANGA (récapitulatif) puis proposée pour SOLITRANS
% (février 2026, exercice intégrateur et lab FAV-LAB-10).
% Sources : ISO/IEC 27035, 27043 ; NIST SP 800-61 r2 et 800-86.
% Faits nouveaux (bible) : coût de l'incident NYANGA 118 M FCFA (47,3 M de
% virements dont 12,6 M recouvrés, enquête 21 M, remplacement et mesures
% 34 M, arrêt des agences 28,3 M) ; durée de la mission SANKOFA 23 jours.
% ============================================================================

\chapter{L'Enquête Intégratrice}
\label{chap:enquete}
\chaptermark{L'enquête intégratrice}

\epigraph{%
    « Une enquête bien conduite se reconnaît à ce qu'un autre\\
    pourrait la refaire à partir du dossier. »%
}{Principe du cours}

\niveauChapitre{Expert}{%
    Maîtrise des Chapitres~\ref{chap:investigation} à~\ref{chap:antiforensique}
}

\begin{boxobjectifs}
\begin{itemize}
    \item Conduire une enquête complète en dix étapes, de l'alerte à la
          clôture.
    \item Articuler réponse à incident, enquête technique et enquête
          judiciaire sur une même affaire.
    \item Tenir le dossier d'enquête de façon qu'il soit reproductible.
    \item Relier l'enquête au programme de gestion des vulnérabilités.
    \item Estimer le coût d'un incident et le présenter à une direction.
\end{itemize}
\end{boxobjectifs}

\bigskip

% ============================================================================
\section*{Avant de commencer : vingt-trois jours}
% ============================================================================

Yaoundé, 6 avril 2025. La mission de SANKOFA EXPERTISE à NYANGA CRÉDIT
s'achève après vingt-trois jours. Mme~MBARGA réunit son équipe pour le
dernier exercice de chaque mission~: relire le dossier comme le relirait un
expert adverse, page par page, en se demandant pour chaque conclusion d'où
elle vient. Elle ne cherche pas les erreurs de l'attaquant, mais les siennes.

\begin{boxalert}
\textbf{Violation typique.} Conduire chaque étape correctement, mais sans
fil~: des acquisitions sans questions, des analyses sans chronologie, une
chronologie sans arbre des causes, un rapport sans recommandations
vérifiables.

\smallskip
\textbf{Règle~:} une enquête suit une méthode écrite, dont chaque étape
produit une pièce du dossier et prépare la suivante. Le dossier, pas la
mémoire des enquêteurs, porte l'enquête.
\end{boxalert}

\separateur

% ============================================================================
\section{Une Méthode en Dix Étapes}
\label{sec:c14-methode}
% ============================================================================

\begin{tableaucours}{cL{3.4cm}L{5.0cm}L{2.6cm}}{%
    \tableauHeader{N°} &
    \tableauHeader{Étape} &
    \tableauHeader{Pièce produite} &
    \tableauHeader{Chapitre}
}
1 & Recevoir l'alerte, isoler & Registre des actions ouvert & \ref{chap:prologue} \\
2 & Cadrer~: mandat, questions, cadre juridique & Lettre de mission,
questions écrites & \ref{chap:investigation}, \ref{chap:normes-droit} \\
3 & Formuler les hypothèses & Liste des hypothèses et pièces réfutantes &
\ref{chap:investigation} \\
4 & Préserver et acquérir & Formulaires de chaîne de possession, scellés &
\ref{chap:preuve}, \ref{chap:acquisition} \\
5 & Examiner disque, mémoire, réseau & Notes d'examen, sorties d'outils &
\ref{chap:disque} à \ref{chap:reseau} \\
6 & Construire la chronologie & Tableau des horloges, chronologie cotée &
\ref{chap:chronologie} \\
7 & Remonter à la vulnérabilité & Matrice des explications, arbre des
causes & \ref{chap:artefact-vuln} \\
8 & Chercher les variantes & Liste des autres instances & \ref{chap:artefact-vuln} \\
9 & Rendre compte & Rapport, annexes, recommandations &
\ref{chap:rapport} \\
10 & Clore et transmettre & Restitution ou transfert des scellés, règles
nouvelles au programme & \ref{chap:gerer} \\
\end{tableaucours}
\vspace{-0.8em}
\begin{center}{\small\itshape Tableau~14.1. La méthode du cours en dix étapes}\end{center}

Les étapes se chevauchent~: on acquiert encore pendant qu'on examine, et
chaque découverte peut rouvrir une hypothèse. La méthode fixe l'ordre des
dépendances, pas un calendrier. Elle respecte les classes de processus de
l'ISO/IEC~27043~\cite{iso27043_2015} et les phases du NIST
SP~800-86~\cite{nist80086}, en y ajoutant explicitement les étapes~7 et~8,
qui relient l'enquête à la gestion des vulnérabilités.

% ============================================================================
\section{Trois Enquêtes, un Dossier}
\label{sec:c14-trois}
% ============================================================================

À NYANGA CRÉDIT, la réponse à incident, l'enquête du cabinet et l'enquête
judiciaire ont avancé ensemble. Leur coordination tient à quelques règles.

\begin{tableaucours}{L{4.4cm}L{7.8cm}}{%
    \tableauHeader{Règle} &
    \tableauHeader{Application à NYANGA}
}
Un seul registre des actions & Ouvert par la RSSI le 14~mars, repris par le
cabinet, versé à la procédure \\
Les originaux vont à la justice & Scellés NC-01 à NC-07 remis à l'OPJ le
16 et le 22~mars~; le cabinet travaille sur ses copies \\
Chacun garde son rapport & Le cabinet rend compte à l'établissement,
l'expert au procureur \\
La reprise d'activité attend la préservation & Contrôleur de secours en
service après la capture de la mémoire et l'image du disque \\
Les communications passent par la cellule & Aucune déclaration publique
avant validation par la juriste \\
\end{tableaucours}
\vspace{-0.8em}
\begin{center}{\small\itshape Tableau~14.2. Coordonner les trois enquêtes}\end{center}

\begin{boxinfo}
\textbf{Les obligations de notification.} Un incident qui expose des données
de clients peut déclencher des obligations de notification envers une
autorité de protection des données et envers les personnes concernées, au
titre de la loi n°~2024/017~\cite{loi2024017}. Un établissement de
microfinance peut aussi devoir informer son autorité de supervision [nature
et délai à vérifier dans la réglementation applicable]. Le détail relève de
l'UE MSN-2S3-JUR01~; l'enquêteur fournit les faits qui permettent de
remplir ces obligations~: quelles données, combien de personnes, depuis
quand.
\end{boxinfo}

% ============================================================================
\section{Le Dossier NYANGA}
\label{sec:c14-dossier}
% ============================================================================

\begin{boxscenario}{Opération NYANGA. Le dossier au 6 avril 2025}
\begin{tableaucours}{L{3.8cm}L{8.4cm}}{%
    \tableauHeader{Rubrique} &
    \tableauHeader{Contenu}
}
Mandat et questions & Lettre de mission du 14~mars~; cinq questions (point
d'entrée, séquence, données sorties, auteurs, mesures) \\
Registre des actions & 214 entrées, du 14~mars 8~h~15 au 6~avril \\
Scellés & NC-01 à NC-07, formulaires, procès-verbaux de remise \\
Notes d'examen & Disque (SRV-CPT01, SRV-AD01, SRV-FIC01), mémoire
(SRV-AD01, SRV-FIC01), réseau (NC-04), sorties brutes des outils en
annexe \\
Chronologie & 14 lignes cotées, tableau des horloges, temps de présence \\
Causes & Matrice E1-E2-E3, arbre des sept causes \\
Variantes & Recherche lancée, programme de septembre \\
Rapport & Version du 6~avril, relue contradictoirement \\
\end{tableaucours}

\medskip
\textbf{Réponses aux cinq questions.} Point d'entrée~: \OUI{}
(VULN-2024-4471). Séquence~: \OUI{}, sauf du 3 au 8~mars (\NEK{}). Données
sorties~: \OUI{} pour le dossier \texttt{Clients} par le siège, \NEK{} pour
d'autres chemins. Auteurs~: \NEK{}. Mesures~: sept recommandations.
\end{boxscenario}

\begin{boxlecture}{Les cinq questions de la mission}
\textbf{Décompte.} Trois questions établies (point d'entrée, séquence
principale, données sorties par le siège), une non tranchée (auteurs), et
deux réserves partielles (trou du 3 au 8~mars, autres chemins de sortie).

\medskip
\textbf{Lecture.} La mission répond à ce que l'établissement doit savoir
pour décider~: fermer la porte, mesurer la fuite, prévenir les clients. Elle
ne répond pas à la question pénale de l'identité, qui dépend de pièces
détenues par des tiers à l'étranger.
\end{boxlecture}

% ============================================================================
\section{Le Coût d'un Incident}
\label{sec:c14-cout}
% ============================================================================

La direction demande toujours combien l'incident a coûté. La réponse se
construit poste par poste, en distinguant ce qui est constaté (factures,
pertes comptabilisées) de ce qui est estimé.

\begin{boxscenario}{Opération NYANGA. Le coût au 30 juin 2025}
\begin{tableaucours}{L{6.6cm}rL{2.6cm}}{%
    \tableauHeader{Poste} &
    \tableauHeader{M FCFA} &
    \tableauHeader{Nature}
}
Virements frauduleux & 47,3 & Constaté \\
Montants recouvrés par les opérateurs & $-12{,}6$ & Constaté \\
Mission d'enquête et d'expertise privée & 21,0 & Constaté (factures) \\
Remplacement de la passerelle et mesures immédiates & 34,0 & Constaté \\
Arrêt partiel des agences (estimation des pertes d'activité) & 28,3 &
Estimé \\
\textbf{Total} & \textbf{118,0} & \\
Notification et accompagnement des clients exposés & Non chiffré & \NEK{} \\
Contentieux éventuels & Non chiffré & \NEK{} \\
\end{tableaucours}

\medskip
\textbf{Analyse} (confiance élevée sur les postes constatés, moyenne sur
l'arrêt des agences). Le coût connu, 118~M~FCFA, est une borne basse. Les
deux postes \NEK{} peuvent le dépasser. À titre de comparaison, la mise à
jour de la passerelle en décembre 2024 aurait demandé une intervention de
deux heures du prestataire.
\end{boxscenario}

% ============================================================================
\section{L'Exercice Intégrateur : SOLITRANS, Février 2026}
\label{sec:c14-solitrans}
% ============================================================================

Le portail des transitaires de SOLITRANS a été compromis en février 2026
(Chapitre~\ref{chap:antiforensique}). Le laboratoire FAV-LAB-10 de NDAP
SUKULU en fournit les pièces~: journaux d'audit de la console du
fournisseur, journaux de l'orchestrateur, journaux de l'application envoyés
au collecteur, instantané des disques, historique des balayages du portail,
tickets de correction. Le lecteur y applique les dix étapes. Trois
questions guident l'exercice~:
\begin{enumerate}
    \item Par quelle faille l'attaquant est-il entré, et cette faille
          était-elle connue du programme de gestion des vulnérabilités~?
    \item Qu'a-t-il modifié, et qu'a-t-il pu lire~?
    \item Le programme de SOLITRANS a-t-il échoué, ou l'attaque est-elle
          passée par une faille hors de sa portée~?
\end{enumerate}

\begin{boxinfo}
\textbf{Une question piège.} La troisième question n'a pas de réponse
évidente. Une faille connue, classée, dont l'exception E-03 était en cours,
n'est pas un échec du programme au même titre qu'une faille ignorée~; mais
une mesure compensatoire défaillante en est un. L'exercice demande de
coter, pas de défendre.
\end{boxinfo}

% ============================================================================
\section{Ce que l'Enquête demande à chaque Métier}
\label{sec:c14-metiers}
% ============================================================================

\begin{boxmetier}{Analyste en réponse à incident}
\textbf{Relire son dossier avant de le remettre}
\begin{enumerate}
    \item Chaque conclusion renvoie-t-elle à des pièces cotées~?
    \item Chaque pièce a-t-elle sa chaîne de possession~?
    \item Chaque outil a-t-il sa version et sa sortie brute en annexe~?
    \item Chaque heure a-t-elle son fuseau et sa dérive~?
    \item Chaque \NEK{} est-il dit, avec ce qui permettrait de le lever~?
\end{enumerate}
\end{boxmetier}

\begin{boxmetier}{RSSI}
Présenter à la direction le coût constaté, le coût estimé et ce qui reste
inconnu, puis le coût de la mesure qui aurait évité l'incident. La
comparaison fait plus pour le budget de la sécurité que toute projection.
\end{boxmetier}

\begin{boxmetier}{Analyste en gestion des vulnérabilités}
Recevoir le dossier d'enquête et répondre par écrit à une question~: le
programme aurait-il dû voir cette faille, et pourquoi ne l'a-t-il pas vue~?
\end{boxmetier}

\begin{boxmetier}{Application OPJ/Expert Judiciaire}
L'expert relit son rapport en se posant la question de l'expert adverse~:
quelle conclusion pourrais-je attaquer, et sur quelle pièce~?
\end{boxmetier}

\separateur

\begin{boxresume}
\textbf{Ce qu'il faut retenir de ce chapitre~:}
\begin{itemize}
    \item Une enquête suit dix étapes, chacune produisant une pièce du
          dossier~; les étapes~7 et~8 la relient à la gestion des
          vulnérabilités.
    \item Réponse à incident, enquête privée et enquête judiciaire
          partagent un registre et des scellés, pas un rapport.
    \item Le coût d'un incident se construit poste par poste, constaté ou
          estimé~; ce qui reste \NEK{} se dit.
    \item À NYANGA, trois des cinq questions sont tranchées, l'identité des
          auteurs reste \NEK{}, et le coût connu atteint 118~M~FCFA.
\end{itemize}
\end{boxresume}

\begin{boxinfo}
\textbf{Points de contrôle.} Vous avez compris ce chapitre si vous savez~:
\begin{itemize}
    \item citer les dix étapes et la pièce que chacune produit~;
    \item organiser la coordination de trois enquêtes sur une affaire~;
    \item présenter un coût d'incident en distinguant constaté, estimé et
          inconnu.
\end{itemize}
\end{boxinfo}

% ============================================================================
\section*{Exercices}
% ============================================================================

\begin{boxexo}[title={\ding{45}\quad Exercice~14.1. Ordonner \niveaubase}]
Remettez dans l'ordre de la méthode~: (a) construire la chronologie~;
(b) formuler les hypothèses~; (c) isoler~; (d) chercher les variantes~;
(e) acquérir~; (f) remonter à la vulnérabilité~; (g) cadrer le mandat~;
(h) examiner~; (i) rendre compte~; (j) clore.
\end{boxexo}

\begin{boxexo}[title={\ding{45}\quad Exercice~14.2. Relecture contradictoire \niveaumoyen}]
Jouez l'expert adverse. Choisissez trois conclusions du dossier NYANGA et,
pour chacune, formulez l'objection la plus forte possible, puis la réponse
que le dossier permet d'y apporter. Indiquez si l'une d'elles devrait
changer de cote.
\end{boxexo}

\begin{boxexo}[title={\ding{45}\quad Exercice~14.3. Le coût \niveaumoyen}]
\begin{enumerate}[(a)]
    \item Vérifiez le total de l'encadré « Le coût au 30 juin 2025 ».
    \item Quelle part du coût connu est estimée plutôt que constatée~?
    \item Rédigez trois phrases pour présenter ce coût au conseil
          d'administration.
\end{enumerate}
\end{boxexo}

\begin{boxexo}[title={\ding{45}\quad Exercice~14.4. SOLITRANS, février 2026 \niveauexpert}]
À partir des pièces du laboratoire FAV-LAB-10, conduisez les dix étapes et
répondez aux trois questions de la section~14.5. Pour chaque réponse,
donnez une cote et les pièces qui la fondent. Rédigez enfin la note d'une
page destinée à M.~NJIKE.
\end{boxexo}

\bigskip

\begin{boxinfo}
\textbf{Pour aller plus loin.}
\begin{itemize}
    \item L'ISO/IEC~27043~\cite{iso27043_2015} et l'ISO/IEC~27035~\cite{iso27035}.
    \item Le NIST SP~800-61~\cite{nist80061r2}.
\end{itemize}
\end{boxinfo}

\begin{boxinfo}
\textbf{Transition.} Le dossier est complet. Il reste à l'écrire pour ceux
qui ne l'ont pas vu se construire, et à le défendre devant ceux qui le
contestent~: Chapitre~\ref{chap:rapport}.
\end{boxinfo}

        % ============================================================================
% Fichier  : partie5/P05_C15_rapport_temoignage.tex
% Titre    : Le Rapport et le Témoignage
% Auteur   : MINKA MI NGUIDJOÏ Thierry Emmanuel
% Version  : 1.0 -- novembre 2026
% Corrigés : corriges/Corr_C15.tex
% ============================================================================
% Sources : ISO/IEC 27042 (interprétation et rapport, paraphrasée) ; ISO/IEC
% 17025 (compétence des laboratoires) ; loi n° 2010/012, art. 56 (rappel).
% Faits nouveaux (bible) : rapport d'expertise de M. ABOUBAKAR déposé le
% 28/05/2025, 46 pages et 9 annexes ; audience du 15/02/2026 devant le
% tribunal de grande instance de la Mifi (Bafoussam) : deux titulaires de
% portefeuilles poursuivis pour recel ; l'expert n'attribue pas l'intrusion.
% ============================================================================

\chapter{Le Rapport et le Témoignage}
\label{chap:rapport}
\chaptermark{Le rapport et le témoignage}

\epigraph{%
    « Un rapport que seul son auteur comprend\\
    n'a convaincu personne. »%
}{Principe du cours}

\niveauChapitre{Intermédiaire}{%
    Maîtrise du Chapitre~\ref{chap:enquete}
}

\begin{boxobjectifs}
\begin{itemize}
    \item Adapter un rapport à son destinataire~: direction, autorité
          judiciaire, équipe technique.
    \item Structurer un rapport d'expertise de façon qu'il soit vérifiable.
    \item Écrire chaque conclusion avec sa constatation, son analyse et son
          niveau de confiance.
    \item Dire les limites de l'examen sans les cacher ni s'y réfugier.
    \item Préparer une déposition et répondre à une question hostile.
\end{itemize}
\end{boxobjectifs}

\bigskip

% ============================================================================
\section*{Avant de commencer : une question simple}
% ============================================================================

Bafoussam, tribunal de grande instance de la Mifi, 15 février 2026. Deux
titulaires de portefeuilles de paiement mobile, qui ont reçu une partie des
fonds détournés, sont poursuivis pour recel. M.~ABOUBAKAR, l'expert requis
par le procureur, dépose à la barre. L'avocat de la défense lui pose une
question simple~: « Monsieur l'expert, pouvez-vous dire au tribunal que mes
clients ont piraté NYANGA CRÉDIT~? » La réponse tient en une phrase, et
l'expert l'a préparée.

\begin{boxalert}
\textbf{Violation typique.} Laisser un rapport, ou une déposition, aller
au-delà de ce que les pièces établissent, parce que le destinataire attend
une réponse ferme.

\smallskip
\textbf{Règle~:} l'expert répond à la question posée, avec ce que les pièces
établissent, et dit clairement ce qu'elles n'établissent pas. Une réponse
\NEK{} bien formulée sert la justice mieux qu'une certitude fragile.
\end{boxalert}

\separateur

% ============================================================================
\section{Trois Destinataires, Trois Rapports}
\label{sec:c15-destinataires}
% ============================================================================

\begin{tableaucours}{L{2.6cm}L{3.4cm}L{3.2cm}L{3.0cm}}{%
    \tableauHeader{Destinataire} &
    \tableauHeader{Ce qu'il doit décider} &
    \tableauHeader{Forme} &
    \tableauHeader{Écueil}
}
Direction & Fermer la porte, financer, informer & Note de deux pages,
recommandations datées & Trop de technique \\
Autorité judiciaire & Établir des faits, qualifier & Rapport d'expertise
complet, annexes & Sortir de la mission \\
Équipe technique & Corriger, détecter & Rapport détaillé, indicateurs,
règles & Pas de priorités \\
\end{tableaucours}
\vspace{-0.8em}
\begin{center}{\small\itshape Tableau~15.1. Trois rapports pour une même enquête}\end{center}

Les trois rapports reposent sur les mêmes constatations. Ils diffèrent par
ce qu'ils retiennent et par leur langue. La note à la direction ne contient
pas de vecteur CVSS~; le rapport d'expertise ne contient pas de
recommandation budgétaire~; le rapport technique contient les deux.

% ============================================================================
\section{Le Rapport d'Expertise}
\label{sec:c15-expertise}
% ============================================================================

L'ISO/IEC~27042~\cite{iso27042_2015} demande que le rapport permette à un
tiers compétent de comprendre ce qui a été fait, d'en évaluer la validité et
de le refaire (paraphrase). Le plan suivant répond à cette exigence.

\begin{tableaucours}{cL{3.6cm}L{8.0cm}}{%
    \tableauHeader{} &
    \tableauHeader{Partie} &
    \tableauHeader{Contenu}
}
1 & Mission & Autorité requérante, fondement (art.~56 de la loi
n°~2010/012), questions posées, mot pour mot \\
2 & Pièces reçues & Scellés, état à réception, empreintes vérifiées \\
3 & Opérations & Chaque opération, datée, avec outils, versions et copies
utilisées \\
4 & Constatations & Les faits, sans interprétation, avec leur source \\
5 & Discussion & Les interprétations, les autres explications examinées,
le niveau de confiance \\
6 & Limites & Ce qui n'a pas pu être examiné, et pourquoi \\
7 & Réponses aux questions & Une réponse par question, cotée \\
8 & Annexes & Sorties brutes, tableaux, chronologie, formulaires \\
\end{tableaucours}
\vspace{-0.8em}
\begin{center}{\small\itshape Tableau~15.2. Le plan d'un rapport d'expertise}\end{center}

\subsection{Écrire une conclusion}

Chaque conclusion du rapport suit le même schéma~: la constatation, avec sa
source~; l'analyse, avec les autres explications examinées~; le niveau de
confiance, et ce qui le fonde. Le vocabulaire de la confiance doit être
défini dans le rapport et tenu partout.

\begin{tableaucours}{L{2.8cm}L{9.4cm}}{%
    \tableauHeader{Niveau} &
    \tableauHeader{Usage dans ce livre}
}
Élevée & Plusieurs pièces indépendantes concordent, aucune explication
concurrente plausible ne subsiste \\
Moyenne & Les pièces concordent mais l'une est fragile, ou une explication
concurrente reste possible \\
Faible & Une seule pièce, ou des explications concurrentes aussi plausibles \\
\end{tableaucours}
\vspace{-0.8em}
\begin{center}{\small\itshape Tableau~15.3. Les niveaux de confiance}\end{center}

\begin{boxscenario}{Opération NYANGA. Une réponse du rapport d'expertise}
\textit{Question~3 de la mission~: « Dire par quel moyen l'accès au système
d'information de NYANGA CRÉDIT a été obtenu. »}

\medskip
\textbf{Constatation.} La passerelle d'accès à distance fonctionnait le
2~mars 2025 avec le micrologiciel 7.2.1 (scellé NC-02, export de
configuration). Le bulletin de l'éditeur du 12~décembre 2024 désigne cette
version comme vulnérable à un contournement de l'authentification. Le
journal de la passerelle (NC-02) enregistre, le 2~mars entre 23~h~46 et
23~h~48, quatre requêtes depuis l'adresse 203.0.113.18 conformes aux
indicateurs publiés par l'éditeur, suivies de l'ouverture d'une session du
compte \texttt{svc\_backup} sans authentification préalable.

\textbf{Discussion.} Deux autres explications ont été examinées~:
l'utilisation d'un mot de passe connu, qui n'explique pas la session sans
authentification~; le vol d'une session existante, écarté faute de session
légitime dans les neuf jours précédents.

\textbf{Réponse} (confiance élevée). L'accès a été obtenu par l'exploitation
d'une vulnérabilité de la passerelle permettant de contourner
l'authentification.
\end{boxscenario}

\subsection{Dire les limites}

La partie consacrée aux limites protège l'expert autant que le lecteur.
Elle énumère ce que l'examen n'a pas couvert~: la mémoire de SRV-CPT01,
jamais capturée~; le système de la passerelle, illisible sans l'éditeur~;
la période du 3 au 8~mars, sans trace~; la provenance des journaux de
l'interface de paiement. Pour chaque limite, elle dit l'effet sur les
conclusions. Une limite non écrite devient, à l'audience, une faille de
l'expertise.

% ============================================================================
\section{L'Indépendance et la Compétence}
\label{sec:c15-independance}
% ============================================================================

L'expert requis tient sa légitimité de deux qualités~: son indépendance à
l'égard des parties, et sa compétence. La première exclut d'avoir travaillé
pour l'une d'elles sur l'affaire. C'est pourquoi le procureur n'a pas requis
Mme~MBARGA, mandatée et payée par NYANGA CRÉDIT, mais M.~ABOUBAKAR. La
seconde se démontre par la formation, l'expérience et la méthode~; pour un
laboratoire, l'accréditation selon l'ISO/IEC~17025~\cite{iso17025_2017}
atteste de la compétence technique et de la maîtrise des méthodes.

\begin{boxinfo}
\textbf{L'enquêteur privé comme témoin.} Mme~MBARGA peut être entendue comme
témoin des opérations qu'elle a réalisées~: acquisitions, chaîne de
possession, premières constatations. Elle ne se substitue pas à l'expert,
et l'expert ne se contente pas de reprendre ses conclusions~: il refait ou
contrôle ce qui fonde les siennes.
\end{boxinfo}

% ============================================================================
\section{Déposer à l'Audience}
\label{sec:c15-audience}
% ============================================================================

L'audience n'est pas un séminaire technique. Le tribunal, les avocats et
les prévenus n'ont pas lu les annexes. L'expert y gagne à préparer trois
choses~: une chronologie d'une page, lisible à distance~; une définition en
une phrase de chaque terme technique qu'il emploiera~; et la liste des
questions hostiles probables, avec leur réponse.

\begin{tableaucours}{L{5.4cm}L{6.8cm}}{%
    \tableauHeader{Question hostile} &
    \tableauHeader{Réponse fondée sur le dossier}
}
« Vos outils ne sont-ils pas faillibles~? » & Les conclusions principales
reposent sur deux sources indépendantes, lues avec des outils dont les
versions et les sorties sont en annexe \\
« Le disque est resté une nuit sur un bureau. » & L'intégrité est établie
par les empreintes~; la conclusion qu'il soutient est confirmée par une
autre pièce, à la chaîne complète \\
« Pouvez-vous exclure une erreur d'horloge~? » & Les dérives ont été
mesurées et corrigées~; les écarts restants sont inférieurs à la minute \\
« Mes clients ont-ils piraté l'établissement~? » & Les pièces ne permettent
pas de le dire (voir ci-dessous) \\
\end{tableaucours}
\vspace{-0.8em}
\begin{center}{\small\itshape Tableau~15.4. Préparer les questions hostiles}\end{center}

\begin{boxscenario}{Opération NYANGA. La réponse de l'expert}
\textbf{Réponse de M.~ABOUBAKAR.} « Non, Maître. Les pièces que j'ai
examinées établissent comment le système a été pénétré et comment les fonds
ont été transférés. Elles ne permettent pas d'identifier les personnes qui
ont commis l'intrusion. Elles établissent, par les relevés de l'opérateur,
que deux des onze portefeuilles qui ont reçu les fonds sont ouverts au nom
de vos clients. Ce que vos clients savaient de l'origine de ces fonds n'est
pas une question technique. »

\medskip
\textbf{Analyse.} L'expert répond à la question, sépare ce qui est établi de
ce qui ne l'est pas, et laisse au tribunal ce qui relève du tribunal.
L'avocat n'obtient ni une attribution qu'il pourrait contester, ni un refus
de répondre qu'il pourrait exploiter.
\end{boxscenario}

\begin{boxlecture}{Les réponses de l'expertise}
\begin{tableaucours}{L{7.6cm}c}{%
    \tableauHeader{Question de la mission} &
    \tableauHeader{Réponse}
}
Moyen de l'accès & \OUI{} \\
Séquence des opérations de l'attaquant & \OUI{} (sauf 3 au 8~mars, \NEK{}) \\
Données copiées & \OUI{} \\
Destination des fonds & \OUI{} (onze portefeuilles) \\
Identité des auteurs de l'intrusion & \NEK{} \\
\end{tableaucours}

\medskip
\textbf{Lecture.} Quatre réponses établies, une non tranchée. Le \NEK{} sur
les auteurs n'affaiblit pas l'expertise~; il en marque l'honnêteté. Les
poursuites pour recel reposent sur d'autres éléments, que la justice réunit
par d'autres moyens.
\end{boxlecture}

% ============================================================================
\section{Ce que le Rapport demande à chaque Métier}
\label{sec:c15-metiers}
% ============================================================================

\begin{boxmetier}{Application OPJ/Expert Judiciaire}
\textbf{Rédiger et défendre une expertise}
\begin{enumerate}
    \item Recopier les questions de la mission, mot pour mot.
    \item Séparer, pour chaque conclusion, constatation, discussion et
          réponse.
    \item Définir les niveaux de confiance et les tenir.
    \item Écrire les limites et leur effet.
    \item Préparer la chronologie d'une page et les questions hostiles.
\end{enumerate}
\end{boxmetier}

\begin{boxmetier}{Analyste en réponse à incident}
Écrire le rapport technique pour qu'il puisse devenir une pièce de
procédure~: chaque affirmation sourcée, chaque opération datée. Il le
deviendra plus souvent qu'on ne le prévoit.
\end{boxmetier}

\begin{boxmetier}{RSSI}
Lire les limites du rapport avant ses conclusions~: elles disent ce que
l'organisation doit encore préparer pour la prochaine enquête.
\end{boxmetier}

\separateur

\begin{boxresume}
\textbf{Ce qu'il faut retenir de ce chapitre~:}
\begin{itemize}
    \item Une même enquête donne trois rapports, pour la direction, la
          justice et l'équipe technique.
    \item Le rapport d'expertise recopie la mission, décrit les opérations,
          sépare constatations et discussion, écrit ses limites et répond
          question par question.
    \item Les niveaux de confiance se définissent et se tiennent.
    \item L'indépendance exclut d'avoir travaillé pour une partie~; la
          compétence se démontre.
    \item À l'audience, l'expert répond à la question posée et laisse au
          tribunal ce qui lui revient.
\end{itemize}
\end{boxresume}

\begin{boxinfo}
\textbf{Points de contrôle.} Vous avez compris ce chapitre si vous savez~:
\begin{itemize}
    \item rédiger une réponse d'expertise en trois temps~;
    \item écrire une section de limites et dire l'effet de chacune~;
    \item répondre à une question hostile sans dépasser les pièces.
\end{itemize}
\end{boxinfo}

% ============================================================================
\section*{Exercices}
% ============================================================================

\begin{boxexo}[title={\ding{45}\quad Exercice~15.1. Quel rapport~? \niveaubase}]
Dites dans quel rapport (direction, justice, technique) figure chacun des
éléments suivants~: (a) le vecteur CVSS de VULN-2024-4471~; (b) le coût de
l'incident~; (c) la sortie brute de \texttt{windows.netscan}~; (d) la
recommandation sur la conservation des journaux~; (e) les questions de la
mission.
\end{boxexo}

\begin{boxexo}[title={\ding{45}\quad Exercice~15.2. Réécrire une conclusion \niveaumoyen}]
Réécrivez en trois temps (constatation, discussion, réponse cotée) la phrase
suivante, tirée d'un brouillon de rapport~: « Il est évident que les
pirates ont volé le mot de passe de M.~KAMGA pour vider les comptes. »
\end{boxexo}

\begin{boxexo}[title={\ding{45}\quad Exercice~15.3. La section des limites \niveaumoyen}]
Rédigez la section « Limites » du rapport d'expertise de NYANGA. Pour chaque
limite, indiquez la ou les conclusions affectées et la manière dont elles le
sont.
\end{boxexo}

\begin{boxexo}[title={\ding{45}\quad Exercice~15.4. Trois questions hostiles \niveauexpert}]
Préparez la réponse de l'expert aux trois questions suivantes, en vous
appuyant sur le dossier~: (a) « Pourquoi n'avez-vous pas examiné la
passerelle elle-même~? »~; (b) « Les journaux de l'interface de paiement
vous ont été transmis par courriel par une société privée~: comment savoir
qu'ils sont authentiques~? »~; (c) « Votre rapport reprend-il simplement
celui du cabinet payé par la victime~? »
\end{boxexo}

\bigskip

\begin{boxinfo}
\textbf{Pour aller plus loin.}
\begin{itemize}
    \item L'ISO/IEC~27042~\cite{iso27042_2015}, sur l'interprétation et la
          présentation des résultats.
    \item L'ISO/IEC~17025~\cite{iso17025_2017}, pour la compétence des
          laboratoires.
\end{itemize}
\end{boxinfo}

\begin{boxinfo}
\textbf{Transition.} Le rapport dit ce que l'enquête sait et ce qu'elle ne
sait pas. Peut-on mesurer cette part d'ignorance, la comparer d'une enquête
à l'autre, la réduire méthodiquement~? C'est la question de recherche du
Chapitre~\ref{chap:auditologie}.
\end{boxinfo}


    % PARTIE 6 · RECHERCHE
    \part{Recherche : Mesurer l'Incertitude d'une Enquête}
    \label{part:recherche}

        % ============================================================================
% Fichier  : partie6/P06_C16_mesurer_incertitude.tex
% Titre    : Mesurer l'Incertitude d'une Enquête
% Auteur   : MINKA MI NGUIDJOÏ Thierry Emmanuel
% Version  : 1.0 -- novembre 2026
% Corrigés : corriges/Corr_C16.tex
% ============================================================================
% Sources : Minka (2026), A0 « Vers l'auditologie » (HAL hal-05701007) et A1
% sur l'instrument trichotomique (HAL hal-05701270) ; Shannon (1948).
% Les articles A2 à A5 du programme sont en préparation et ne sont pas cités.
% Mesures reprises du corpus (NRAC, ACSI) : mu = part de OUI, U = part de
% NEK, H = entropie de Shannon de la répartition OUI/NON/NEK, en bits.
% Faits nouveaux (bible) : douze propositions d'enquête P1 à P12 cotées au
% 14/03, au 29/03 et au 06/04/2025 (0/0/12 ; 4/2/6 ; 5/2/5).
% ============================================================================

\chapter{Mesurer l'Incertitude d'une Enquête}
\label{chap:auditologie}
\chaptermark{Mesurer l'incertitude}

\epigraph{%
    « Ce que l'on ne sait pas se compte,\\
    comme ce que l'on sait. »%
}{Principe du cours}

\niveauChapitre{Recherche}{%
    Maîtrise de l'ensemble du cours
}

\begin{boxobjectifs}
\begin{itemize}
    \item Représenter l'état d'une enquête par un ensemble de propositions
          cotées \OUI{}, \NON{} ou \NEK{}.
    \item Calculer et interpréter les mesures $\mu$, $U$ et $H$.
    \item Suivre l'évolution de l'incertitude au cours d'une enquête.
    \item Distinguer un \NEK{} réductible d'un \NEK{} irréductible dans le
          cadre d'une mission.
    \item Formuler une question de recherche sur l'incertitude en
          investigation numérique.
\end{itemize}
\end{boxobjectifs}

\bigskip

% ============================================================================
\section*{Avant de commencer : « avez-vous avancé~? »}
% ============================================================================

Bafoussam, 29 mars 2025. Le directeur général de NYANGA CRÉDIT demande à
Mme~MBARGA où en est l'enquête. Elle pourrait répondre par une liste de
découvertes. Elle répond par trois nombres~: sur les douze propositions que
l'enquête doit trancher, quatre sont établies, deux réfutées, six encore
ouvertes. Le 14~mars au soir, les douze étaient ouvertes. Le directeur
comprend ce que l'enquête sait, et ce qu'il lui reste à apprendre.

\begin{boxalert}
\textbf{Violation typique.} Rendre compte de l'avancement d'une enquête par
ce qui a été trouvé, sans dire ce qui reste inconnu ni si cet inconnu
diminue.

\smallskip
\textbf{Règle~:} l'état d'une enquête se décrit par l'ensemble des
propositions qu'elle doit trancher, chacune cotée. L'avancement se mesure à
la baisse de la part des \NEK{}, pas au nombre de découvertes.
\end{boxalert}

\separateur

% ============================================================================
\section{L'Instrument Trichotomique}
\label{sec:c16-instrument}
% ============================================================================

Tout au long du cours, chaque conclusion a reçu une cote \OUI{}, \NON{} ou
\NEK{}. Cet instrument vient du programme de recherche de
l'auditologie~\cite{minka2026auditologie}, qui propose de traiter l'audit
des systèmes d'information, et plus largement toute démarche qui conclut
sur des preuves, comme une science appliquée de la réduction de
l'incertitude. L'article qui formalise l'instrument~\cite{minka2026nek}
montre qu'une réponse binaire, oui ou non, force l'auditeur à transformer
son ignorance en affirmation, et qu'une troisième réponse, \NEK{}, rend
cette ignorance visible et mesurable.

\begin{boxdef}
Soit un ensemble de $n$ propositions, cotées, avec $n_{O}$ réponses \OUI{},
$n_{N}$ réponses \NON{} et $n_{K}$ réponses \NEK{}. On définit~:

\smallskip
\textbf{\termecle{Taux d'établissement}}~: $\mu = n_{O}/n$, la part des
propositions établies.

\smallskip
\textbf{\termecle{Taux d'incertitude}}~: $U = n_{K}/n$, la part des
propositions que les preuves ne permettent pas de trancher.

\smallskip
\textbf{\termecle{Entropie}}~: $H = -\sum_{c} p_{c} \log_{2} p_{c}$, où
$p_{c}$ est la part de chaque cote, exprimée en bits~\cite{shannon1948}.
Elle mesure la dispersion des réponses~: nulle si toutes les cotes sont
identiques, maximale ($\log_{2} 3 \approx 1{,}585$~bit) si elles sont
également réparties.
\end{boxdef}

Dans un audit, \OUI{} signifie qu'un contrôle est en place et \NON{} qu'il
ne l'est pas~; $\mu$ y mesure la conformité. Dans une enquête, \OUI{} et
\NON{} sont deux formes de savoir~: une proposition réfutée est aussi
tranchée qu'une proposition établie. La part de savoir y vaut donc $1 - U$,
et $U$ devient la mesure principale.

% ============================================================================
\section{Les Douze Propositions de NYANGA}
\label{sec:c16-nyanga}
% ============================================================================

\begin{boxscenario}{Opération NYANGA. Trois instantanés de l'enquête}
\begin{tableaucours}{cL{6.2cm}ccc}{%
    \tableauHeader{N°} &
    \tableauHeader{Proposition} &
    \tableauHeader{14/03} &
    \tableauHeader{29/03} &
    \tableauHeader{06/04}
}
P1 & L'attaquant est entré par la passerelle & \NEK{} & \OUI{} & \OUI{} \\
P2 & Par l'exploitation de VULN-2024-4471 & \NEK{} & \NEK{} & \OUI{} \\
P3 & La clé de l'interface a été lue sur SRV-CPT01 & \NEK{} & \OUI{} &
\OUI{} \\
P4 & Les identifiants de \texttt{adm.kamga} ont été extraits de la copie
du 9~mars & \NEK{} & \NEK{} & \NEK{} \\
P5 & Les données sorties sont le dossier \texttt{Clients} & \NEK{} & \OUI{}
& \OUI{} \\
P6 & Les virements ont été émis hors du réseau du siège & \NEK{} & \OUI{} &
\OUI{} \\
P7 & Un employé a participé & \NEK{} & \NEK{} & \NEK{} \\
P8 & L'attaquant a agi du 3 au 8~mars & \NEK{} & \NEK{} & \NEK{} \\
P9 & Les auteurs sont identifiés & \NEK{} & \NEK{} & \NEK{} \\
P10 & La passerelle contient un code persistant & \NEK{} & \NEK{} & \NEK{} \\
P11 & Les virements sont passés par le lien du siège & \NEK{} & \NON{} &
\NON{} \\
P12 & \texttt{upd.exe} était présent sur SRV-FIC01 en 2017 & \NEK{} & \NON{} &
\NON{} \\
\end{tableaucours}
\end{boxscenario}

\begin{boxlecture}{L'incertitude de l'enquête, trois dates}
\begin{tableaucours}{L{2.0cm}cccc}{%
    \tableauHeader{Date} &
    \tableauHeader{$n_{O}/n_{N}/n_{K}$} &
    \tableauHeader{$\mu$} &
    \tableauHeader{$U$} &
    \tableauHeader{$H$ (bit)}
}
14/03 & 0/0/12 & 0 & 1 & 0 \\
29/03 & 4/2/6 & 0,333 & 0,500 & 1,459 \\
06/04 & 5/2/5 & 0,417 & 0,417 & 1,483 \\
\end{tableaucours}

\medskip
\textbf{Calcul pour le 29~mars.} $p_{O} = 4/12$, $p_{N} = 2/12$, $p_{K} =
6/12$~:
\[
H = -\left(\tfrac{1}{3}\log_{2}\tfrac{1}{3} + \tfrac{1}{6}\log_{2}\tfrac{1}{6}
+ \tfrac{1}{2}\log_{2}\tfrac{1}{2}\right) \approx 0{,}528 + 0{,}431 + 0{,}500
\approx 1{,}459~\text{bit}.
\]

\textbf{Lecture.} $U$ passe de 1 à 0,500 puis à 0,417~: l'enquête avance.
$H$, lui, augmente~: il était nul quand tout était inconnu, et la
répartition des cotes devient plus mélangée à mesure que l'enquête tranche.
$H$ ne mesure donc pas l'ignorance, mais la diversité des réponses. La
règle du corpus s'applique~: on lit $U$ d'abord, $\mu$ ensuite, et $H$ pour
contrôler que la répartition n'est pas dégénérée.

\medskip
\hfill\lecture{0{,}417}{0{,}417}{1{,}483}
\end{boxlecture}

% ============================================================================
\section{Deux Sortes de NEK}
\label{sec:c16-nek}
% ============================================================================

Les cinq \NEK{} du 6~avril ne sont pas de même nature. Certains peuvent être
levés par une pièce que l'enquête est en mesure d'obtenir~; d'autres ne le
peuvent pas dans le cadre de la mission.

\begin{tableaucours}{cL{5.2cm}L{2.4cm}L{3.8cm}}{%
    \tableauHeader{N°} &
    \tableauHeader{Ce qui lèverait le \NEK{}} &
    \tableauHeader{Nature} &
    \tableauHeader{Coût ou condition}
}
P4 & Aucune pièce~: l'extraction a eu lieu hors du système & Irréductible &
Hors de portée de toute enquête technique sur le système \\
P7 & Pièces sur les personnes (témoignages, téléphonie) & Réductible hors
mission & Enquête judiciaire \\
P8 & Journaux des agences, s'ils existent encore & Réductible & Quelques
jours d'examen \\
P9 & Données des hébergeurs de 203.0.113.18, .47 et .62 & Réductible hors
mission & Entraide internationale, mois \\
P10 & Image du système de la passerelle par l'éditeur & Réductible &
Réquisition de l'éditeur, semaines \\
\end{tableaucours}
\vspace{-0.8em}
\begin{center}{\small\itshape Tableau~16.1. Les \NEK{} de NYANGA au 6~avril}\end{center}

Cette distinction change la décision. Il est rationnel de lever P8, peu
coûteux. Il est inutile de promettre de lever P4. Lever P9 et P7 appartient
à d'autres acteurs, et le rapport doit le dire pour que la direction ne
l'attende pas de son cabinet.

\begin{boxinfo}
\textbf{La valeur d'une pièce.} Chaque pièce supplémentaire a un coût
(temps, argent, délai) et une valeur~: les \NEK{} qu'elle peut lever, et
l'importance de ces \NEK{} pour la décision. Choisir les pièces à chercher
revient à comparer ces valeurs. Formaliser cette comparaison pour
l'investigation numérique est l'une des questions ouvertes du programme de
recherche.
\end{boxinfo}

% ============================================================================
\section{Questions Ouvertes pour la Recherche}
\label{sec:c16-recherche}
% ============================================================================

L'auditologie a été proposée pour l'audit des systèmes d'information. Son
transfert à l'investigation numérique et à l'analyse de vulnérabilités
soulève des questions qui peuvent fonder un mémoire de master ou une thèse.

\begin{enumerate}
    \item \textbf{La granularité des propositions.} Les mesures dépendent du
          découpage de l'enquête en propositions. Comment définir un
          découpage stable, pour que deux enquêteurs obtiennent des $U$
          comparables sur la même affaire~?
    \item \textbf{La pondération.} Toutes les propositions ne pèsent pas
          autant pour la décision. Faut-il pondérer $U$, et selon quel
          critère qui ne soit pas arbitraire~?
    \item \textbf{La confiance.} Le cours associe à chaque analyse un niveau
          de confiance (faible, moyenne, élevée) distinct de la cote. Comment
          articuler ces deux échelles~?
    \item \textbf{Le programme de vulnérabilités.} Un actif hors inventaire
          est un \NEK{} (Chapitre~\ref{chap:identifier}). Le $U$ d'un parc,
          calculé sur la présence de failles, peut-il servir d'indicateur de
          pilotage, à côté du nombre de failles~?
    \item \textbf{L'attribution.} L'identité des auteurs reste souvent
          \NEK{}. Peut-on décrire formellement les conditions qui la rendent
          réductible~?
\end{enumerate}

\begin{boxinfo}
\textbf{Construire un sujet de mémoire sur l'incertitude}
\begin{enumerate}
    \item Choisir un objet~: une enquête, un programme de vulnérabilités,
          un type de pièce.
    \item Définir l'ensemble des propositions et la règle de cotation.
    \item Recueillir des cas réels anonymisés ou des cas construits,
          documentés.
    \item Mesurer $U$, $\mu$ et $H$, et leur évolution.
    \item Confronter les mesures à un critère extérieur (décision prise,
          erreur constatée a posteriori).
    \item Discuter les limites de la mesure, et proposer une amélioration.
\end{enumerate}
\end{boxinfo}

% ============================================================================
\section{Ce que la Mesure demande à chaque Métier}
\label{sec:c16-metiers}
% ============================================================================

\begin{boxmetier}{Analyste en réponse à incident}
Ouvrir chaque mission par la liste des propositions à trancher, la coter
chaque semaine et rendre compte de l'avancement par $U$. La liste elle-même
oblige à dire, dès le départ, ce que l'on cherche.
\end{boxmetier}

\begin{boxmetier}{Application OPJ/Expert Judiciaire}
Le tableau des propositions cotées, annexé au rapport, donne au tribunal une
vue d'ensemble honnête~: ce qui est établi, ce qui est réfuté, ce qui ne
peut pas l'être.
\end{boxmetier}

\begin{boxmetier}{Analyste en gestion des vulnérabilités}
Publier, à côté du nombre de failles, la part du parc dont on ne sait rien.
C'est le $U$ du programme.
\end{boxmetier}

\begin{boxmetier}{RSSI}
Demander à chaque enquête et à chaque programme ce qu'ils ont trouvé, puis
ce qu'ils ignorent encore et ce qu'il en coûterait de le savoir.
\end{boxmetier}

\separateur

\begin{boxresume}
\textbf{Ce qu'il faut retenir de ce chapitre~:}
\begin{itemize}
    \item L'état d'une enquête se décrit par des propositions cotées
          \OUI{}, \NON{} ou \NEK{}.
    \item $\mu$ est la part des \OUI{}, $U$ la part des \NEK{}, $H$
          l'entropie de la répartition~; dans une enquête, $U$ est la mesure
          principale.
    \item À NYANGA, $U$ passe de 1 à 0,417 en trois semaines, tandis que
          $H$ augmente~: $H$ mesure la diversité, pas l'ignorance.
    \item Un \NEK{} est réductible ou non dans le cadre d'une mission~; la
          distinction guide le choix des pièces à chercher.
    \item Le transfert de l'auditologie à l'investigation ouvre des
          questions de recherche sur la granularité, la pondération et la
          confiance.
\end{itemize}
\end{boxresume}

\begin{boxinfo}
\textbf{Points de contrôle.} Vous avez compris ce chapitre si vous savez~:
\begin{itemize}
    \item calculer $\mu$, $U$ et $H$ pour un ensemble de cotes~;
    \item expliquer pourquoi $H$ peut augmenter pendant qu'une enquête
          avance~;
    \item classer un \NEK{} comme réductible ou irréductible.
\end{itemize}
\end{boxinfo}

% ============================================================================
\section*{Exercices}
% ============================================================================

\begin{boxexo}[title={\ding{45}\quad Exercice~16.1. Calculer \niveaubase}]
Calculez $\mu$, $U$ et $H$ pour les répartitions suivantes de dix
propositions~: (a) 10/0/0~; (b) 5/5/0~; (c) 4/3/3~; (d) 0/0/10.
\end{boxexo}

\begin{boxexo}[title={\ding{45}\quad Exercice~16.2. Le 6 avril \niveaumoyen}]
\begin{enumerate}[(a)]
    \item Vérifiez le calcul de $H$ au 6~avril.
    \item Si l'examen des journaux des agences établit l'activité du 3 au
          8~mars (P8 passe à \OUI{}), que deviennent $\mu$, $U$ et $H$~?
    \item Et si l'éditeur établit l'absence de code persistant dans la
          passerelle (P10 passe à \NON{})~?
\end{enumerate}
\end{boxexo}

\begin{boxexo}[title={\ding{45}\quad Exercice~16.3. Le U d'un parc \niveaumoyen}]
Reprenez la couverture de SOLITRANS en mai 2025 (Chapitre~\ref{chap:identifier})~:
1\,512 actifs, dont 170 non analysés et 22~\% des serveurs analysés sans
authentification. Proposez une manière de calculer un $U$ du parc sur la
question « l'actif présente-t-il une faille de classe P1~? », et discutez-en
les limites.
\end{boxexo}

\begin{boxexo}[title={\ding{45}\quad Exercice~16.4. Un sujet de recherche \niveauexpert}]
Choisissez l'une des cinq questions ouvertes de la section~16.4 et rédigez
une proposition de sujet de mémoire d'une page~: question, hypothèse,
données, méthode de mesure, critère de validation, limites attendues.
\end{boxexo}

\bigskip

\begin{boxinfo}
\textbf{Pour aller plus loin.}
\begin{itemize}
    \item T.~E.~Minka Mi Nguidjoï, « Vers l'auditologie »~\cite{minka2026auditologie}.
    \item T.~E.~Minka Mi Nguidjoï, sur l'instrument
          trichotomique~\cite{minka2026nek}.
    \item C.~E.~Shannon, « A Mathematical Theory of
          Communication »~\cite{shannon1948}.
\end{itemize}
\end{boxinfo}

\begin{boxinfo}
\textbf{Fin du cours.} La nuit du 14~mars a commencé par vingt-trois
virements et une question~: par où sont-ils entrés~? Elle s'achève sur une
règle d'inventaire, cinq \NEK{} assumés et une passerelle sous scellé. Entre
les deux, l'enquête a préservé, acquis, analysé, daté, remonté à la faille
et rendu compte, en disant à chaque étape ce qu'elle savait et ce qu'elle
ignorait.
\end{boxinfo}


% --------------------------------------------------------------------------
% ANNEXES
% --------------------------------------------------------------------------
\appendix

    % ============================================================================
% Fichier  : annexes/Ann_A_glossaire.tex
% Titre    : Glossaire
% Version  : 1.0 -- novembre 2024
% ============================================================================
% Définitions reprises des chapitres, abrégées ; le chapitre de référence est
% indiqué.
% ============================================================================

\chapter{Glossaire}
\label{ann:glossaire}

Ce glossaire reprend, sous une forme abrégée, les définitions données dans
les chapitres. Le numéro renvoie au chapitre où le terme est défini et
illustré.

\begin{description}
    \item[Acquisition à chaud] Recueil sur une machine allumée (mémoire,
          volume déverrouillé), qui laisse une empreinte documentée de
          l'outil (Ch.~\ref{chap:acquisition}).
    \item[Acquisition à froid] Copie d'un support après extinction, derrière
          un bloqueur d'écriture (Ch.~\ref{chap:acquisition}).
    \item[Analyse] Interprétation d'une ou de plusieurs constatations,
          assortie d'un niveau de confiance (Ch.~\ref{chap:prologue}).
    \item[Analyse de variantes] Recherche, dans tout le parc, de la même
          vulnérabilité, de la même faiblesse ou de la même cause profonde
          qu'un incident a révélées (Ch.~\ref{chap:artefact-vuln}).
    \item[Anti-forensique] Techniques visant à empêcher, retarder ou tromper
          l'investigation (Ch.~\ref{chap:antiforensique}).
    \item[Artefact] Trace qu'un système laisse en fonctionnant normalement~:
          entrée de journal, clé de registre, raccourci (Ch.~\ref{chap:disque}).
    \item[Autorité de numérotation (CNA)] Organisme habilité à attribuer des
          identifiants CVE (Ch.~\ref{chap:vulnerabilite}).
    \item[Balise] Connexion régulière d'un programme malveillant vers son
          serveur de commande (Ch.~\ref{chap:reseau}).
    \item[Bloqueur d'écriture] Dispositif qui laisse passer les lectures
          vers un support et refuse les écritures (Ch.~\ref{chap:acquisition}).
    \item[Cause contributive] Condition qui a aggravé un incident ou facilité
          l'une de ses étapes (Ch.~\ref{chap:artefact-vuln}).
    \item[Cause immédiate] Événement qui a directement permis l'incident
          (Ch.~\ref{chap:artefact-vuln}).
    \item[Cause profonde] Défaut d'organisation qui a laissé subsister la
          cause immédiate et les causes contributives
          (Ch.~\ref{chap:artefact-vuln}).
    \item[Chaîne de possession] Documentation continue de chaque détenteur,
          transfert, lieu et opération d'une pièce (Ch.~\ref{chap:preuve}).
    \item[Constatation] Fait vérifiable, rattaché à une source identifiée
          (Ch.~\ref{chap:prologue}).
    \item[Conversation] Échange entre une paire d'adresses et de ports, unité
          d'analyse d'une capture réseau (Ch.~\ref{chap:reseau}).
    \item[Cote OUI, NON, NEK] Réponse à une proposition fermée~: établie,
          contredite, ou non tranchable avec les preuves disponibles
          (Ch.~\ref{chap:prologue}, \ref{chap:auditologie}).
    \item[CVE] Identifiant unique d'une vulnérabilité publiée
          (Ch.~\ref{chap:vulnerabilite}).
    \item[CVSS] Système de notation de la gravité des vulnérabilités, de 0 à
          10~; version~4.0 avec groupes Base, Menace, Environnement,
          Supplémentaire (Ch.~\ref{chap:vulnerabilite}).
    \item[CWE] Catalogue des types de faiblesses logicielles
          (Ch.~\ref{chap:vulnerabilite}).
    \item[Empreinte cryptographique] Valeur de taille fixe calculée par une
          fonction de hachage, qui change si un seul bit de la donnée change
          (Ch.~\ref{chap:preuve}).
    \item[Entropie (H)] Mesure, en bits, de la dispersion des cotes d'un
          ensemble de propositions (Ch.~\ref{chap:auditologie}).
    \item[EPSS] Probabilité estimée qu'une CVE soit exploitée dans les trente
          jours (Ch.~\ref{chap:vulnerabilite}).
    \item[Exception] Décision écrite acceptant, pour une durée limitée et avec
          une mesure compensatoire, qu'une vulnérabilité reste ouverte
          (Ch.~\ref{chap:gerer}).
    \item[Faiblesse] Type d'erreur de conception ou de programmation
          susceptible de créer des vulnérabilités (Ch.~\ref{chap:vulnerabilite}).
    \item[Faux négatif] Vulnérabilité présente et non signalée
          (Ch.~\ref{chap:identifier}).
    \item[Faux positif] Vulnérabilité signalée et absente
          (Ch.~\ref{chap:identifier}).
    \item[Force probante] Poids qu'un décideur accorde à une pièce recevable
          (Ch.~\ref{chap:preuve}).
    \item[Indicateur de compromission] Trace caractéristique d'une attaque
          ou de l'exploitation d'une faille, publiée pour en permettre la
          détection (Ch.~\ref{chap:artefact-vuln}).
    \item[Investigation numérique] Démarche méthodique qui identifie,
          recueille, préserve, examine et analyse des traces numériques pour
          répondre à des questions de fait (Ch.~\ref{chap:investigation}).
    \item[KEV] Catalogue des vulnérabilités dont l'exploitation est constatée,
          tenu par la CISA (Ch.~\ref{chap:vulnerabilite}).
    \item[MFT] Table principale de NTFS, qui décrit chaque fichier par un
          enregistrement (Ch.~\ref{chap:disque}).
    \item[Menace] Source potentielle d'atteinte (Ch.~\ref{chap:vulnerabilite}).
    \item[Modèle de responsabilité partagée] Répartition de la sécurité entre
          le fournisseur et le client d'un service en nuage
          (Ch.~\ref{chap:antiforensique}).
    \item[Pivot] Élément commun qui relie des événements dans une
          chronologie~: adresse, compte, empreinte, heure
          (Ch.~\ref{chap:chronologie}).
    \item[Preuve numérique] Trace recueillie, préservée et examinée de façon à
          soutenir ou réfuter une proposition de fait devant un décideur
          (Ch.~\ref{chap:investigation}).
    \item[Recevabilité] Aptitude d'une pièce à être versée au débat
          (Ch.~\ref{chap:preuve}).
    \item[Répétabilité] Propriété d'un examen qui donne le même résultat,
          refait par la même personne avec les mêmes outils
          (Ch.~\ref{chap:investigation}).
    \item[Reproductibilité] Propriété d'un examen qui donne un résultat
          compatible, refait par une autre personne avec d'autres outils
          (Ch.~\ref{chap:investigation}).
    \item[Risque] Combinaison de la vraisemblance qu'une menace exploite une
          vulnérabilité et de la gravité de l'atteinte
          (Ch.~\ref{chap:vulnerabilite}).
    \item[SSVC] Méthode de décision de la CISA, qui conduit à suivre, suivre
          de près, traiter ou agir (Ch.~\ref{chap:vulnerabilite}).
    \item[Système de fichiers] Organisation des données d'un volume en
          fichiers et répertoires (Ch.~\ref{chap:disque}).
    \item[Table de partitions] Structure qui découpe un disque en volumes
          (MBR, GPT) (Ch.~\ref{chap:disque}).
    \item[Taux d'établissement ($\mu$)] Part des propositions cotées \OUI{}
          (Ch.~\ref{chap:auditologie}).
    \item[Taux d'incertitude ($U$)] Part des propositions cotées \NEK{}
          (Ch.~\ref{chap:auditologie}).
    \item[Temps de présence] Durée entre la première action établie de
          l'attaquant et la détection~; borne inférieure
          (Ch.~\ref{chap:chronologie}).
    \item[Trace numérique] Donnée qui garde la marque d'une action ou d'un
          état passé (Ch.~\ref{chap:investigation}).
    \item[Vulnérabilité] Défaut d'une version précise d'un produit, exploitable
          pour porter atteinte à la sécurité (Ch.~\ref{chap:vulnerabilite}).
\end{description}

    % ============================================================================
% Fichier  : annexes/Ann_B_outils.tex
% Titre    : Référentiel des Outils
% Version  : 1.0 -- novembre 2024
% ============================================================================

\chapter{Référentiel des Outils}
\label{ann:outils}

Les outils cités dans le cours sont rangés par usage. La plupart sont libres
ou gratuits. L'annexe n'est pas une liste d'achats~: un outil ne vaut que par
la méthode de celui qui l'emploie, sa version notée et sa validation sur une
pièce de référence (Chapitre~\ref{chap:normes-droit}, ISO/IEC~27041).

\begin{boxalert}
\textbf{Règle de l'autorisation écrite.} Aucun outil d'acquisition,
d'analyse ou de balayage ne s'emploie sur un système sans l'autorisation
écrite de celui qui en a le droit, et, le cas échéant, de l'hébergeur ou du
prestataire. Sans autorisation, l'usage peut constituer un accès non
autorisé au sens de l'article~65, alinéa~2, de la loi n°~2010/012.
\end{boxalert}

\section*{Acquisition}

\begin{tableaucours}{L{3.6cm}L{5.0cm}L{3.6cm}}{%
    \tableauHeader{Outil} &
    \tableauHeader{Usage} &
    \tableauHeader{Chapitre}
}
Bloqueur d'écriture matériel & Acquisition à froid sans écriture sur la
source & \ref{chap:acquisition} \\
FTK Imager & Images de disques (brut, E01), captures mémoire, calcul
d'empreintes & \ref{chap:acquisition} \\
\texttt{dd}, \texttt{dc3dd}, \texttt{ewfacquire} & Images brutes ou E01 sous
Linux & \ref{chap:acquisition} \\
WinPmem, DumpIt & Capture mémoire sous Windows & \ref{chap:acquisition} \\
LiME, AVML & Capture mémoire sous Linux & \ref{chap:acquisition} \\
\texttt{sha256sum} & Calcul et vérification d'empreintes &
\ref{chap:preuve} \\
\end{tableaucours}

\section*{Analyse}

\begin{tableaucours}{L{3.6cm}L{5.0cm}L{3.6cm}}{%
    \tableauHeader{Outil} &
    \tableauHeader{Usage} &
    \tableauHeader{Chapitre}
}
The Sleuth Kit, Autopsy & Systèmes de fichiers, fichiers supprimés,
artefacts & \ref{chap:disque} \\
Volatility~3 & Analyse de captures mémoire & \ref{chap:memoire} \\
Wireshark, tshark & Analyse de captures réseau & \ref{chap:reseau} \\
Zeek & Journaux par protocole à partir d'une capture & \ref{chap:reseau} \\
log2timeline, Plaso & Super-chronologies & \ref{chap:chronologie} \\
Visionneuses d'événements Windows & Lecture des journaux EVTX &
\ref{chap:disque} \\
\end{tableaucours}

\section*{Vulnérabilités}

\begin{tableaucours}{L{3.6cm}L{5.0cm}L{3.6cm}}{%
    \tableauHeader{Outil ou ressource} &
    \tableauHeader{Usage} &
    \tableauHeader{Chapitre}
}
Scanners de vulnérabilités (libres ou commerciaux) & Balayage authentifié ou
non & \ref{chap:identifier} \\
Calculateur CVSS du FIRST & Score à partir d'un vecteur &
\ref{chap:vulnerabilite} \\
EPSS (FIRST) & Probabilité d'exploitation, interface de consultation &
\ref{chap:vulnerabilite} \\
Catalogue KEV (CISA) & Liste des failles exploitées &
\ref{chap:vulnerabilite}, \ref{chap:gerer} \\
NVD, CVE, CWE & Descriptions et références des failles &
\ref{chap:vulnerabilite} \\
CIS Benchmarks & Références de configuration & \ref{chap:identifier} \\
\end{tableaucours}

\section*{Dans NDAP SUKULU}

Les dix laboratoires du cours n'exigent aucune installation~: leurs pièces
(images réduites, captures, journaux, tableurs de balayage) se lisent avec
les outils du bureau d'analyse de la plateforme ou, pour la version en
ligne, avec un tableur et un éditeur de texte. Chaque laboratoire fournit
son manifeste d'empreintes SHA-256, à vérifier avant tout examen.

    % ============================================================================
% Fichier  : annexes/Ann_C_modeles.tex
% Titre    : Modèles de Documents
% Version  : 1.0 -- novembre 2026
% ============================================================================

\chapter{Modèles de Documents}
\label{ann:modeles}

Les modèles suivants reprennent les documents utilisés dans le cours. Ils se
recopient tels quels dans un tableur ou un traitement de texte, et
s'adaptent à l'organisation.

\section*{M1. Registre des actions}

\begin{tableaucours}{L{1.6cm}L{1.6cm}L{2.0cm}L{3.6cm}L{3.0cm}}{%
    \tableauHeader{N°} &
    \tableauHeader{Date, heure} &
    \tableauHeader{Auteur} &
    \tableauHeader{Action} &
    \tableauHeader{Résultat, empreinte}
}
001 & & & & \\
002 & & & & \\
\end{tableaucours}

\noindent Chapitres~\ref{chap:prologue} et~\ref{chap:enquete}.

\section*{M2. Formulaire de chaîne de possession}

\begin{enumerate}
    \item Identifiant de la pièce et numéro de scellé.
    \item Description~: nature, marque, modèle, numéro de série, capacité.
    \item Provenance~: système, emplacement, propriétaire.
    \item Recueil~: date, heure, lieu, personne, méthode, outil et version.
    \item Empreintes~: algorithme, valeur, heure, témoin.
    \item Conditionnement~: enveloppe, sac, scellé, photographie.
    \item Transferts~: remettant, destinataire, date, heure, motif, état du
          scellé, signatures.
    \item Examens~: personne, date, opération, copie utilisée.
    \item Conservation~: lieu, conditions.
    \item Sort final~: restitution, destruction, archivage, autorisation.
\end{enumerate}

\noindent Chapitre~\ref{chap:preuve}.

\section*{M3. Lettre de mission d'enquête}

\begin{itemize}
    \item Mandant et enquêteur~; date~; fondement (contrat, réquisition).
    \item Questions de l'enquête, numérotées.
    \item Périmètre~: systèmes, comptes, période.
    \item Autorisations~: signataires, hébergeurs et prestataires concernés.
    \item Règles de conservation des pièces et de remise à l'autorité
          judiciaire.
    \item Livrables et destinataires.
\end{itemize}

\noindent Chapitres~\ref{chap:investigation} et~\ref{chap:normes-droit}.

\section*{M4. Matrice des hypothèses concurrentes}

\begin{tableaucours}{L{5.6cm}cccc}{%
    \tableauHeader{Pièce} &
    \tableauHeader{H1} & \tableauHeader{H2} & \tableauHeader{H3} &
    \tableauHeader{H4}
}
& & & & \\
\end{tableaucours}

\noindent C~: compatible~; I~: incompatible~; N~: neutre. On retient
l'hypothèse qui a le moins d'incompatibilités (Chapitre~\ref{chap:investigation}).

\section*{M5. Tableau des horloges}

\begin{tableaucours}{L{3.6cm}cccL{2.6cm}}{%
    \tableauHeader{Source} &
    \tableauHeader{Fuseau} &
    \tableauHeader{Dérive} &
    \tableauHeader{Précision} &
    \tableauHeader{Synchronisée}
}
& & & & \\
\end{tableaucours}

\noindent Chapitre~\ref{chap:chronologie}.

\section*{M6. Fiche de constat de vulnérabilité}

\begin{itemize}
    \item Identifiant (CVE), faiblesse (CWE), produit et version.
    \item Actif~: nom, propriétaire, exposition, importance.
    \item Preuve du constat~: méthode (authentifiée ou non), élément lu.
    \item Vérification~: confirmé, faux positif, non vérifiable.
    \item Vecteur CVSS et score, avec la mention B, BT, BE ou BTE.
    \item EPSS du jour~; inscription au catalogue KEV et date.
    \item Classe de priorité, échéance, responsable.
    \item Correction~: voie choisie, date, preuve de vérification.
\end{itemize}

\noindent Chapitres~\ref{chap:vulnerabilite} à~\ref{chap:gerer}.

\section*{M7. Demande d'exception}

\begin{itemize}
    \item Faille, classe, actif.
    \item Motif, avec sa preuve.
    \item Mesure compensatoire, sa date de mise en place et sa
          vérification.
    \item Durée et date de revue.
    \item Propriétaire du risque et décideur, conformes à la politique.
\end{itemize}

\noindent Chapitre~\ref{chap:gerer}.

\section*{M8. Plan d'un rapport d'expertise}

Mission (questions mot pour mot)~; pièces reçues~; opérations~;
constatations~; discussion~; limites~; réponses aux questions, cotées~;
annexes (Chapitre~\ref{chap:rapport}).

\section*{M9. Tableau des propositions cotées}

\begin{tableaucours}{cL{6.6cm}ccc}{%
    \tableauHeader{N°} &
    \tableauHeader{Proposition} &
    \tableauHeader{Date 1} &
    \tableauHeader{Date 2} &
    \tableauHeader{Date 3}
}
P1 & & & & \\
\end{tableaucours}

\noindent Avec, à chaque date, $\mu$, $U$ et $H$
(Chapitre~\ref{chap:auditologie}).

    % ============================================================================
% Fichier  : annexes/Ann_D_textes.tex
% Titre    : Textes et Normes de Référence
% Version  : 1.0 -- novembre 2024
% ============================================================================
% Les extraits de la loi n° 2010/012 sont recopiés du texte officiel détenu
% dans COURS\Normes\loi_2010-012_cybersecurite_cybercriminalite.pdf (lu le
% 08/11/2024).
% ============================================================================

\chapter{Textes et Normes de Référence}
\label{ann:textes}

Cette annexe donne, pour chaque texte cité, son degré de vérification.
Trois degrés sont utilisés~: \emph{texte officiel lu} (article cité entre
guillemets d'après le texte officiel)~; \emph{paraphrasé} (norme sous
licence, dont le contenu est résumé sans citation)~; \emph{à vérifier}
(information dont l'auteur n'a pas pu contrôler la source à la date de
rédaction).

\section*{Droit camerounais}

\begin{tableaucours}{L{3.8cm}L{5.2cm}L{3.2cm}}{%
    \tableauHeader{Texte} &
    \tableauHeader{Articles utilisés} &
    \tableauHeader{Degré}
}
Loi n°~2010/012 du 21~décembre 2010 relative à la cybersécurité et à la
cybercriminalité & 26 al.~4, 29 al.~1 à~3, 30 al.~1, 52, 53, 54, 55, 56,
65 al.~2, 69, 72, 91 al.~3 & Texte officiel lu \\
Loi n°~2010/012 & 26 al.~3, 57, 58, 66, 68, 71 & Texte officiel lu, cités
par résumé \\
Loi n°~2024/017 du 23~décembre 2024 (données personnelles) & Mentionnée
seulement & Renvoi à l'UE JUR01 \\
Code de procédure pénale & Preuve, expertise & À vérifier \\
Réglementation de la supervision des établissements de microfinance &
Notification d'incident & À vérifier \\
\end{tableaucours}

\section*{Extraits de la loi n°~2010/012 non reproduits dans les chapitres}

\begin{boxjuris}
\textbf{Article~66, alinéa~1}~: \textit{« Est puni d'un emprisonnement de
deux (02) à cinq (05) ans et d'une amende de 1.000.000 (un million) à
2.000.000 (deux millions) F CFA ou de l'une de ces deux peines seulement,
celui qui entraîne la perturbation ou l'interruption du fonctionnement d'un
réseau de communications électroniques ou d'un équipement terminal, en
introduisant, transmettant, endommageant, effaçant, détériorant, modifiant,
supprimant ou rendant inaccessibles les données. »}

\medskip
\textbf{Article~68, alinéa~1}~: \textit{« Est puni d'un emprisonnement de
cinq (05) à dix (10) ans et d'une amende de 10.000.000 (dix millions) à
50.000.000 (cinquante millions) F CFA ou de l'une de ces deux peines
seulement, celui qui accède ou se maintient, frauduleusement, dans tout ou
partie d'un réseau de communications électroniques ou d'un système
d'information en transmettant, endommageant, provoquant une perturbation
grave ou une interruption du fonctionnement dudit système ou dudit
réseau. »}

\medskip
\textbf{Article~71}~: \textit{« Est puni d'un emprisonnement de deux (02) à
cinq (05) ans et d'une amende de 1.000.000 (un million) à 25.000.000 (vingt
cinq millions) F CFA, celui qui introduit sans droit, des données dans un
système d'information ou dans un réseau de communications électroniques en
vue de supprimer ou de modifier les données qui en sont contenues. »}
\end{boxjuris}

\section*{Conventions internationales}

\begin{tableaucours}{L{4.6cm}L{4.2cm}L{3.4cm}}{%
    \tableauHeader{Texte} &
    \tableauHeader{Objet} &
    \tableauHeader{Degré}
}
Convention de Budapest (Conseil de l'Europe, 2001) & Cybercriminalité,
conservation rapide, entraide & Statut du Cameroun à vérifier \\
Convention de Malabo (Union africaine, 2014) & Cybersécurité et données
personnelles~; en vigueur en 2023 & Statut du Cameroun à vérifier \\
\end{tableaucours}

\section*{Normes et référentiels}

\begin{tableaucours}{L{3.6cm}L{5.6cm}L{3.0cm}}{%
    \tableauHeader{Texte} &
    \tableauHeader{Objet} &
    \tableauHeader{Degré}
}
ISO/IEC 27037:2012 & Identification, recueil, acquisition, préservation &
Paraphrasé \\
ISO/IEC 27041:2015 & Adéquation des méthodes & Paraphrasé \\
ISO/IEC 27042:2015 & Analyse et interprétation & Paraphrasé \\
ISO/IEC 27043:2015 & Principes et processus d'investigation & Paraphrasé \\
ISO/IEC 27035 & Gestion des incidents & Paraphrasé \\
ISO/IEC 27002:2022, mesure 8.8 & Gestion des vulnérabilités techniques &
Paraphrasé \\
ISO/IEC 17025:2017 & Compétence des laboratoires & Paraphrasé \\
ISO/IEC 29147, 30111 & Divulgation et traitement des vulnérabilités &
Paraphrasé \\
NIST SP 800-86, 800-61 r2, 800-115, 800-40 r4 & Forensique, incidents,
évaluation, correctifs & Publications libres, paraphrasées \\
RFC 3227 & Ordre de volatilité & Publication libre \\
CVSS v4.0 et v3.1 (FIRST) & Notation de la gravité & Spécification libre \\
EPSS (FIRST), KEV et SSVC (CISA) & Probabilité, exploitation, décision &
Ressources libres \\
\end{tableaucours}

    % ============================================================================
% Fichier  : annexes/Ann_E_bibliographie.tex
% Titre    : Bibliographie Commentée
% Version  : 1.0 -- novembre 2024
% ============================================================================

\chapter{Bibliographie Commentée}
\label{ann:bibliographie}

Les références complètes figurent dans la bibliographie en fin d'ouvrage.
Cette annexe indique, partie par partie, par où commencer.

\section*{Partie I. Fondements}

E.~Casey~\cite{casey2004} pose les notions de trace et de preuve numérique
avec de nombreux exemples judiciaires. Le guide NIST SP~800-86~\cite{nist80086},
court et gratuit, est la meilleure entrée dans les phases d'une enquête.
R.~Heuer~\cite{heuer1999} reste la référence sur les biais de l'analyste et
l'analyse des hypothèses concurrentes. L'ISO/IEC~27037~\cite{iso27037_2012}
et l'ISO/IEC~27043~\cite{iso27043_2015} donnent le cadre normatif.

\section*{Partie II. Méthode}

Pour le disque, B.~Carrier~\cite{carrier2005} explique les systèmes de
fichiers de l'intérieur~; on le lit avec Autopsy ouvert. Pour la mémoire,
M.~H.~Ligh et ses coauteurs~\cite{ligh2014memory} couvrent Windows, Linux et
macOS~; la documentation de Volatility~3~\cite{volatility3} met à jour les
commandes. Pour le réseau, R.~Bejtlich~\cite{bejtlich2013nsm} pense la
supervision, C.~Sanders~\cite{sanders2017packet} apprend à lire les paquets.
La RFC~3227~\cite{rfc3227} tient en quelques pages et se relit avant chaque
intervention.

\section*{Partie III. Cadre normatif et juridique}

La loi n°~2010/012~\cite{loi2010012} se lit dans son texte officiel,
chapitre de procédure d'abord. Les Conventions de
Budapest~\cite{budapest2001} et de Malabo~\cite{ua2014malabo} éclairent
l'entraide.

\section*{Partie IV. Vulnérabilités}

Les spécifications du CVSS~4.0~\cite{first2023cvss4}, la documentation de
l'EPSS~\cite{firstepss}, le catalogue KEV~\cite{cisakev} et le guide
SSVC~\cite{cisa2022ssvc} sont gratuits et courts. Le NIST
SP~800-40~\cite{nist80040} traite la gestion des correctifs comme une
maintenance planifiée~; le NIST SP~800-115~\cite{nist800115} distingue
évaluation et test d'intrusion.

\section*{Partie V. Défis avancés}

L'ISO/IEC~27042~\cite{iso27042_2015} pour l'interprétation et le rapport~;
l'ISO/IEC~27035~\cite{iso27035} et le NIST SP~800-61~\cite{nist80061r2} pour
les enseignements tirés d'un incident.

\section*{Partie VI. Recherche}

Les deux articles fondateurs du programme de l'auditologie sont déposés sur
HAL~: « Vers l'auditologie »~\cite{minka2026auditologie} et l'article sur
l'instrument trichotomique~\cite{minka2026nek}. Les articles suivants du
programme sont en préparation. C.~E.~Shannon~\cite{shannon1948} reste la
source de la mesure d'entropie.

    % ============================================================================
% Fichier  : annexes/Ann_F_contacts.tex
% Titre    : Contacts et Ressources
% Version  : 1.0 -- novembre 2026
% ============================================================================

\chapter{Contacts et Ressources}
\label{ann:contacts}

\section*{Au Cameroun}

\begin{tableaucours}{L{4.6cm}L{7.6cm}}{%
    \tableauHeader{Organisme} &
    \tableauHeader{Rôle pour le cours}
}
Agence nationale des technologies de l'information et de la communication
(ANTIC) & Agents habilités à enquêter (loi n°~2010/012, art.~52)~; audits de
sécurité~; équipe nationale de réponse aux incidents \\
Parquet et officiers de police judiciaire & Plainte, réquisitions, scellés,
entraide \\
Experts près les cours d'appel & Expertise judiciaire sur réquisition
(art.~56) \\
\end{tableaucours}

\section*{Ressources internationales}

\begin{tableaucours}{L{4.6cm}L{7.6cm}}{%
    \tableauHeader{Ressource} &
    \tableauHeader{Adresse}
}
Programme CVE & \url{https://www.cve.org/} \\
Catalogue CWE & \url{https://cwe.mitre.org/} \\
Base NVD & \url{https://nvd.nist.gov/} \\
CVSS et EPSS (FIRST) & \url{https://www.first.org/cvss/},
\url{https://www.first.org/epss/} \\
Catalogue KEV (CISA) &
\url{https://www.cisa.gov/known-exploited-vulnerabilities-catalog} \\
Matrice ATT\&CK & \url{https://attack.mitre.org/} \\
Volatility~3 & \url{https://github.com/volatilityfoundation/volatility3} \\
The Sleuth Kit et Autopsy & \url{https://www.sleuthkit.org/} \\
\end{tableaucours}

\section*{Le cours}

Dépôt du cours~: \url{https://github.com/MaletYon/cours-fav}. NDAP SUKULU en
ligne~: \url{https://maletyon.github.io/cours-fav/ndap-sukulu/}. Compte de
démonstration~: matricule \texttt{ETU001}, mot de passe \texttt{eTU001}.
Contact~: \href{mailto:maletyon@proton.me}{maletyon@proton.me}.


% --------------------------------------------------------------------------
% ÉLÉMENTS FINAUX
% --------------------------------------------------------------------------
\backmatter

    \nocite{*}
    \printbibliography[heading=bibintoc, title={Références bibliographiques}]

    \printindex

% ==============================================================================
% 4E DE COUVERTURE
% ==============================================================================
\clearpage
\thispagestyle{empty}
\pagecolor{CouleurPrimaire}

\begin{tikzpicture}[remember picture, overlay]
    \fill[CouleurAccent]
        (current page.south west) rectangle
        ([yshift=1.8cm] current page.south east);
    \fill[CouleurAccent]
        ([yshift=-1.4cm] current page.north west) rectangle
        ([yshift=-1.7cm] current page.north east);
\end{tikzpicture}

\vspace*{2.5cm}

\begin{center}
    {\color{CouleurAccent}\rule{0.6\linewidth}{1pt}}\\[1.5ex]
    {\color{white}\Large\bfseries À propos de ce cours}\\[1.5ex]
    {\color{CouleurAccent}\rule{0.6\linewidth}{1pt}}
\end{center}

\vspace{1cm}

\begin{center}
\begin{minipage}{0.75\textwidth}
\color{white!85!gray}\setstretch{1.4}
Une enquête numérique part d'une trace et remonte vers une cause. Une
analyse de vulnérabilités part d'une faiblesse connue et cherche à la
fermer avant qu'elle ne laisse une trace. Ce cours enseigne les deux, et
la manière dont l'une nourrit l'autre.

\medskip

Il progresse en six parties, précédées d'un prologue~: les fondements de
l'investigation et de la preuve, la méthode (acquisition, disque, mémoire,
réseau, chronologie), le cadre normatif et juridique, les vulnérabilités
(identifier, coter, corriger), les défis avancés et la recherche sur
l'incertitude d'une enquête.

\medskip

Il s'adresse aux étudiants du Master Professionnel en Sécurité des Systèmes
Numériques et aux professionnels qui exercent comme analystes forensiques,
analystes en réponse à incident, analystes en gestion des vulnérabilités ou
RSSI.
\end{minipage}
\end{center}

\vfill

\begin{center}
    {\color{CouleurAccent}\rule{0.7\linewidth}{1.2pt}}\\[1.2ex]
    {\color{white}\large\bfseries MINKA MI NGUIDJOÏ Thierry Emmanuel}\\[0.4ex]
    {\color{white!80!gray}\small
        CISA~·~CISM~·~CGEIT~·~CRISC~·~CDPSE~·~ISO 27001~·~ISO 22301~·~ISO 9001}\\[0.3ex]
    {\color{white!70!gray}\footnotesize\itshape
        Chercheur au laboratoire d'Ingénierie Mathématiques et des Systèmes d'Information}\\[0.3ex]
    {\color{white!70!gray}\footnotesize\itshape
        Chercheur~|~Cryptologue~|~Auditeur Senior~|~Expert Judiciaire~|~GRC \& Sécurité de l'Information}\\[1.2ex]
    {\color{CouleurAccent}\rule{0.7\linewidth}{1.2pt}}
\end{center}

\vspace{1cm}

\afterpage{\nopagecolor}
\end{document}
